%% Theory of Computing
%% File: v015a006.tex
%% Publication date: October 11, 2019
%% Authors: Anindya De, Elchanan Mossel, and Joe Neeman


%
%
%
%

%
%
%

%

\documentclass{toc}

%
%
%
\tocdetails{%
    title = {Noise Stability is Computable and Approximately Low-Dimensional},
    runningtitle = {Noise Stability is Low-Dimensional}, %
%
%
%
%
  number_of_pages = {NN}, 
  number_of_bibitems = {NN},
  number_of_figures = {NN},
  conference_version = {CCC17}, %
    %
    %
%
%
%
  author = {Anindya De, Elchanan Mossel, and Joe Neeman},
    %
    %
    %
  authorlist = {Anindya De, Elchanan Mossel, Joe Neeman},
    %
    %
    %
    %
    %
    %
    %
%
%
  %
  %
  %
  %
%
  acmclassification = {G.3, G.1.6},
    %
    %
  amsclassification = {60},
    %
    %
  keywords = {noise stability, Gaussian surface area, computability},
    %
    %
    %
    %
    %
}   %

    %
    %
    %
    %
    %
    %
    %
    %
    %
    %
    %
    %
    %
    %
    %
    %
    %
    %
    %
    %
    %
    %
    %
    %
    %
    %
    %
    %
    %

%
%
\tocdetails{%
 volume = 15,
 number = 6,
 year = 2019,
 specissue={\cccBH},  %
%
%
%
%
%
%
 received = {September 13, 2017},
 revised = {July 7, 2019},
 published = {October 11, 2019},
%
%
 doi = {10.4086/toc.2019.v015a006},
}   %

\newcommand{\E}{{\bf E}}
\newcommand{\Inner}{\mathrm{In}}
\newcommand{\Outer}{\mathrm{Out}}
\newcommand{\num}{\mathrm{num}}
\newcommand{\bP}{{\bf P}}
\newcommand{\bQ}{{\bf Q}}
\newcommand{\bR}{{\bf R}}
\newcommand{\fh}{{f_{\mathsf{junta}}}}
\newcommand{\fhv}{{f_{\mathsf{junta,v}}}}
\newcommand{\inr}[2]{\langle #1, #2 \rangle}
\newcommand{\calF}{\mathcal{F}}
\newcommand{\Stab}{\mathsf{Stab}}
\newcommand{\Var}{\mathrm{Var}}
\newcommand{\grade}[1]{{{#1}}}
\newcommand{\ignore}[1]{}
\newcommand{\Coeff}{\mathsf{Coeff}}
\newcommand{\Num}{\mathsf{Num}}

\newtheorem{cond}[theorem]{Condition}
\newtheorem{obs}[theorem]{Observation}
\newtheorem*{definition*}{Definition}
\newtheorem*{lemma*}{Lemma}
\newenvironment{proofof}[1]{\par{\noindent \bf Proof of #1:}}{\qed\par}

\begin{document}

\begin{frontmatter}%
  
%
    \title{Noise Stability is Computable and Approximately Low-Dimensional\titlefootnote{A conference version of this paper appeared in the \href{http://drops.dagstuhl.de/portals/extern/index.php?semnr=16039}{Proceedings of the 32nd Computational Complexity Conference, 2017}.
    }} %

%
%

%
%
%

%
%
%
%
%
%
%
%
%
%
%
%


%
%
%
%

%

\author[de]{Anindya De\thanks{Supported by NSF grant CCF 1926872 (transferred from CCF 1814706). Most of the work done as a faculty at Northwestern University. }}
\author[mossel]{Elchanan Mossel\thanks{Supported by DMS-1737944 and  ONR N00014-16-1-2227, N00014-17-1-2598. }}
\author[neeman]{Joe Neeman\thanks{Funded by a Deutsche
Forschungs-gemeinschaft (DFG, German Research Foundation) under Germany’s Excellence Strategy
GZ 2047/1, Projekt-ID 390685813 and a Sloan Fellowship.}}

%
%
%
%

%
%
%
%
%
%
\begin{abstract}
    The notion of Gaussian noise stability plays an important role in hardness
    of approximation in theoretical computer science as well as in the theory
    of voting.
    The Gaussian noise stability of a partition of $\R^n$ is simply the probability
    that two correlated Gaussian vectors both fall into the
%
same part. In many   %
applications, the goal is to find an optimizer of noise stability among all possible partitions of $\R^n$  to $k$ parts with given Gaussian measures $\mu_1,\ldots,\mu_k$. We call a partition $\epsilon$-optimal, if its noise stability is optimal up to an additive $\epsilon$.
In this paper, we give a computable function $n(\epsilon)$
such that an $\epsilon$-optimal partition exists in $\R^{n(\epsilon)}$.
This result has implications for the computability of certain problems
in non-interactive simulation, which are addressed in a subsequent
%
%
paper.  %
\end{abstract}

\iffalse %
%
%
%
%
%
%
%
%
%
%
%
%
%
%
%
%
%
%
%
%
%

\tocarxivcategory{math.PR}

\fi %

\end{frontmatter}

\section{Introduction} 

\paragraph{Isoperimetric theory and noise stability.}
Isoperimetric problems have been studied in mathematics since antiquity, and
they have become central to  mathematics and theoretical computer
science.
Our interest in this work is in a modern version of the isoperimetric problem: that of
maximizing Gaussian noise stability. This problem was originally studied by Borell in
the context of stochastic processes~\cite{Borell:85}.
Benjamini, Kalai, and Schramm~\cite{BKS:99} introduced Boolean noise stability
(both the notion and the terminology) to computer science: roughly speaking, a
Boolean function---that is, a function $\{0, 1\}^n \to \{0, 1\}$---is
\emph{noise sensitive} if slightly rerandomizing the inputs has the effect of
completely rerandomizing the output. This notion naturally extends to functions
$\{0, 1\}^n \to \{0, \dots, k - 1\}$, or, equivalently, to partitions of $\{0,
1\}^n$ into $k$ parts; this extension and its applications in theoretical computer
science
%
are   %
discussed in~\cite{IsakssonMossel:12}.

Boolean noise stability has a natural Gaussian analogue, in which
the Boolean hypercube $\{0, 1\}^n$ is replaced by $\R^n$, and the Boolean
noise is replaced by Gaussian noise.  The \emph{invariance principle} of~\cite{MOO10}
%
shows     %
a deep connection between these two notions of noise
stability. In particular, it allows one to translate results about Gaussian
noise stability into results about Boolean noise stability, which
are typically harder. This connection is responsible for many results
in hardness of approximation~\cite{KKMO07}. For more on the applications of
noise stability in the theory of computing, see the book by O'Donnell~\cite{ODonnell:book}.

We equip $\mathbb{R}^n$ with the standard
$n$-dimensional Gaussian measure, denoted by $\gamma_n$.
In particular, for any $x \in \mathbb{R}^n$, the density $\gamma_n(x)$ is given by 
\[
\gamma_n(x) = \frac{1}{(2\pi)^{n/2}}  \eee^{-{\Vert x \Vert_2
^2}/{2}}.
\]
 The Ornstein-Uhlenbeck operator $P_t$ is defined for $t \in [0, \infty)$ and bounded, measurable $f: \mathbb{R}^n \rightarrow \mathbb{R}$ by
$$
(P_t f)(x) = \int_{y \in \mathbb{R}^n} f(\eee^{-t} \cdot x + \sqrt{1- \eee^{-2t}} \cdot y) d \gamma_n(y).
$$
We define the (Gaussian) noise stability of $f$ by $\mathsf{Stab}_t(f) :=
\mathbf{E}[f \cdot P_t f + (1-f) \cdot (1 - P_t f)]$, where the expectation is with respect to
$\gamma_n$. If $f$ denotes the indicator function of a set (call it
$A$), then $\mathsf{Stab}_t(f)$ is the probability that two
$\eee^{-t}$-correlated Gaussian random variables fall either both in $A$ or both in $A^c$. In
the limit $t \rightarrow 0^+$, noise stability is related to the Gaussian
surface area of the set $A$~\cite{Ledoux:94}.

\paragraph{Halfspaces are the most stable sets.}
Since Borell's work in the 1980s~\cite{Borell:85},
halfspaces have been known to
maximize noise stability subject to a measure constraint.
We will state this
fact in a slightly convoluted way, in order to more easily generalize it.  Let
$\Delta_k$ denote the standard simplex in $\R^k$ (\ie, the convex hull
of  the standard unit vectors $\{\mathbf{e}_1, \ldots, \mathbf{e}_k\}$).
Using $[k]$ to denote $\{1,\ldots, k\}$, we observe that any function with range $[k]$ naturally embeds in $\Delta_k$ by identifying $i
\in [k]$ with $\mathbf{e}_i$. 
Moreover, we extend the Ornstein-Uhlenbeck operator to act on vector-valued
functions in the obvious way: if $f = (f_1, \dots, f_k): \R^n \to \R^k$
then $P_t f = (P_t f_1, \dots, P_t f_k)$.
Finally, we say that $f = (f_1, f_2): \R^n \to \Delta_2$ is a \emph{halfspace} if
there exist $a, b \in \R^n$ such that $f_1(x) = 1_{\{\langle x - a, b\rangle \le 0\}}$
for all $x \in \R^n$.
(Here, $\langle v, w \rangle$ denotes
the Euclidean inner product $\sum_{i=1}^n v_i \cdot w_i$.)

\begin{theorem}[Borell]~\label{thm:Borell}
For any $g: \mathbb{R}^n \rightarrow \Delta_2$ and any $t \ge 0$, every halfspace
$f: \mathbb{R}^n \rightarrow \Delta_2$ with $\mathbf{E}[f] = \mathbf{E}[g]$
satisfies 
$$
\mathbf{E}[\langle f, P_t f \rangle] \ge \mathbf{E}[\langle g, P_t g \rangle]. 
$$
\end{theorem}

\expref{Theorem}{thm:Borell} has many applications in computational complexity.
Most famously, it can be combined with an
\emph{invariance principle}~\cite{MOO10}
in order to prove a number of tight
%
%
%
%
%
%
%
Unique-Games-hardness of approximation results~\cite{KKMO04}.

\paragraph{More parts?}
It is straightforward to extend the notion of noise stability to partitions
with more than two parts. Namely, a partition of $\R^n$ can be described by $f:
\mathbb{R}^n \rightarrow [k]$. Identifying $[k]$ with $\{\mathbf{e}_1, \dots,
\mathbf{e}_k\}$ as above, we define the noise stability of such a partition by
$\mathbf{E}[\langle f, P_t f \rangle]$.
One may then ask for an analogue of \expref{Theorem}{thm:Borell} where
$\Delta_2$ is replaced by $\Delta_k$ for some $k \ge 3$.

Even the three-part version of \expref{Theorem}{thm:Borell} turns out to be
significantly harder than the two-part one. We will try to explain why, by analogy
with isoperimetric problems. First of all, the Euclidean, isoperimetric
analogue for $k=3$ is known as the ``double bubble'' problem.
The well-known  ``double bubble conjecture'' states that the minimum
total surface area of two bodies separating and enclosing two given (Lebesgue)
volumes is achieved by two spheres meeting at $120^\circ$. After being open
for more than a century, this problem was settled rather recently~\cite{HMRR:00, HMRR:02, RHLS:03, Reichardt:08}.
For the Gaussian space,~\cite{Corwin} showed that for some small  positive
constant $c>0$, the Gaussian surface area of three-partitions is minimized by
the ``standard simplex partition'' as long as the measure of each part is within $1/3 \pm c$. 

Since halfspaces both maximize noise stability and minimize Gaussian surface
area, and since the standard simplex partition is known to minimize multi-part
Gaussian surface area in certain cases, it seems natural to guess that the
standard simplex partition also maximizes multi-part noise stability.  This was
explicitly conjectured in~\cite{KKMO04}, in the special case that all of the
parts in the partition have equal measure.  However, a somewhat surprising recent result~\cite{heilman2014} showed that the standard
simplex partition fails to maximize multi-part noise stability \emph{unless}
all of the parts have equal measure.  On the other hand, there is also some
support for the conjecture in the equal-measure case: Heilman~\cite{Heilman:14}
showed that the conjecture is true if the noise rate $t$ is larger than some explicit function of the ambient dimension $n$. 

\paragraph{Approximate noise stability of multipartitions?}
In light of the uncertainty about optimal partitioning for $k \ge 3$, one can
ask a more modest question. Given $k \ge 3$, $t>0$, and prescribed measures $v \in \Delta_k$ for
the $k$ parts, let $\alpha_n = \alpha_n(k, v)$ be the maximal noise stability that can
be obtained in $\R^n$ under these constraints. Since the Gaussian measure
is a product measure, $\alpha_n$ is clearly non-decreasing in $n$.
Since it is bounded above by one, it has a limit as $n \to \infty$.

Our main result is that there is a computable
$n_0=n_0(k,t,\epsilon)$ such that $\alpha_{n_0} \ge \alpha_{m} - \epsilon$ for
all $m \in \mathbb{N}$. Although the bound on $n_0$ that we give is not
particularly good (in fact, it is not primitive recursive), the key point is that it is  computable.
As a consequence, up to error $\epsilon$, the noise-stable partition is
also  computable.
We conclude the introduction by an open question: 
\begin{question}
    Fix $k \ge 3$ and measure constraints $v \in \Delta^k$. Does there exist $n_0$ such that $\alpha_{n_0} = \alpha_{n}$ for all $n \ge n_0$? 
\end{question}
Our current techniques are not suitable for addressing the question above. 

\subsection{Acknowledgments} 
The authors started this research at the workshop on ``Information theory and concentration phenomena'' at the IMA, Minneapolis in 2015. The authors would like to thank the organizers of this workshop and the IMA for their hospitality. 

\section{Main theorem and overview of proof technique} \label{sec:theorem}
In order to state the main theorem, we first need to recall the notion of a
\emph{polynomial threshold function.} %
A function  $f: \mathbb{R}^n \rightarrow \{0,1\}$ is said to be a degree-$d$ PTF if there exists a polynomial $p: \mathbb{R}^n \rightarrow \mathbb{R}$ of degree-$d$ such that $f(x) = 1$ if and only if $p(x) > 0$. We will need a $k$-ary generalization of this definition. We note that there are several possible ways to generalize the notion of PTFs to $k$-ary PTFs and our particular choice is dictated by the convenience of using the relevant results from~\cite{DS14}. 

\begin{definition}
A function $f: \mathbb{R}^n \rightarrow [k]$ is said to be a multivariate PTF
if there exist polynomials $p_1, \ldots, p_k : \mathbb{R}^n \rightarrow
\mathbb{R}$ such that 
$$
f(x) = \begin{cases}
j  & \textrm{if} \  p_j(x)>0 \ \textrm{ and  for } i \neq j, \  p_{i}(x) \le 0, \\
1 & \textrm{otherwise.}\end{cases}
$$
In this case, we denote $f = \mathsf{PTF}(p_1, \ldots, p_k)$. Further, $f$ is said to be a degree-$d$ multivariate PTF if   $p_1, \ldots, p_k$ are  degree-$d$ polynomials.
\end{definition}
We now state the main theorem of this paper. We set the convention that, unless
explicitly mentioned otherwise, the underlying distribution is $\gamma_n$, the
standard $n$-dimensional Gaussian measure. 
Likewise, given any measurable function $f: \R^n \to \R^k$, $\mathbf{E}[f]$ denotes its (vector-valued) expectation (with respect to $\gamma_n$) and $\mathsf{Var}(f)$ denotes its covariance matrix.  
For a vector $v \in \R^k$ and a number $p \ge 1$, $\|v\|_p$ denotes the
$\ell^p$ norm, $(\sum_i |v_i|^p)^{1/p}$.

\begin{theorem}~\label{thm:main-stab}
There exist computable functions $n_0 = n_0(t, k, \epsilon)$ and $d= d(t, k, \epsilon)$ such that the following holds for all $d$, $k$, $t>0$ and $\epsilon>0$. 
Let $f: \mathbb{R}^n \rightarrow [k]$ such that $\mathbf{E}[f] = \mu \in \mathbb{R}^k$. For $t>0$ and $\epsilon>0$,  there is a degree-$d$ PTF $g: \mathbb{R}^{n_0} \rightarrow [k]$ such that 
\begin{enumerate}
\item $\Vert \mathbf{E}[f] - \mathbf{E}[g] \Vert_1 \le \epsilon$, 
\item $\mathbf{E}[\langle g, P_t g \rangle] \ge \mathbf{E}[\langle f, P_t f \rangle] - \epsilon$.
\end{enumerate}
%
We note  %
that while the functions $n_0(\cdot, \cdot, \cdot)$ and $d(\cdot, \cdot, \cdot)$ are computable, $n_0$ is not primitive recursive.
\end{theorem}

Note that the above theorem automatically implies that a function $g$ satisfying the above properties can be  computed up to some additional error $\epsilon$. This is because the set of degree-$d$ PTFs on $\mathbb{R}^{n_0}$ has a finite $\epsilon$-cover (whose elements 
can be computed). 

\begin{remark}
    As explained above, we may interpret the expression $\E[\langle f, P_t f
      \rangle]$
in  %
    \expref{Theorem}{thm:Borell} in terms of Gaussian random variables
    taking values in $\R^n$. The correlation of these random
    variables is $\eee^{-t}$ (which is always positive), but in many applications
    it is useful for allow for negative correlations also.
    For example, \expref{Theorem}{thm:Borell} remains true with negative
    correlations, but with the inequality reversed.

    We do not know how to prove the analogue of \expref{Theorem}{thm:main-stab} in
    the case of negative correlations, although we do discuss some related
    results in the follow-up paper~\cite{DMN:17b}. The reason for this failure
    is that a certain trick which works in the $k=2$ setting fails for $k\ge3$:
    suppose instead of maximizing $\E[\langle f, P_t f \rangle]$, we attempt to
    minimize $\E[\langle f_1, P_t f_2 \rangle]$ where $f_1$ and $f_2$ are not
    required to be the same function. In the setting of
    \expref{Theorem}{thm:Borell}, the minimizers miraculously satisfy $f_1(x) =
    f_2(-x)$ for all $x$, and it follows that $f_1$ minimizes the noise stability
    for the \emph{negative} correlation $-\eee^{-t}$. As far as we know,
    this miracle does not occur for $k \ge 3$.
\end{remark}

\paragraph{Proof sketch.}
The proof of \expref{Theorem}{thm:main-stab} consists of the following main steps: 

\paragraph{From general partitions to PTF.} 
The first step in the proof is to show that given any $f: \mathbb{R}^n \rightarrow [k]$, there is a multivariate PTF $g': \mathbb{R}^n \rightarrow [k]$ which meets the two criteria in \expref{Theorem}{thm:main-stab} and has degree $d= d(t, k, \epsilon)$, for some computable function $d( t, k ,\epsilon)$.  
In other words, $g'$ satisfies $\Vert \mathbf{E}[f] - \mathbf{E}[g'] \Vert_1 \le \epsilon$ and  $\mathsf{Stab}_t(g') \ge \mathsf{Stab}_t(f')  - \epsilon$. 
 This is done in \expref{Section}{sec:reduction}. Note that main difference between the desired conclusion of \expref{Theorem}{thm:main-stab} and what is accomplished in this step is that the ambient dimension remains $n$ as opposed to a bounded dimension $n_0$. 
 
Why is this true? The basic intuition is that if $f$ is noise stable then it should have most of its Hermite expansion weight at low degree. 
Therefore we should be able to replace $f$ with the PTF where the polynomial is the truncated expansion of $f$. 
There are a number of challenges in formalizing this intuition: 1. We cannot rule out that a positive fraction of the weight of $f$ is at high degrees (perhaps as large as $n$).
2. It is not clear that the PTF obtained this way is noise stable nor that 3. It has the right expected value. 

Our analysis proceeds as follows. We would like to construct $g'$ from $f$ by ``rounding''
$P_\delta f$ for some small $\delta>0$. The advantage of $P_\delta f$ over $f$ is that $P_\delta f$ is guaranteed to have Hermite coefficients that decay at a certain rate.
The rounding of $P_\delta f$ can be performed given some $a \in \mathbb{R}^n$ by considering the function
$g_a: \mathbb{R}^n \to [k]$ which takes the value $i$ whenever $i$ is the largest
coordinate of $P_\delta f - a$. It is not hard to prove that it is possible to choose $a$ such
that $\E[g_a] = \E[f]$; moreover, one can show that this function $g_a$ has better noise
stability than $f$ does. The main obstacle is that the function $g_a$ is not a PTF\@. 
Unfortunately the Hermite decay of $P_\delta f$ does not translate to Hermite decay of $g_a$. 
Instead we use a randomized construction to show that for most $a$'s, $g_a$ has
Hermite decay and can therefore be well approximated by a PTF\@. The
analysis of this construction uses the co-area formula~\cite{federer1969} and gradient
bounds~\cite{Ledoux:94} in the Gaussian space  and draws on ideas from~\cite{Neeman14, KNOW14}.


\paragraph{From PTF in dimension $n$ to a small PTF of
%
bounded-degree    %
polynomials.}   

Given the function $g': \mathbb{R}^n \rightarrow [k]$ of degree $d= d( t, k, \epsilon)$, our next goal is to show  that it is possible to obtain a PTF $g$ on some $n_0 = n_0 (d, k, \epsilon)$ variables such that (i) $\Vert \mathbf{E}[g] - \mathbf{E}[g'] \Vert_1 \le \epsilon$ and (ii) $|\langle g, P_t g \rangle - \langle g', P_t g' \rangle| \le \epsilon$. This part builds on and extends the theory and results of~\cite{DS14}. The key notion introduced in~\cite{DS14} is that of an \emph{eigenregular} polynomial. Namely, a polynomial is said to be $\delta$-eigenregular if for the canonical tensor $\mathcal{A}_p$ associated with the polynomial, the ratio of  the maximum singular value to its Frobenius norm is at most $\delta$ (the tensor notions are explicitly defined later). 

  The key advantage of this definition is that as shown in~\cite{DS14}, when $\delta \rightarrow 0$, the distribution of $p$ (under $\gamma_n$) converges to a normal.  In other words, eigenregular polynomials obey a central limit theorem.  In fact, given $k$ polynomials $p_1, \ldots, p_k$ which are $\delta$-eigenregular, they also obey a multidimensional central limit theorem.  
  
The \emph{regularity lemma} from~\cite{DS14} implies that the polynomials
$p_1,\ldots,p_k$ can be jointly expressed as bounded-size
polynomials in eigenregular homogenous polynomials. More precisely,
for any $\delta > 0$ there is a collection of ``inner'' polynomials
$\{\Inner(p_{s, q, \ell })\}$ for $1 \le s \le k$, $1 \le q \le d$ and $1 \le
\ell \le \num(s,q)$ (where $\num(\cdot, \cdot)$ is an explicitly defined
function), and a collection of ``outer'' polynomials $\{\Outer(p_s)\}$
for $1 \le s \le k$ such that
\[
    p_s = \Outer(p_s) (\{\Inner(p_{s,q, \ell})\}_{q \in [d], \ell \in [\num(s,q)]}).
\]
Moreover, every inner polynomial is $\delta$-eigenregular and homogeneous, and
each outer polynomial $\Outer(p_s)$ has arity
$\num(s) = \sum_{q=1}^d \num(s,q)$
bounded by a computable function of $d$, $k$, and $\delta$.

In~\cite{DS14}, the regularity lemma was used to conclude that the joint
distribution of $p_1,\ldots,p_k$ can be approximated in a bounded dimension as
we can replace each of the inner polynomials by a one dimensional Gaussian. For
our application things are more delicate, as we are not only interested in the
joint distribution of $p_1,\ldots,p_k$ but also in the noise stability of
$p_1,\ldots,p_k$.  For this reason it is important for us to maintain the
degrees of the inner polynomials (each of which is homogenous) and not replace
them with Gaussians.

\paragraph{A small PTF representation.}
In the final step of the proof, we maintain $\Outer(p_s)$ and show how the polynomials 
$\{\Inner(p_{s,q, \ell})\}$ can be replaced by a collection of polynomials $\{\Inner(r_{s,q, \ell})\}$ in bounded dimensions (\ie, bounded in terms of $d$, $k$ and $\epsilon$).
The fact that a collection of homogenous polynomials can be replaced by polynomials in bounded dimensions is a tensor analogue of the fact that for any $k$  vectors in $\mathbb{R}^n$, there exist $k$ vectors in $\mathbb{R}^k$ with the same matrix of inner products. Once such polynomials are found, it is not hard to construct eigenregular polynomials from them by averaging the polynomials over independent copies of random variables.  


\section{Applications} 

Given the wide applicability of Borell's isoperimetric result to combinatorics and theoretical computer science, we believe that \expref{Theorem}{thm:main-stab} will also be widely applicable. We will now point out some applications of this theorem. First, by combining \expref{Theorem}{thm:main-stab} with the invariance principle~\cite{MOO10}, we  
derive a weak $k$-ary analogue of ``Majority is Stablest.''
The analogue of the Ornstein-Uhlenbeck operator is the so-called
\emph{Bonami operator}
%
%
(also referred to  %
as the Bonami-Beckner operator) defined as follows: for $\rho \in [-1,1]$ and $x
\in \{-1,1\}^n$, let $\mathcal{D}_{\rho}(x)$ be the product distribution over
$\{-1,1\}^n$ such that for $y \sim \mathcal{D}_{\rho}(x)$, for all $1 \le i \le
n$, $\mathbf{E}[x_i y_i] =\rho$. 
Then, for any $f: \{-1,1\}^n \rightarrow \mathbb{R}^k$,
\[
  T_{\rho} f(x) = \mathbf{E}_{y \sim D_{\rho}(x)} [f (y)].
\]
Likewise, for any $i \in [n]$ and $z \in \{-1,1\}^{n-1}$, let
$f_{z,-i}: \{-1,1\} \rightarrow \mathbb{R}^k $ denote the function
obtained by restricting all but the $i$-{th} coordinate to $z$. Then,
$\mathsf{Var}(f_{z,-i}) = \mathbf{E}_{x \in \{-1,1\}}[\Vert
f_{z,-i}(x) - \mathbf{E}_{z}[f_{z,-i}] \Vert_2^2]$.  Define the
influence of the $i$-{th} coordinate on $f$ by
$$
\mathsf{Inf}_i(f)= \mathbf{E}_{z \in \{-1,1\}^{n-1}} [\mathsf{Var}(f_{z,-i})]. 
$$
%

\begin{theorem}~\label{thm:Plurality}
There exist computable functions $n_0 = n_0(k,\epsilon,\rho)$ and
$C = C(k,\epsilon,\rho)$ such that the following holds. 
Given any $k \in \mathbb{N}$ and $\rho \in [0,1],\epsilon>0$ and 
$\boldsymbol{\mu} = (\mu_1, \ldots, \mu_k) \in \Delta_k$ and $n \geq
n_0$, there is a function $g = g_{\boldsymbol{\mu}, \rho, \epsilon}: \{-1,1\}^{n} \rightarrow [k]$ such that
\begin{enumerate}
    \item $\max_i \mathsf{Inf}_i(g) \le C/\sqrt{n}$;
    \item $|\Pr_{x \in \{-1,1\}^{n_0}}[g(x) = i ] - \mu_i| \leq  C / \sqrt{n}$; and
    \item for any $f: \{-1,1\}^n \rightarrow [k]$ satisfying $|\Pr_{x \in
\{-1,1\}^n} [f(x) = i ] - \mu_i| \leq \epsilon$ and $\max_i \mathsf{Inf}_i(f)
\le \epsilon$, 
\[
\mathbf{E}[\langle f, T_{\rho} f \rangle] \le \mathbf{E}[\langle g, T_{\rho} g \rangle] + 2 k \epsilon. 
\]
\end{enumerate}
Further, the function $g$ can be computed given the parameters $\boldsymbol{\mu}$, $\rho$ and $\epsilon$. 
\end{theorem}

In particular, the case where $(\mu_1,\ldots,\mu_k) = (1/k,\ldots,1/k)$ implies
that in a tied election between $k$ alternatives, we can find an
$\epsilon$-optimally robust
%
noise-stable   %
voting rule, where the stability is
with respect of each voter randomizing their vote independently with
probability $1-\rho$. Note that in the actual ``Majority is Stablest'' theorem,
the function $g$ is known explicitly (it is the ``majority'' function), whereas
in \expref{Theorem}{thm:Plurality} we only know that $g$ is computable. An anonymous
referee suggested that we call \expref{Theorem}{thm:Plurality} a ``Something is Stablest''
theorem.

The proof of the theorem, which is omitted, is by now a standard reduction from a Gaussian noise stability result to a discrete one~\cite{MOO10, IsakssonMossel:12, DMN:sos-journal}. In one direction, the invariance principle immediately bounds the discrete stability of any partition by the Gaussian stability. In the other direction, starting from an $\epsilon$-optimal Gaussian partition, each Gaussian can be replaced by a normalized sum of independent variables to obtain a discrete partition with nearly the same stability.  
The same result holds for other domains, for example for $f,g : [k]^n \to [k]$.

\subsection{Relationship to rounding of SDPs}
We do not know of any connection between our results and those of Raghavendra and Steurer~\cite{RaghavendraSteurer:09} who showed that for any CSP, there is a rounding algorithm that is optimal up to $\epsilon$, whose running time is polynomial in the instance size and doubly exponential in $1/\epsilon$. 
It is natural to suspect that the two results are related, since there are well-known connections between SDP rounding and Gaussian noise stability.
However, the usual analysis relating rounding to Gaussian noise stability seems
to require that halfspaces maximize noise stability for all possible values of
the noise. In our results, this property does not necessarily hold.

In the other direction, it is tempting to try to cast the noise stability
problem as an optimization problem on a ``Gaussian graph:'' divide $\R^n$
into %
%
%
tiny pieces $A_1, \dots, A_m$, and consider the weighted  %
graph with vertices $1, \dots, m$ and an edge between vertices $i$ and $j$
that has weight $\E [1_{A_{i}} P_t 1_{A_j}]$ (where $1_A$ denotes the indicator
function of the set $A$). If the partition is fine enough, maximizing noise
stability is approximately equivalent to finding an appropriate optimal
partition of this graph, and so one might hope to apply the
results of Steurer and Raghavendra to obtain computable bounds on the dimension
where an almost optimal solution can be achieved. It is hard to implement this
approach for two reasons: first, we do not know the SDP solution for the
Gaussian graph; second, we are interested in the optimal solution and it is not
clear what is the relation between the best integral solution and the SDP
solution for the Gaussian graph. 

While we do not see how to formally relate the two works, connecting the two
(if possible) would be likely to yield important insights.

\subsection{Non-interactive correlation distillation}
Next, we discuss a basic problem in information theory and communication
complexity which was recently considered in the work of Ghazi, Kamath and
Sudan~\cite{GKS:16}.  Let there be two non-communicating players Alice and Bob
who have access to independent samples from a joint distribution $\bP$ on the set $\mathcal{A} \times \mathcal{B}$. In other words, Alice
(resp.\ Bob) has access to $(x_1, x_2, \ldots, )$ (resp.\ $(y_1, y_2, \ldots)$)
such that $x_i \in \mathcal{A}$, $y_i \in \mathcal{B}$ and for each $i \in
\mathbb{N}$, $(x_i, y_i)$ is distributed according to $\bP$, and the random
variables $\{(x_i, y_i)\}_{i=1}^n$ are mutually independent (as $i$ varies).
Let $\boldsymbol{\mu} = (\mu_1, \ldots, \mu_k) \in \Delta_k$ and
$\boldsymbol{\nu} = (\nu_1, \ldots, \nu_k) \in \Delta_k$.  What is the maximum
$\kappa \in [0,1]$ such that Alice and Bob can non-interactively jointly sample
a distribution $\bQ$ on $[k] \times [k]$ such that the distribution of the
marginals of Alice and Bob are $\boldsymbol{\mu}$ and $\boldsymbol{\nu}$,
%
respectively, and they sample the same output with probability $\kappa$? 

To formulate this problem more precisely, we introduce some notation.  Let $x^n
= (x_1, \dots, x_n)$ and $y^n = (y_1, \dots, y_n)$, where each pair $(x_i,
y_i)$ is drawn from $\bP$, independently with respect to $i$.  Now, note that a
non-interactive protocol for Alice and Bob is equivalent to a pair $(f,g)$
where 
$f: \mathcal{A}^n \rightarrow [k]$ and $g: \mathcal{B}^n \rightarrow [k]$ (for some $n \in \mathbb{N}$). 
In this terminology, the question now becomes the following: given
$\boldsymbol{\mu}$, $\boldsymbol{\nu}$, do there exist $n$ and $f:
\mathcal{A}^n \rightarrow [k]$ and $g: \mathcal{B}^n \rightarrow [k]$ such that
$f(x^n) \sim \boldsymbol{\mu}$,
$g(y^n) \sim \boldsymbol{\nu}$, and
$\Pr(f(x^n) = g(y^n)) = \kappa$?
(Here, $f(x^n) \sim \boldsymbol{\mu}$ means that $\boldsymbol{\mu}$ is the law
of $f(x^n)$.)

Before we state the main result of~\cite{GKS:16} and our extension, we consider
a motivating example. Let $\mathcal{A}= \mathcal{B}= \mathbb{R}$ and let $\bP$ be
the law of two $\rho$-correlated standard Gaussians. Let $k=2$ and
$\boldsymbol{\mu} = \boldsymbol{\nu}$.  Then, Borell's isoperimetric theorem
(\expref{Theorem}{thm:Borell})
states that the maximum achievable $\kappa$ is given by 
\[
\kappa = \Pr_{x,y}[f(x) = f(y)],
\]
where $f : \mathbb{R} \rightarrow \Delta_2$ is a halfspace with measure $\boldsymbol{\mu}$. Thus, in the above case, $n=1$ suffices and $f=g$ is the halfspace whose measure is $\boldsymbol{\mu}$. We now state the main result of~\cite{GKS:16}. For a probability distribution $\bP$, we let $|\bP|$ denote the size of some standard encoding of $\bP$. 

\begin{theorem}[{Ghazi-Kamath-Sudan}]~\label{thm:GKS}
    Let $(\mathcal{A} \times \mathcal{B}, \bP)$ be a probability space, fix $\delta > 0$,
    and let $x^n$ and $y^n$ be as above. 
There is an algorithm running in time $O_{| \bP|, \delta}(1)$ such that
given $\boldsymbol{\mu}$ and $\boldsymbol{\nu}$ in $\Delta_2$ and a parameter $\kappa \in [0,1]$, it distinguishes between the following two cases: 
\begin{enumerate}
  \item There exist $n \in \mathbb{N}$,
    $f: \mathcal{A}^n \rightarrow \{0, 1\}$, and
    $g: \mathcal{B}^n \rightarrow \{0, 1\}$ such that $f(x^n) \sim
    \boldsymbol{\mu}$, $g(y^n) \sim \boldsymbol{\nu}$, and $\Pr(f(x^n) =
    g(y^n)) \ge \kappa - \delta$. Moreover, 
    there is a computable function $n_0$ such that $ n \le n_0  
 ( \bP, \delta)$. Similarly, $f$ and $g$ can also be computed given $\bP$, $\boldsymbol{\mu}$, $\boldsymbol{\nu}$ and $\delta$ as inputs. 


  \item For any $n \in \mathbb{N}$ and $f: \mathcal{A}^n \rightarrow \{0, 1\}$
    and $g: \mathcal{B}^n \rightarrow \{0, 1\}$
    that satisfy $f(x^n) \sim
    \boldsymbol{\mu}$ and $g(y^n) \sim \boldsymbol{\nu}$, $\Pr(f(x^n) =
    g(y^n)) \le \kappa - 8\delta$.
\end{enumerate}
\end{theorem}
In other words, the above theorem states that there is an algorithm which, given $\bP$, target marginals $\boldsymbol{\mu}$, $\boldsymbol{\nu}$, and correlation $\kappa$, can distinguish between two cases: (a) In the first case, it is possible for Alice and Bob to non-interactively simulate a distribution $\bR$ which has the correct marginals and achieves the correlation $\kappa$ up to $\delta$. In this case, 
there is a computable bound on the number of copies of $\bP$ required and the algorithm also outputs the functions $f,g$ used for the non-interactive simulation. 
(b) In the second case, no non-interactive protocol between Alice and Bob can simulate a distribution $\bR$ such that it has the correct marginals and the correlation is more than $\kappa  - 8 \delta$.  

We remark that the requirement for the marginals to match \emph{exactly} is
unimportant. Indeed, any protocol where the marginals match approximately
can be ``fixed'' to have exact marginals, with a small loss in the correlation.

The main restriction of \expref{Theorem}{thm:GKS} is that the output of
the non-interactive protocol is a pair of bits. \expref{Theorem}{thm:main-stab}
immediately implies the following modification of \expref{Theorem}{thm:GKS}:
The distribution $\bP$ is a Gaussian measure and $\{0,1\}$ is replaced by $[k]$.
In the best of both worlds, we would be able to replace $\{0, 1\}$ by $[k]$
without adding the assumption that $\bP$ is a Gaussian measure.
Using the methods we develop here, this also turns out to be possible; this is done in a follow-up work~\cite{DMN:17b}. 

\section{Preliminaries} 
We will start by defining some technical preliminaries which will be useful for the rest of the paper. 
\begin{definition}
For $k \in \mathbb{N}$ and $1\le i \le k$, let $\mathbf{e}_i \in \R^k$ be the unit vector along coordinate $i$ and let $\Delta_k$ be the convex hull formed by $\{ \mathbf{e}_i \}_{1\le i \le k}$. 
\end{definition}
A function $f: \mathbb{R}^n \rightarrow \mathbb{R}^k$ is said to be in $L^2(\gamma_n)$ if $\int_x \Vert f(x)\Vert_2^2 \cdot \gamma_n(x) dx$ is finite.  In this paper, unless explicitly mentioned otherwise, (a) the domain is always equipped with the standard Gaussian measure $\gamma_n$ and (b) all functions are in $L^2(\gamma_n)$. 
 A key property of such functions is that they admit the so-called Hermite expansion. To define this notion, let us first define a family of polynomials $H_q: \mathbb{R} \rightarrow \mathbb{R}$ (for $q \ge 0$) as 
$$
H_0(x) = 1 ; \ H_1(x) = x ; \  H_q(x) = \frac{(-1)^q}{\sqrt{q!}}  \cdot \eee^{x^2/2}  \cdot \frac{d^q}{dx^q} \eee^{-x^2/2}. 
$$
This family of polynomials is referred to as the \emph{Hermite polynomials}. Let $\mathbb{Z}^{\ast}$ denote the subset of non-negative integers. For $S \in \mathbb{Z}^{\ast n}$, define $H_S: \mathbb{R}^n \rightarrow \mathbb{R}$ as
$$
H_S(z) = \prod_{i=1}^n H_{S_i}(z_i). 
$$
It is well known that the set $\{H_S\}_{S \in \mathbb{Z}^{\ast n}}$ forms an orthonormal basis for $L^2(\gamma_n)$. In other words, every $f \in L^2(\gamma_n)$
may be written as
$$
f = \sum_{S \in \mathbb{Z}^{\ast n}} \widehat{f}(S) \cdot H_S,
$$
where the sum converges respect to $L^2(\gamma_n)$. We call $\widehat{f}(S) = (\widehat{f_1}(s) \ldots, \widehat{f_k}(S)) \in
\mathbb{R}^k$ the \emph{Hermite coefficients}, and we refer to the formula above as the \emph{Hermite expansion} of $f$. 
Because the $H_S$ form an orthonormal basis, we obtain \emph{Parseval's identity}:
\begin{equation}\label{eq:parseval}
    \int_{\R^n} \Vert f (x) \Vert_2^2  \ \gamma_n(x) dx = \sum_{S \in \mathbb{Z}^{\ast n}} \Vert \widehat{f}(S) \Vert_2^2.
\end{equation}
Finally, for $f: \mathbb{R}^n \rightarrow \mathbb{R}^k$ and $d \in \mathbb{N}$, we define the degree-$d$ truncation of $f$ by ``throwing away''
all Hermite coefficients at levels higher than $d$:
$$
f_{\le d} (x) = \sum_{S: |S| \le d} \widehat{f}(S) \cdot H_S(x). 
$$
We also define $\mathsf{W}^{\le d} [f] = \Vert f_{\le d} \Vert_2^2$ and $\mathsf{W}^{> d} [f] = \sum_{|S|>d} \Vert \widehat{f}(S) \Vert_2^2$. 
\subsubsection*{Ornstein-Uhlenbeck operator} 
\begin{definition}
The Ornstein-Uhlenbeck  operator $P_t$ is defined for $t \in [0, \infty)$ and any bounded, measurable function $f: \mathbb{R}^n \rightarrow \mathbb{R}^k$ by
$$
(P_t f)(x) = \int_{y \in \mathbb{R}^n} f(\eee^{-t} \cdot x + \sqrt{1- \eee^{-2t}} \cdot y) d \gamma_n(y).
$$
%
\end{definition}
Note that if $f$ takes values in $\Delta_k \subset \R^k$, then so does $P_t f$ for every $t>0$. 
A basic fact about  the Ornstein-Uhlenbeck operator is that the functions $\{H_S\}$ are eigenfunctions of this operator. We leave the proof of the next proposition to the reader. 
\begin{proposition}
For $S \in \mathbb{Z}^{\ast n}$,
$P_t H_S = \eee^{-t \cdot |S|} \cdot H_S$. 
\end{proposition}
We define the noise stability of $f: \mathbb{R}^n \rightarrow \Delta_k$ with noise rate $t \in [0, \infty)$ by
\[
\mathsf{Stab}_t(f)  = \mathbf{E}[\langle f , P_t f \rangle]. 
\]
\subsubsection*{Multivariate polynomial threshold functions}
We recall the definition of polynomial threshold functions from the introduction (primarily to define an associated quantity which shall be useful later). 
\begin{definition*}
A function $f: \mathbb{R}^n \rightarrow [k]$ is said to be a multivariate PTF if there exist polynomials $p_1, \ldots, p_k : \mathbb{R}^n \rightarrow \mathbb{R}$ such that 
$$
f(x) = \begin{cases}
j  & \textrm{if} \  p_j(x)>0 \ \textrm{ and  for } i \neq j, \  p_{i}(x) \le 0, \\
1 & \textrm{otherwise.}\end{cases}
$$
We then denote $f$ as $f = \mathsf{PTF}(p_1, \ldots, p_k)$. Further, $f$ is
said to be a degree-$d$ multivariate PTF if   $p_1, \ldots, p_k$ are  degree-$d$ polynomials. For the rest of this paper, we will use ``PTF'' to mean
``multivariate PTF,'' and the underlying $k$ will be clear from the context. Also, let
us define $\mathsf{Collision} (f ) \subseteq \mathbb{R}^n$ defined as 
\[
\mathsf{Collision} (f) = \big\{x \in \mathbb{R}^n: |\{ j: p_j (x) > 0\}| \neq 1 \big\}. 
\]
\end{definition*}

For technical reasons, we will also require the PTFs to satisfy a certain regularity property:
\begin{definition} A degree-$d$ multivariate PTF $f = \mathsf{PTF}(p^{(1)}, \ldots, p^{(k)})$ is said to be $(d, \delta)$-balanced if each $p^{(i)}$ has variance $1$ and $|\mathbf{E}[p^{(i)}]| \le \log^{d/2} (k \cdot d/\delta)$. 
\end{definition}
Observe that the first condition (namely, $\mathsf{Var}(p^{(i)}) =1$) can be
achieved without loss of generality by simply rescaling all the polynomials to
have variance $1$. This rescaling does not change the value of the PTF at any
point $x$. While the condition on expectation is non-trivial, the next
proposition says that any multivariate PTF can be assumed to be
$(d,\delta)$-balanced while only changing the value of the PTF at
$\delta$-fraction of places. 
\begin{proposition}\label{prop:balanced}
Let $f : \mathbb{R}^n \rightarrow [k]$ satisfy $f = \mathsf{PTF}(p^{(1)},
\ldots, p^{(k)})$ for some
%
collection of
%
degree-$d$ polynomials $p^{(1)}, \dots, p^{(k)}$.
For any $\delta > 0$,
there is a $(d,\delta)$-balanced PTF $g = \mathsf{PTF}(q^{(1)}, \ldots, q^{(k)})$
defined on $\R^n$ such that $\Pr_{x} [g(x) \neq  f(x) ] \le \delta$ and
$\Pr_{x}[ x \in \mathsf{Collision}(g)] \le \Pr_{x} [x \in
\mathsf{Collision}(f)] + \delta$.
%
%
\end{proposition}
To prove this proposition, we will require the following standard concentration bound for low-degree polynomials which follows as a consequence of the well-known hypercontractive inequality---\eg, Lemma~2.2 in~\cite{DFKO06} proves the analogue of this statement for the Boolean hypercube. To get the statement for Gaussians, one can just substitute each one-dimensional Gaussian $x_i$ by $((y_{i,1} + \cdots + y_{i,m})/\sqrt{m})$ where each $y_{i,j}$ is an unbiased $\{\pm 1\}$ random variable and then take the limit as $m \rightarrow \infty$. %
\begin{theorem}\label{thm:hyper}
Let $p: \mathbb{R}^n \rightarrow \mathbb{R}$ be a degree-$d$ polynomial. Then, for any $t>0$, 
$$
\Pr_{x} \big[|p(x) - \mathbf{E}[p(x)]| \ge t \cdot \sqrt{\mathsf{Var}[p]}\big]  \leq d \cdot \eee^{-t^{2/d}}.
$$
\end{theorem}

\begin{proof}[Proof of \expref{Proposition}{prop:balanced}.]
As we have observed, we can assume without loss of generality that the polynomials $p^{(i)}$ have variance $1$. We define the polynomials $q^{(i)}$ as follows. If $|\mathbf{E}[p^{(i)}] | \leq \log^{d/2}(d \cdot k / \delta)$, then we set $q^{(i)} = p^{(i)}$. Else, let $b_i  = \mathsf{sign} (\mathbf{E}[p^{(i)}])$ and define 
\[
q^{(i)} = p^{(i)}  -\mathbf{E}[p^{(i)}]  + b_i \cdot \log^{d/2}(d \cdot k / \delta). 
\]
%
%
According to this definition, $\mathsf{sign}(q^{(i)}(x)) \ne \mathsf{sign}(p^{(i)}(x))$ implies
that $|q^{(i)}(x) - \E[q^{(i)}]| \ge \log^{d/2} (d k/\delta)$.
By \expref{Theorem}{thm:hyper},
\[
\Pr_{x \sim \gamma_n} \big[ \mathsf{sign}(q^{(i)}(x)) \neq  \mathsf{sign}(p^{(i)}(x)) \big] \le \delta/k. 
\]
Note that if $f(x) \ne g(x)$ then there is at least one $i$ for which $\mathsf{sign}(p^{(i)}(x)) \ne \mathsf{sign}(q^{(i)}(x))$.
By a union bound, it follows that $\Pr_{x \sim \gamma_n}[f(x) \ne g(x)] \le \delta$. Similarly, if $x \in \mathsf{Collision}(g)$ but
$x \not \in \mathsf{Collision}(f)$ then there is at least one $i$ for which $\mathsf{sign}(p^{(i)}(x)) \ne \mathsf{sign}(q^{(i)}(x))$.
Using another union bound,
\[
  \Pr_{x \sim \gamma_n} [ x \in \mathsf{Collision}(f) ] \le \Pr_{x}[x \in \mathsf{Collision}(g)] + \delta. \qedhere
  \]
\end{proof}

\subsubsection*{Convention} 
%
At  several places in the paper, we will establish bounds on one quantity in terms of others without the explicit mention of the dependence. For example, if we state that a quantity $d= d(k,\epsilon)$, we mean that that $d$ can be bounded as a computable function of $k$ and $\epsilon$ (even if the function $d(\cdot, \cdot)$ is simple, in the interest of clarity, we do not explicitly write it).  

%

\section{Proof strategy for \expref{Theorem}{thm:main-stab}}
We now formally state the main results required to prove \expref{Theorem}{thm:main-stab}. The first theorem says that given any $k$-ary function with a given set of measures for each of the $k$-partitions, there is a low-degree multivariate PTF which up to an error $\epsilon$ has 
(a)  the same measures for the induced partitions and (b) is no less noise stable at a fixed noise rate $t$.
In particular, we have the following theorem. 
\begin{theorem}\label{thm:1}
  Let $f: \mathbb{R}^n \rightarrow [k]$ satisfy $\E[f] = \boldsymbol{\mu}$ where $\boldsymbol{\mu} = (\mu_1, \ldots, \mu_k)$. Then, for every
  $\epsilon>0$ and $t >0$, there exists a multivariate PTF $g': \mathbb{R}^n \rightarrow
  [k]$ of degree $d = d(t, k, \epsilon)$ such that
  \begin{itemize}
    \item $\Vert \mathbf{E}[g'] - \boldsymbol{\mu} \Vert_1 \le \epsilon$,
    \item $\mathbf{E}[\langle g', P_t g' \rangle] \ge\mathbf{E}[ \langle f, P_t f \rangle] - \epsilon$,
    \item $\Pr [x \in \mathsf{Collision} (g')] \le \epsilon$, and
  \item $g'$ is $(d, \epsilon)$-balanced.
  \end{itemize}
\end{theorem}
The key difference between \expref{Theorem}{thm:1} and \expref{Theorem}{thm:main-stab} is that the number of variables in PTF $g'$ is the same as that in $f$. The above statement corresponds to the first step (``from general partitions to PTF'') of the proof of \expref{Theorem}{thm:main-stab} sketched in \expref{Section}{sec:theorem}. The next theorem statement formalizes the second and third steps sketched in \expref{Section}{sec:theorem}. In particular, it states that given a degree-$d$ multivariate PTF over $n$ variables, there is another multivariate PTF which induces approximately the same partition sizes and has approximately the same noise stability (at any fixed noise rate $t$) but the new PTF is only over some $O_{d,t}(1)$ variables.
\begin{theorem}\label{thm:reduction}
  Let $f: \mathbb{R}^n \rightarrow [k]$ be a degree-$d$, $(d,\epsilon)$-balanced  PTF with $\E_{x}[f(x)] = \boldsymbol{\mu}$ where $\boldsymbol{\mu} = (\mu_1, \ldots, \mu_k)$.  Further, let us assume that 
$\Pr[x \in \mathsf{Collision} (f)] \le \epsilon/(40 k^2)$.  
  Then, for every $\epsilon>0$ and $t > 0$, there exists a degree-$d$ 
PTF $\fh: \mathbb{R}^{n_0} \rightarrow [k]$ such that, \begin{itemize} \item  $\Vert \mathbf{E}_{x} [\fh(x)] - \boldsymbol{\mu} \Vert_1 \le \epsilon$, \item $ \mathbf{E}[\langle \fh, P_t \fh \rangle] \ge \mathbf{E}[\langle f, P_t f \rangle] - \epsilon$. %
\end{itemize} Further, $n_0 = n_0 (d,k,\epsilon, t)$ is a computable function.   
\end{theorem}
We now see how \expref{Theorem}{thm:1} and \expref{Theorem}{thm:reduction} yield \expref{Theorem}{thm:main-stab}. 
\subsubsection*{Proof of \expref{Theorem}{thm:main-stab}} 
Let us  assume that given measure $\boldsymbol{\mu} = (\mu_1, \ldots, \mu_k)$ and noise rate $t>0$, the most
%
noise-stable    %
partition is $f: \mathbb{R}^n \rightarrow [k]$. Then, applying \expref{Theorem}{thm:1} (with error parameter $\epsilon/(80k^2)$), we obtain a PTF $g_1$ of degree $d = O_{t, k,\epsilon}(1)$ such that it is $(d,\epsilon/2)$ balanced and $\Pr[x \in \mathsf{Collision}(g_1)]\le \epsilon/(80k^2)$. Further, note that $\Vert \mathbf{E}[g_1] - \mathbf{E}[f] \Vert_1 \le \epsilon/2$ and $\mathbf{E}[\langle g_1, P_t g_1 \rangle] \ge\mathbf{E}[ \langle f, P_t f \rangle] - \epsilon/2$. 

We next apply \expref{Theorem}{thm:reduction} to function $g_1$ to obtain $\fh$ which is a degree-$d$ PTF on $n_0 = O_{t,k, \epsilon}(1)$ variables such that $\Vert \mathbf{E}[g_1] - \mathbf{E}[\fh] \Vert_1 \le \epsilon/2$ and 
$\mathbf{E}[\langle \fh, P_t \fh \rangle] \ge \mathbf{E}[\langle g_1, P_t g_1 \rangle] - \epsilon/2$. 

Combining these two facts, we obtain $\Vert \mathbf{E}[g_1] - \mathbf{E}[\fh] \Vert_1 \le \epsilon$ and 
$\mathbf{E}[\langle \fh, P_t \fh \rangle] \ge \mathbf{E}[\langle f, P_t f \rangle] - \epsilon$. Setting $g = \fh$ concludes the proof. 

\section{Reduction from arbitrary functions to PTFs: Proof of \expref{Theorem}{thm:1}}\label{sec:reduction}
In this section, we will prove \expref{Theorem}{thm:1} in two parts.
The first step is to show that
we can replace $f$ by a function with a lower bound on the Hermite decay:
\begin{lemma}\label{lem:1} 
  For every $f: \mathbb{R}^n \rightarrow [k]$, every $\epsilon>0$, and every $t > 0$, there exists a function
 $h: \mathbb{R}^n \rightarrow [k]$ satisfying the following: 
 \begin{itemize}
   \item $\E[h] = \E[f]$,
 \item $\mathbf{E}[\langle h, P_t h \rangle] \ge\mathbf{E}[ \langle f, P_t f \rangle ]- \epsilon$, 
 \item $\mathsf{W}^{> d} [h] \le \epsilon$ for a computable $d = d(t, k, \epsilon)$. 
 \end{itemize}
\end{lemma} 
The second step in the proof of \expref{Theorem}{thm:1} goes from the
Hermite decay to an actual PTF. 
\begin{lemma}\label{lem:2}
Let $h : \mathbb{R}^n \rightarrow [k]$ satisfy $\mathsf{W}^{>d}[h] \le \epsilon$. Then, there is a PTF $g: \mathbb{R}^n \rightarrow [k]$ of degree $d$ such that 
$ \mathbf{E}_{x } [g(x) \neq  h(x) ] \le k^2 \cdot  \epsilon$ and $\Pr_{x } [x \in \mathsf{Collision} (h)]  \le k^2 \cdot \epsilon$.
\end{lemma}
It is easy to see that \expref{Theorem}{thm:1} follows in a straightforward way by combining \expref{Lemma}{lem:1} and \expref{Lemma}{lem:2}. 
Note that the condition of $g$ being $(d,\epsilon)$-balanced can be obtained by  applying \expref{Proposition}{prop:balanced}.
We begin with the proof of \expref{Lemma}{lem:2}, which is easier.

\begin{proof}[Proof of \expref{Lemma}{lem:2}]
    By identifying $i \in [k]$ with $\mathbf{e}_i$ as before, we view $h$ as a
    function taking values in $\R^k$. Define $h^{(0)} : \mathbb{R}^n
    \rightarrow \mathbb{R}^k $ by $h^{(0)} (x) = h(x) -({1}/{k}) \cdot
    \mathbf{1}$, where $\mathbf{1} = (1, \ldots, 1) \in \R^k$. Note that for
    every $x \in \R^n$, $h^{(0)}(x)$
    is a $k$-dimensional vector with
    the entry $(k-1)/k$ in one of the coordinates and $-1/k$ in every other
    coordinate. Call such a point \emph{a $k$-lattice point}.  

Let $h^{(0)}_{\le d} : \mathbb{R}^n \rightarrow \mathbb{R}^k$ be the function
obtained by truncating the Hermite expansion of $h^{(0)}$ at degree $d$.
Consider the degree-$d$ polynomials $p_1, \ldots, p_k: \mathbb{R}^n \rightarrow
\mathbb{R}$ where $p_i$ is the $i$-{th} coordinate of $h^{(0)}_{\le d}$, and
define $g = \mathsf{PTF}(p_1, \ldots, p_k)$. We claim that $g$ satisfies
the conditions of the lemma.

Consider any point $x \in
\mathbb{R}^n$. 
\begin{itemize}
\item If $x \in \mathsf{Collision}(g)$, we claim that $\Vert h^{(0)} (x) -
    h^{(0)}_{\le d}(x)  \Vert_2^2 \ge k^{-2}$. To see this, note that if $x \in
    \mathsf{Collision}(g)$, there are two possibilities: either $p_i(x) \le 0$
    for all $1 \le i \le k$ or there are two distinct $i,j$ such that
    $p_i(x) > 0$ and $p_j(x) > 0$. In either case, it
    easily follows that the $\ell_2$ distance from
    $h^{(0)}_{\le d}(x)$ to any $k$-lattice point is at least $1/k$.

\item If $x \not \in \mathsf{Collision}(g)$ and $g(x) \neq  h(x)$, then again
    we claim that $\Vert h^{(0)} (x) - h^{(0)}_{\le d}(x)  \Vert_2^2 \ge
    k^{-2}$. To see this, suppose that $g(x) = \mathbf{e}_i \ne \mathbf{e}_j = h(x)$.
    Then the $j$th coordinate of $h^{(0)}(x)$ is $(k-1)/k$, but on the other hand
    $p_j(x) \le 0$ because $g(x) = \mathbf{e}_i$ and $x \not \in \mathsf{Collision}(g)$.
    Hence,
    $\Vert h^{(0)} (x) - h^{(0)}_{\le d}(x)\Vert_2^2 \ge {(k-1)^2}/{k^2}$. 
\end{itemize}
Combining these two, we get 
\[
\mathbf{E}_{x } \big[ \Vert h^{(0)} (x) - h^{(0)}_{\le d}(x)\Vert_2^2 \big] \geq \frac{1}{k^2} \cdot  \big( \Pr_{x } [x \in \mathsf{Collision}(g)] +  \Pr_{x } [x \not \in \mathsf{Collision}(g) \ \wedge \ g(x) \neq  h(x) ]\big). 
\]
As $h^{(0)}$ and $h$ differ only in the Hermite coefficient for $S = \emptyset$, the left hand side is the same as $ \mathsf{W}^{\ge d}[h]$. Combining this with the above inequality achieves all the stated guarantees. 
\end{proof}

The proof of \expref{Lemma}{lem:1} is longer, and so we begin with an outline.
The first observation is that
bounding $\E[|\nabla h|]$ implies a bound on $\mathsf{W}^{> d}[h]$.
This follows a standard spectral argument, and is stated as \expref{Corollary}{corr:gradient}.
The second observation is that the function $h$ obtained by thresholding
$P_t f$ at a suitable value (chosen, for example, so that $\E[h] = \E[f]$)
satisfies $\mathbf{E}[\langle h, P_t h\rangle] \ge \mathbf{E}[ \langle f, P_t f\rangle]$. This is stated as \expref{Lemma}{lem:rounding-stability}.

Based on the previous paragraph, it seems like we would like to bound $\E[|\nabla h|]$ where $h$ is obtained
by thresholding $P_t f$; let $a \in \R^k$ be the desired threshold value, so that thresholding $P_t f$ at $a$
produces a partition with the right measures. It turns out, unfortunately, that for $h$ defined in this way,
$\E[|\nabla h|]$ could be
arbitrarily large. A key insight of~\cite{KNOW14} is that (using the co-area formula and gradient bounds on $P_t f$)
the partition produced by thresholding at a random value near $a$ has bounded expected surface area. In particular,
although thresholding exactly at $a$ might be a bad idea,
there exist many good nearby values at which to threshold.
Based on this observation, we construct $h$ in two steps. In the first step, we define $h$ by thresholding $P_t f$,
but only on the set of $x \in \R^n$
for which $P_t f(x)$ is not too close to $a$. By choosing ``not too close'' in a suitable random way, the observation
of~\cite{KNOW14} implies that this step only contributes a bounded amount to $\E[|\nabla h|]$. Since the first step
is almost the same as just thresholding $P_t f$ at $a$, it is consistent with our desire that
$\mathbf{E}[\langle h, P_t h\rangle ]\ge \mathbf{E}[\langle f, P_t f\rangle] - \epsilon$.

In the second step, we partition the remaining part of $\R^n$ by chopping it with halfspaces of the correct size.
Since halfspaces have a bounded surface area, this also contributes a bounded amount to $\E[|\nabla h|]$.
Crucially, this step does not
destroy the value of $\langle h, P_t h\rangle$; fundamentally, this is because $P_t f$ is almost constant on the set
we are partitioning. This finishes the outline of the proof of \expref{Lemma}{lem:1}. We now start with some preliminaries required in the proof of this lemma. 

\subsection{Surface area and spectrum}

In our outline of the proof of \expref{Lemma}{lem:1}, we claimed that a control of
$\E[|\nabla h|]$ implies a control of $\mathsf{W}^{> d}[h]$. Here we prove that
claim, using a theorem of Ledoux~\cite{Ledoux:94} that gives a lower bound on
the gradient in terms of the noise sensitivity. 

%
\begin{theorem}[Ledoux]  %
  \label{thm:Ledoux}
    For a differentiable function $f: \mathbb{R}^n \rightarrow [0,1]$ and any $t > 0$,
\[
\mathop{\mathbf{E}}_{x \sim \gamma_n} [f \cdot (f - P_t f)]  \le \frac{\arccos (\eee^{-t})}{\sqrt{2\pi}} \cdot \mathop{\mathbf{E}}_{x \sim \gamma_n}[|\nabla f|]. 
\]
\end{theorem}
Using this theorem, it is possible to establish an upper bound on $\mathsf{W}^{\ge d}[f]$ in terms of $|\nabla f|$ as done below. 
\begin{corollary}\label{corr:gradient}
    For any differentiable $f: \mathbb{R}^n \rightarrow [0,1]$ be of bounded variation (in Gaussian measure) and any $d \in \mathbb{N}$,
\[
\mathsf{W}^{\ge d} [f]  \leq O\bigg( \frac{1}{\sqrt{d}}\bigg) \cdot \mathbf{E}[|\nabla f|].
\]
\end{corollary}
\begin{proof}
\begin{eqnarray*}
\mathbf{E}[ f (f- P_t f)] = \sum_{j=0}^{\infty} (1- \eee^{-t \cdot j}) \mathsf{W}^{j}[f] \ge (1-\eee^{-d \cdot t}) \cdot  \mathsf{W}^{ \ge d}[f].
\end{eqnarray*}
Choose $d=1/t$ and apply \expref{Theorem}{thm:Ledoux}, we get 
\[
\mathsf{W}^{ \ge d}[f] \leq O(1) \cdot \mathbf{E}[ f (f- P_{\frac{1}{d}} f)] \le  \frac{\arccos (\eee^{-\frac1d})}{\sqrt{2\pi}} \cdot \mathop{\mathbf{E}}_{x \sim \gamma_n}[|\nabla f|] \le O\bigg( \frac{1}{\sqrt{d}}\bigg) \cdot \mathbf{E}[|\nabla f|].\qedhere
\]
\end{proof}

\subsection{Various lemmas for thresholding}

Recall that our proof of \expref{Lemma}{lem:1} is based on thresholding $P_t f$. Before proving it, we
introduce various properties of the thresholding procedure.
The first lemma states that the ``best'' way to round a $\Delta_k$-valued function to
a $\{\mathbf{e}_1, \dots, \mathbf{e}_k\}$-valued function while preserving its expectation is simply to threshold it.
Here, ``best'' means that we are trying to maximize the correlation between the original function
and the rounded one.

\begin{lemma}\label{lem:rounding-correlation}
  Let $f: \R^n \to \Delta_k$, $g: \R^n \to \{\mathbf{e}_1, \dots, \mathbf{e}_k\}$, and $z \in \R^k$.
  Suppose that whenever $g(x) = \mathbf{e}_i$, we have
  \[
    i \in \mathrm{argmax}\{f_j(x) - z_j: j = 1, \dots, k\}.
  \]
  Then $g$ maximizes $\E[\langle f, h \rangle]$ among all
  $h: \R^n \to \{\mathbf{e}_1, \dots, \mathbf{e}_k\}$ satisfying $\E [h] = \E[g]$.
\end{lemma}

\begin{proof}
  Let $h: \R^n \to \{\mathbf{e}_1, \dots, \mathbf{e}_k\}$ satisfy $\E[h] = \E[g]$. By the defining assumption of $g$,
  $\inr{g}{f-z} \ge \inr{h}{f-z}$ pointwise. Hence,
  \[
    \E[\inr{h}{f}] = \E[\inr{h}{f-z}] + \inr{\E[h]}{z} \ge \E[\inr{g}{f-z}] + \inr{\E[g]}{z} = \E[\inr{g}{f}]. \qedhere
  \]
\end{proof}

Let us assign a notation to the sort of thresholding involved in \expref{Lemma}{lem:rounding-correlation}.

\begin{definition}
  For $f: \R^n \to \Delta_k$ and $z \in \R^k$, let $\calF_z(f)$ denote the set of functions $g: \R^n \to \{\mathbf{e}_1, \dots, \mathbf{e}_k\}$
  such that
  \[
      \langle g(x), f(x) - z \rangle = \max_j \{f_j(x) - z_j\}
  \]
  for every $x \in \R^n$.
\end{definition}

We say that functions in $\calF_z(f)$ are obtained by ``thresholding'' $f$ around the value $z$.

The next step is to show that a function obtained by thresholding $P_t f$ is always at least as noise-stable (with parameter $t$)
as the original function $f$. 

\begin{lemma}\label{lem:rounding-stability}
  Let $f: \R^n \to \{\mathbf{e}_1, \dots, \mathbf{e}_k\}$ and take $t > 0$. If $g \in \calF_z(f)$ for some $z \in \R^k$
  satisfies $\E[g] = \E[f]$
  then $\Stab_t(g) \ge \Stab_t(f)$.
\end{lemma}

\begin{proof}
  Since $\E[g] = \E[f]$, \expref{Lemma}{lem:rounding-correlation} implies that
  $\E[\langle g, P_t f\rangle] \ge \E[\langle f, P_t f\rangle] = \Stab_t(f)$. On the other hand,
  the Cauchy-Schwarz inequality and the semi-group property of $P_t$ imply that
  \[
    \E[\langle g, P_t f \rangle]
    = \E[\langle P_{t/2} g, P_{t/2} f \rangle]
    \le \sqrt{\Stab_t(g) \Stab_t(f)}.
  \]
  The claim follows.
\end{proof}

Next, we come to the importance of the thresholding parameter $z$: it allows us to tune the expectation of the thresholded functions. In particular, in order to make \expref{Lemma}{lem:rounding-stability} useful we need to show that we can always find a thresholded function with the same expectation as the original function.

\begin{lemma}\label{lem:rounding-existence}
  For any $f: \R^n \to \Delta_k$, there exists some $z\in \R^k$ and $g \in \calF_z(f)$ such that $\E[g] = \E[f]$.
\end{lemma}

It will actually be convenient to prove a more general Lemma about splitting up the mass of a general probability measure.
Given $z \in \R^k$ and $i \in [k]$, let $A_i(z) \subset \Delta_k$ be the set $\{x \in \Delta_k: x_i - z_i = \max_j x_j - z_j\}$. The correspondence between this notion and the definition of $\calF_z(f)$ is that $g \in \calF_z(f)$ if and only if $g(x) = \mathbf{e}_i$ implies $f(x) \in A_i(z)$.

\begin{lemma}\label{lem:rounding-prob-measure}
  For every probability measure $\zeta$ on $\Delta_k$ and every $q \in \Delta_k$,
  there exists $z \in \R^k$ such that $\zeta(\bigcup_{i \in I} A_i(z)) \ge \sum_{i \in I} q_i$ for every $I \subset [k]$.
\end{lemma}

Before proving \expref{Lemma}{lem:rounding-prob-measure}, we will show how it implies \expref{Lemma}{lem:rounding-existence}.
For this, we introduce one more ingredient: a weighted version of Hall's marriage theorem.
It will be used to show that we can divide a certain amount of ``missing'' volume into pieces of the right size,
while preserving certain constraints.

\begin{lemma}\label{lem:hall}
  Let $G = (U, V, w, E)$ be a finite, vertex-weighted, bipartite graph. That is, $U$ and $V$ are finite sets, $w$ is a weight function $w: U \sqcup V \to (0, \infty)$, and $E \subset U \times V$ is the
  set of edges. Suppose that for every $U' \subset U$, $\sum_{u \in U} w(u) \le \sum_{v \sim U'} w(v)$, where $v \sim U'$ if there exists $u \in U'$ such that $(u, v) \in E$. Then
  there exists $p: V \to \Delta_{|U|}$ such that $p_u(v) > 0$ implies $u \sim v$, and such that
  \[
    \sum_{v \in V} p_u(v) w(v) = w(u)
  \]
  for every $u \in U$.
\end{lemma}
In the case that $w$ takes values in $\{0, 1\}$, \expref{Lemma}{lem:hall} follows from the usual formulation of Hall's marriage theorem
(which guarantees in addition that $p: V \to \{\mathbf{e}_1, \dots, \mathbf{e}_{|U|}\}$). When
$w$ is integer-valued (or, by scaling, rational-valued), it follows by applying Hall's marriage
theorem to the graph in which $u$ is replicated $w(u)$ times. The general case follows by
a simple approximation argument.


\begin{proof}[Proof of \expref{Lemma}{lem:rounding-existence}]
    Let $\zeta$ be the distribution of $f$ and let $q = \E[f]$. Applying \expref{Lemma}{lem:rounding-prob-measure},
  we may take $z \in \R^k$ such that
  $\zeta(\bigcup_{i \in I} A_i(z)) \ge \sum_{i \in I} q_i$ for every $I \subset [k]$. For
  $I \subset [k]$, let $B_I$ be the set of $x \in \R^n$ such that $\{i: f(x_i) - z_i = \max_j f(x_j) - z_j\} = I$. Then $\{B_I: I \subset [k]\}$ is a
  partition of $\R^n$. Moreover, if we define $w(I) = \gamma_n(B_I)$ then
  \begin{equation}\label{eq:hall-assumption}
    \sum_{J: J \cap I \ne \emptyset} w(J) = \gamma_n\Big(\bigcup_{J: J \cap I \ne \emptyset} B_J\Big) = \zeta\Big(\bigcup_{i \in I} A_i(z)\Big) \ge \sum_{i \in I} q_i.
  \end{equation}
  Now we construct a bipartite graph with $U = [k]$ and $V = \{I: I \subset [k]\}$. The pair $(i, I)$
  is an edge if $i \in I$. We define $w$ on $V$ as above, and we define $w$ on $U$ by $w(i) = p_i$.
  According to~\eqref{eq:hall-assumption}, the condition of \expref{Lemma}{lem:hall}
  holds; hence, there exists $p: V \to \Delta_k$ such that $p_i(I) > 0$ only when $i \in I$, and
  such that $\sum_{I \ni i} p_i(I) \gamma_n(B_I) = q_i$. Now we will define $g$: for each 
  $I \subset [k]$, partition $B_I$ arbitrarily into sets $\{B_{I,i}: i \in I\}$, where
  $\gamma_n(B_{I,i}) = p_i(I) \gamma_n(B_I)$. This may be done because $\gamma_n$ has no atoms.
  Finally, define $g$ to be $e_i$ on every $B_{I,i}$.
  Then the condition $\sum_{I \ni i} p_i(I) \gamma_n(B_I) = q_i$ ensures that $\gamma_n(\{g = e_i\}) = q_i$.
  Moreover, note that $g(x) = e_i$ implies that $x \in B_I$ for some $I$, which implies
  that $f_i(x) - z_i = \max_j (f_j(x) - z_j)$. This completes the construction of $g$.
\end{proof}

\begin{proof}[Proof of \expref{Lemma}{lem:rounding-prob-measure}]
  We will assume initially that $\zeta$ is absolutely continuous with respect to the
  uniform measure on $\Delta_k$. As a consequence, the function $\psi: \Delta_k \to \R^k$ defined by
  $\psi_i(z) = \zeta(A_i(-kz))$
  is continuous. Moreover, the image of $\psi$ is in $\Delta_k$, and we need prove that it is all of $\Delta_k$.
  
  For $I \subsetneq \{1, \dots, k\}$, let $F_I = \{x \in \Delta_k: x_i = 0 \text{ for all } i \in I\}$. Note that
  every face of $\Delta_k$ is of the form $F_I$ for some $I \subsetneq \{1, \dots, k\}$.
  Next, we claim that $\psi$ maps $F_I$ into $F_I$ for every $I$. Indeed, if $z \in F_I$ then there is at least one $j \not \in I$
  such that $z_j \ge 1/k$. For this $j$ and any $i \in I$, if $x$ is in the interior of $\Delta_k$ then
  \[
    x_i + k z_i = x_i < x_j + k z_j.
  \]
  It follows that $A_i(-kz)$ does not intersect the interior
  of $\Delta_k$; hence, $\psi_i(z) = 0$. Since this holds for all $i \in I$, $\psi(z) \in F_I$.
  By Sperner's lemma, any map from $\Delta_k$ into itself that leaves all faces invariant must be onto, and so
  $\psi$ is onto, as claimed.

  To complete the proof, we must eliminate the assumption that $\zeta$ is absolutely continuous with
  respect to the uniform measure. For an arbitrary probability measure $\zeta$, let $\zeta_n$ be a sequence of
  probability measures that are absolutely continuous with respect to the uniform measure, and which converge to $\zeta$ in distribution. (For example, we may construct $\zeta_n$ by first pushing $\zeta$ forward under the map $x \mapsto ({1}/{2n}) (1/k, \dots, 1/k) + (1 - 1/n) x$, and by convolving the push-forward with a smooth bump function supported on $({1}/{2n}) \Delta_2$.) Define $\psi_n: \Delta_k \to \Delta_k$
  by $\psi_n(z) = (\zeta_n(A_1(z)), \dots, \zeta_n(A_k(z)))$.
  By the previous argument, for every $n$ there is some $z_n \in \Delta_k$ such that $\psi(z_n) = q$. Since
  $\Delta_k$ is compact, we may pass to a subsequence and thereby assume that $z_n$ converges to some limit $z_\infty$.
  Note that
  \[
    \bigcup_{i \in I} A_i(z_\infty) = \bigcap_{n=1}^\infty \Big(\bigcup_{i \in I} \bigcup_{m \ge n} A_i(z_m) \cup A_i(z_\infty)\Big).
  \]
  Moreover, $\bigcup_{m \ge n} A_i(z_n) \cap A_i(z_\infty)$ is closed for every $n$, since $z_\infty$ is the only
  limit point of the sequence $z_n$. It follows that
  \[
    \mu\Big(\bigcup_{i \in I} A_i(z_\infty)\Big)
    = \lim_{n \to \infty} \mu\Big(\bigcup_{i \in I} \bigcup_{m \ge n} A_i(z_m) \cup A_i(z_\infty)\Big)
    \ge \lim_{n \to \infty} \lim_{\ell \to \infty} \mu_\ell\Big(\bigcup_{i \in I} \bigcup_{m \ge n} A_i(z_m) \cup A_i(z_\infty)\Big)
    \ge \sum_i q_i,
  \]
  since $\sum_{i \in I} \mu_\ell(A_i(z_\ell)) = \sum_{i \in I} q_i$ for every $I$ and $\ell$
  (using the fact that $\mu$ has a density, and so it assigns no mass to $A_i(z_\ell) \cap A_j(z_\ell)$).
  \end{proof}

%

\subsection{Proof of \expref{Lemma}{lem:1}}

We introduce two final ingredients before beginning the proof of \expref{Lemma}{lem:1}: the first is a gradient bound on $P_t f$
that is due (in a much more precise form) to Bakry and Ledoux~\cite{BakryLedoux:96}.
\begin{theorem}\label{thm:bakry-ledoux}
  If $f: \R^n \to [0, 1]$ then for any $t > 0$, $P_t f$ is differentiable and $\|\nabla P_t f\|_\infty \le C/\sqrt{t}$.
\end{theorem}

The second ingredient is the co-area formula in Gaussian space. To state this, let $\gamma_n^{+}$ denote the Gaussian surface area in $n$-dimensions. The co-area formula (see~\cite{federer1969} for a reference) is stated next. For any locally Lipschitz $f: \mathbb{R}^n \rightarrow \mathbb{R}$ and for any continuous, compactly supported $\mu : \mathbb{R} \rightarrow \mathbb{R}$ 
\[
\int_{\mathbb{R}} \mu(t) \cdot \gamma_n^{+} (\{x: f(x) \ge t\}) dt = \int_{\mathbb{R}^n} \mu(f(x)) |\nabla(f(x))| d \gamma(x).
\]
The co-area formula relates the Gaussian surface of a Boolean function (specified by $\{x: f(x) \ge t\}$) to the gradient of $f$. We remark that the co-area formula is sometimes written with $\gamma_n^+(\{x: f(x) \ge t\})$ replaced by the Gaussian-weighted $(n-1)$-dimensional Hausdorff measure of $\{x: f(x) = t\}$. The two formulations are equivalent, because these two quantities are equal for almost every $t$ (see, for example,~\cite[Chapter 13]{Maggi:12}).


\begin{proof}[Proof of \expref{Lemma}{lem:1}]
We begin by applying \expref{Lemma}{lem:rounding-existence} to $P_t f$: let $z$ and $g \in \calF_z(P_t f)$
be such that $\E[g] = \E[P_t f] = \E[f]$. According to \expref{Lemma}{lem:rounding-stability}, $\Stab_t(g) \ge \Stab_t(f)$.
However, $\E[|\nabla g|]$ cannot be controlled in general; therefore, we will move to an approximation of $g$.

Let $\mu$ be a probability measure on $[0, \epsilon]$ with continuous density bounded by $2/\epsilon$.
By the co-area formula, for any $i \ne j$ we have
\begin{align*}
  &\int_\R \mu(t) \gamma^+(\{x: P_t f_i(x) - z_i > P_t f_j(x) - z_j + t\})\, dt \\
  &= \int_{\R^n} \mu(P_t f_i(x) - z_i - P_t f_j(x) + z_j) |\nabla P_t f_i(x) - \nabla P_t f_j(x)|\, d\gamma(x) \\
  &\le \frac{2}{\epsilon} \int_{\R^n} |\nabla P_t f_i(x) - \nabla P_t f_j(x)|\, d\gamma(x) \\
  &\le \frac{C}{\epsilon \sqrt t},
\end{align*}
where the last inequality follows from \expref{Theorem}{thm:bakry-ledoux}.
In particular, there exist some $y_{ij}^+$ and $y_{ij}^-$ in $[0, \epsilon]$ such that if
$A_{ij}^{\pm} = \{x \in \R^n: P_t f_i(x) - z_i > P_t f_j(x) - z_j \pm y_{ij^{\pm}}\}$ then
\[
\gamma^+(A_{ij}^\pm) \le C(\epsilon)/\sqrt t.
\]
We repeat this construction of $y_{ij}$ and $A_{ij}$ for every ordered pair $(i, j)$ with $i \ne j$.
Define $A_i = \bigcap_{j \ne i} A_{ij}^+$, and note that $A_i \subset \{x: g(x) = \mathbf{e}_i\}$.
On the other hand, $A_{ij}^- \supset \{x: g(x) = \mathbf{e}_i\}$ for every $i, j$.
Next, define
\begin{align*}
  C_i &= \bigcap_{j \ne i} A_{ij}^-, \\
  C_I &= \bigcup_{i \in I} C_i, \\
  B_I &= C_I \setminus \bigcup_{J \supsetneq I} C_J \setminus \bigcup_{i \in I} A_i.
\end{align*}
The meaning of these sets is the following: $A_i$ is the set where $f_i - z_i$ is significantly larger than
any other
$f_j - z_j$. On $C_i$, $f_i - z_i$ is almost $\max_j f_j - z_j$;
on $C_I$, $f_i - z_i$ is almost maximal for \emph{every} $i \in I$; and on $B_I$, the set
of $i$ for which $f_i - z_i$ is almost maximal is exactly $I$. Importantly, the collection
of all $A_i$ and $B_I$ form a partition of $\R^n$. Our basic strategy will be to set $h$ to be $\mathbf{e}_i$
on $A_i$, and then to define $h$ on the remaining part of the space in order to satisfy two properties:
$\E[h] = \E[f]$ and $h(x) = \mathbf{e}_i$ only if $x \in B_I$ for some $I \ni i$.

Since the Gaussian surface area obeys the inequalities $\gamma_n^+(A \cap B) \le \gamma_n^+(A) + \gamma_n^+(B)$
and $\gamma_n^+(A \cup B) \le \gamma_n^+(A) + \gamma_n^+(B)$, and since $B_I$ and $A_i$ are defined
using a finite (depending on $k$) number of intersections and unions, it follows that
\begin{equation}\label{eq:A-B-surface}
  \begin{aligned}
    \gamma^+(A_i) &\le C(\epsilon, k)/\sqrt t, \\
    \gamma^+(B_I) &\le C(\epsilon, k)/\sqrt t.
  \end{aligned}
\end{equation}

Now, for any $I \subset [k]$,
\[
  \bigcup_{J: J \cap I \ne \emptyset} B_J \:\cup\: \bigcup_{i \in I} A_i
\]
contains the set of $x$ for which $g(x) \in \{\mathbf{e}_i: i \in I\}$. It follows that
\begin{equation}\label{eq:hall-assumption-2}
  \sum_{J: J \cap I \ne \emptyset} \gamma_n(B_J)
  \ge \sum_{i \in I} \E [g_i] - \gamma_n(A_i).
\end{equation}
Consider the bipartite graph where $U = [k]$ and $V = \{I: I \subset [k]\}$, and $(i, I) \in E$ if $i \in I$.
We assign the weights $w(i) = \E[g_i] - \gamma_n(A_i)$ to $i \in U$ and $w(I) = \gamma_n(B_I)$ for $I \in V$
(note that $w(i) \ge 0$ because the construction of $A_i$ ensures that $A_i \subset \{x: g(x) = \mathbf{e}_i)\}$).
Equation~\eqref{eq:hall-assumption-2} ensures that this weighted graph satisfies the hypothesis of \expref{Lemma}{lem:hall},
and so there exists $p: V \to \Delta_k$ with $p_i(I) > 0$ only if $i \in I$, and with
\[
  \sum_{I \ni i} p_i(I) \gamma_n(B_I) = \E[g_i] - \gamma_n(A_i).
\]

Finally, we will use $p$ to define $h$. First, for every $I$ and $i \in I$, let $B_{I,i}$ be a set of the form
$\{x \in \R^n: a \le x_1 \le b\}$ such that $\gamma_n(B_{I,i} \cap B_I) = p_i(I) \gamma_n(B_I)$; moreover, we choose $B_{I,i}$
such that $\gamma_n(B_{I,i} \cap B_{I,j}) = 0$ when $i \ne j$. Then we
set $h(x)$ to equal $\mathbf{e}_i$ on the set $A_i \cup \bigcup_{I \ni i} (B_{I,i} \cap B_I)$. By the defining property of $p$,
$h$ satisfies $\E[h_i] = \E[g_i] = \E[f_i]$. Moreover, $h(x) = \mathbf{e}_i$ only on a subset of $A_i \cup \bigcup_{I \ni i} B_I$,
and on this set $P_t f_i(x) - z_i \ge \max_j P_t f_j(x) - z_j - \epsilon = \langle g(x), P_t f(x) - z\rangle - \epsilon$.
It follows that
\[
  \E[\langle h, P_t f\rangle] \ge \E[\langle g, P_t f\rangle] - \epsilon \ge \Stab_t(f) - \epsilon.
\]
By the Cauchy-Schwarz inequality,
\[
  \sqrt{\Stab_t(h)} \ge \sqrt{\Stab_t(f)} - \frac{\epsilon}{\sqrt{\Stab_t(f)}},
\]
which implies that $\Stab_t(h) \ge \Stab_t(f) - 2\epsilon$.

Finally, we address the surface area of $h$.  Recall that
\[
  \{x: h(x) = \mathbf{e}_i\} = A_i \cup \bigcup_{I,i} B_{I,i} \cap B_I.
\]
The number of terms on the right is some constant depending on $k$. By~\eqref{eq:A-B-surface},
$\gamma_n^+(A_i)$ and $\gamma_n^+(B_I)$ are bounded by $C(\epsilon, t, k)$. Since $B_{I, i}$
is an intersection of two halfspaces, its Gaussian surface
area is at most a constant. It follows that $\gamma_n^+(\{x: h(x) = \mathbf{e}_i\}) \le C(\epsilon, k)/\sqrt t$ for every $i$.
Hence, $\E[|\nabla h|] \le C(\epsilon, k)/\sqrt t$.

After applying \expref{Corollary}{corr:gradient}, this completes the proof except that $d = d(k, \epsilon, t)$
instead of the claimed $d(k, \epsilon)$, where $d(k, \epsilon, t)$ blows up as $t \to 0$.
To eliminate this dependence on $t$, it suffices to note that the claim is trivial if $t \ll (\epsilon/k)^2$.
Indeed, one can easily construct $h$ with $\E[h] = \E[f]$, $\Stab_t(h) \ge \E[|h|^2] - O(k \sqrt t)$
and $\E[|\nabla h|] \le O(k)$. For example, we could take the pieces $\{x: h(x) = \mathbf{e}_i\}$ to be parallel
slabs of the form $\{x: a_i \le x_1 \le b_i\}$, where $\{a_i\}$ and $\{b_i\}$ are chosen so that the slabs have the
correct volumes. If $t \ll (\epsilon/k)^2$ then such an example will satisfy the claim of the lemma.
\end{proof}

\section{Reduction from PTFs to PTFs on a constant number of variables}\label{sec:junta}

In this section, we are going to prove \expref{Theorem}{thm:reduction}. We restate it here again for the convenience of the reader. \\
\textbf{\expref{Theorem}{thm:reduction}.}%
\emph{
  Let $f: \mathbb{R}^n \rightarrow [k]$ be a degree-$d$, $(d,\epsilon)$-balanced  PTF with $\E_{x}[f(x)] = \boldsymbol{\mu}$ where $\boldsymbol{\mu} = (\mu_1, \ldots, \mu_k)$.  Further, let us assume that 
$\Pr[x \in \mathsf{Collision} (f)] \le \epsilon/(40 k^2)$.  
  Then, for every $\epsilon>0$ and $t > 0$, there exists a degree-$d$ 
PTF $\fh: \mathbb{R}^{n_0} \rightarrow [k]$ such that
\begin{itemize}
    \item $\Vert \mathbf{E}_{x} [\fh(x)] - \boldsymbol{\mu} \Vert_1 \le \epsilon$, and
    \item $ \mathbf{E}[\langle \fh, P_t \fh \rangle] \ge \mathbf{E}[\langle f, P_t f \rangle] - \epsilon$.
\end{itemize}
}

Our main workhorse for this section is going to be two structural theorems for
low-degree polynomials proven in~\cite{DS14}. In order to state these theorems,
we will need to introduce several notions related to tensors and polynomial decomposition.

\subsection{Polynomials and tensors}
We give a brief description of the connection between symmetric tensors and polynomials under the Gaussian distribution. The interested reader may consult the book~\cite{Janson:97} for a detailed background. 
Let $\mathcal{H}$ denote the Hilbert space $\mathbb{R}^n$ and let
$\mathcal{H}^{\otimes q}$ be used to denote the $q$-ary tensor product of
$\mathcal{H}$, \ie, the space of all multilinear functions $\mathcal{H}^q \to \R$.
For $f_1, \dots, f_q \in \mathcal{H}$, their tensor product $f_1 \otimes \cdots \otimes f_q$
is the element of $\mathcal{H}^{\otimes q}$ defined by
\[
    (f_1 \otimes \cdots \otimes f_q)(h_1, \dots, h_q) = \prod_{i=1}^q \inr{f_i}{h_i}.
\]
It is straightforward to check that $\mathcal{H}^{\otimes q}$ is spanned
by elements of the form $f_1 \otimes \cdots \otimes f_q$; we therefore define
an inner product on $\mathcal{H}^{\otimes q}$ by setting
\[
    \inr{f_1 \otimes \cdots \otimes f_q}{g_1 \otimes \cdots \otimes g_q} = \prod_{i=1}^q \inr{f_i}{g_i}
\]
and extending to $\mathcal{H}^{\otimes q}$ by bilinearity.
The norm of a tensor is defined by $\|f\|_F^2 = \inr{f}{f}$. 

An element $f  \in \mathcal{H}^{\otimes q}$ is said to be
symmetric if it is invariant under any permutation $\sigma: [q] \rightarrow
[q]$, in the sense that $f(h_1, \dots, h_q) = f(h_{\sigma(1)}, \dots
h_{\sigma(q)})$ for every $h_1, \dots, h_q \in \mathcal{H}$. Let
$\mathcal{H}^{\odot q}$ denote the symmetric subspace of
$\mathcal{H}^{\otimes q}$. 
%
%
%
%
We now describe a map
between the space $\mathcal{H}^{\odot q}$ and polynomials. Recall that $H_q$
denotes
a Hermite polynomial.
\begin{definition}
    The iterated It\^{o} integral $I_q$ maps $\mathcal{H}^{\odot q}$ as follows: Let $h \in \mathcal{H}$ be a unit vector and note that  $h^{\otimes q} \in \mathcal{H}^{\odot q}$. Then, 
$I_q(h^{\otimes q}) = H_q( \langle h , x \rangle)$ where $x \in \mathbb{R}^n$. The map $I_q$ is extended linearly to $\mathcal{H}^{\odot q}$. 
\end{definition}
For the convenience of the reader, here we describe a basis of the space $\mathcal{H}^{\odot q}$. Consider an unordered multiset $S = \{s_1, \ldots, s_q\} \subseteq [n]$ of size $q$, and define $\Phi_S \in \mathcal{H}^{\odot q}$ by
$$
\Phi_S = \sum_{\sigma \in \mathsf{Sym}_q} e_{s_{\sigma(1)}} \otimes \cdots \otimes e_{s_{\sigma(q)}}.
$$
It is easy to check (and we omit the proof) that the tensors $\Phi_S$ form a basis for $\mathcal{H}^{\odot q}$:
\begin{proposition}\label{prop:basis}
Let $\mathcal{S}$ be the set of unordered multisets of $[n]$ of size $q$. Then, $\{\Phi_S \}_{S \in \mathcal{S}}$ forms an orthogonal basis of $\mathcal{H}^{\odot q}$.
%
%
\end{proposition}
A fundamental property of the map $I_q$ is that it is an isometry between the space of symmetric tensors and polynomials endowed with the standard normal measure (see~\cite{Janson:97} for a proof). 
\begin{proposition}\label{prop:4}
$\mathop{\mathbf{E}}_{x \sim \gamma_n} [I_q(f) \cdot I_p(g)]  = \delta_{p=q} \cdot \langle f,  g \rangle$.
\end{proposition} 
We will  refer to the range of $I_q$ as the Wiener chaos $\mathcal{W}^{q}$. Based on \expref{Proposition}{prop:4}, $I_q$ is a bijective map from $H^{\odot q}$ to $\mathcal{W}^q$. 
On the other hand, it is easy to show that any polynomial $p: \mathbb{R}^n \rightarrow \mathbb{R}$ of degree at most $d$ can be expressed uniquely as 
\begin{equation}\label{eq:tensor-decomp}
p = \sum_{q=0}^d I_q(f_q) \ \ \textrm{where}  \ \ f_q \in \mathcal{H}^{\odot q}. 
\end{equation}

In order to apply the results from~\cite{DS14}, we will often restrict our
attention to multilinear polynomials, \ie, those for which every variable
appears with degree at most $1$. We call a tensor $f \in \mathcal{H}^{\odot q}$
\emph{square-free} if $I_q(f)$ is a multilinear polynomial. Square-free
tensors in $\mathcal{H}^{\odot q}$ may also be characterized as the set of tensors $f$
for which $f(h_1, \dots, h_q) = 0$ whenever $h_1, \dots, h_q$ contains a repeated element.
Alternatively, the space of square-free tensors in $\mathcal{H}^{\odot q}$ is the span
of $\{\Phi_S\}$ where $S$ ranges over all \emph{sets} (not multisets).
It is easy to check that if $p: \mathbb{R}^n \rightarrow \mathbb{R}$ is a multilinear polynomial, 
then the tensors $\{f_q\}_{0 \le q\le d}$ appearing in the decomposition of $p$ in (\ref{eq:tensor-decomp}) are all square-free.

\subsection{It\^{o}'s multiplication formula}

We now state the formula for product of
two polynomials $I_p(f)$ and $I_q(g)$ in terms of $f$ and $g$. For a reference,
see Nourdin's  survey~\cite{Nourdin2013}. First, we define tensor contraction:
for $f \in \mathcal{H}^{\otimes p}$, $g \in \mathcal{H}^{\otimes q}$, and
$r \le p \wedge q$, define $f \ {\otimes}_r \ g \in \mathcal{H}^{\otimes p+q-2r}$ by
\[
    (f\ {\otimes}_r\ g) (t_1, \ldots, t_{p+q-2r})= \sum_{1 \le z_1, \ldots, z_r \le n} f(t_1, \ldots, t_{p-r}, z_{1}, \ldots, z_r) \cdot g(t_{p-r+1}, \ldots, t_{p+q-r}, z_{1}, \ldots, z_r).
\]
The symmetrized contraction product, denoted by $f \ \widetilde{{\otimes}_r} \
g \in \mathcal{H}^{p+q-2r}$ is defined as the symmetrization of the tensor $f \
{\otimes}_r \ g$:
\[
    (f \ \widetilde{\otimes_r} \ g)(t_1, \dots, t_{p+q-2r}) = \frac{1}{(p+q-2r)!} \sum_\sigma (f \ \otimes_r\ g)(t_{\sigma(1)}, \dots, t_{\sigma(p+q-2r)}),
\]
where the sum ranges over all permutations on $[p+q-2r]$.

\begin{proposition}[It\^{o}'s multiplication
    formula]\label{prop:Ito}
\[
I_p(f) \cdot I_q(g) = \sum_{r=0}^{p \wedge q} r! \cdot \binom{p}{r} \cdot \binom{q}{r}  \cdot \frac{\sqrt{(p+q-2r)!}}{\sqrt{p! \cdot q!}}\cdot I_{p+q-2r} (f \ \widetilde{{\otimes}_r} \ g). 
\]
\end{proposition}

We now list a basic fact about contraction products of tensors which shall be
helpful later. In interpreting this fact, it is useful to think of tensors as
tables which can be ``flattened'' into matrices. That is, if we fix an
orthonormal basis $e_1, \dots, e_n$ of $\mathcal{H}$ then the $n^q$ elements
$\{e_{i_1} \otimes \cdots \otimes e_{i_q}: 1 \le 1_1, \dots, i_q \le n\}$ form
a basis for $\mathcal{H}^{\otimes q}$. For any $1 \le r < q$, we can view $f
\in \mathcal{H}^{\otimes q}$ as a $n^r \times n^{q-r}$ matrix, where the rows
are indexed by $[n]^r$, the columns are indexed by $[n]^{q-r}$, and the entry
in position $(i_1, \dots, i_r)$, $(j_1, \dots, j_{q-r})$ is the coefficient in
$f$ of the basis element $e_{i_1} \otimes \cdots \otimes e_{i_r} \otimes
e_{j_1} \otimes \cdots \otimes e_{j_{q-r}}$.  The advantage of this point of
view is that contraction can be viewed as matrix multiplication: if $M_f$ is
the $n^{p - r} \times n^r$ matrix representing $f \in \mathcal{H}^{\otimes p}$ and
$M_g$ is the $n^r \times n^{q-r}$ matrix representing $g \in \mathcal{H}^{\otimes
q}$ then $M_f M_g$ is the $n^{p - r} \times n^{q - r}$ matrix representing $f
\otimes_r g$.  (Note that each of these matrix representations depends on the
choice of a basis for $\mathcal{H}$, but that the contraction itself does not.)

The following fact follows immediately from the previous paragraph and the
Cauchy-Schwarz inequality:
\begin{fact}~\label{fact:contraction}
For $f \in \mathcal{H}^{\otimes p}$ and $g \in \mathcal{H}^{\otimes q}$ and $0 \le r \le p \wedge q$, $\Vert f \otimes_r g \Vert_F \le \Vert f \Vert_F \cdot \Vert g \Vert_F$. Consequently, $\Vert f \widetilde{\otimes}_r g \Vert_F \le \Vert f \Vert_F \cdot \Vert g \Vert_F$. 
\end{fact}

%
%
%
%
%
%
%

The following polynomial anticoncentration theorem is due to Carbery and Wright. 
\begin{theorem}\emph{\cite{CW:01}}\label{thm:CW}
Let $p: \mathbb{R}^n \rightarrow \mathbb{R}$ be a degree-$d$ polynomial.  Then, for $\epsilon>0$, 
$$
\sup_{\theta \in \mathbb{R}} \Pr_{x} [|p(x) -\theta| \le \epsilon \cdot \sqrt{\mathsf{Var}[p]}] = O(d \cdot \epsilon^{1/d}). 
$$
\end{theorem}
An immediate consequence is that $\Pr_{x} [|p(x)| > d^{-O(d)}] \ge {1}/{2}$ for a degree-$d$
polynomial $p$ with $\mathsf{Var}[p] = 1$. With this observation and \expref{Theorem}{thm:hyper},
it is easy to deduce the following bound. (See Lemma~5 in~\cite{DS14}.)
\begin{theorem}\label{thm:combine-hyper}
Let $a, b: \mathbb{R}^n \rightarrow \mathbb{R}$ be degree-$d$ polynomials. Further, $\mathbf{E}_x [a(x) - b(x)]=0$ and $\mathsf{Var}[a-b] \le (\tau/d)^{3d} \cdot \mathsf{Var}[a]$. Then, $\Pr_{x}[\mathsf{sign}(a(x)) \neq  \mathsf{sign}(b(x))] = O(\tau)$. 
\end{theorem}

Before we proceed further, we will make a minor simplifying assumption, namely that the polynomials involved in defining $f$ in \expref{Theorem}{thm:reduction} can be assumed to be multilinear. As we have said earlier, this is because the results in~\cite{DS14} are stated for multilinear polynomials (though it should be easily possible to carry it over to non-multilinear polynomials). In order to show this, we will use the following simple lemma. 
\begin{lemma}\label{lem:multilinear} 
Let $f: \mathbb{R}^n \rightarrow [k]$ be a degree-$d$, $(d,\epsilon)$-balanced PTF\@. Then, for any $\epsilon>0$ and $t > 0$, there exists a degree-$d$, $(d,2\epsilon)$-balanced multilinear PTF $f_{\mathsf{multi}}: \mathbb{R}^\ell \rightarrow \mathbb{R}$ such that 
\begin{itemize}
\item $\Vert \mathbf{E}_{x } [f(x) ] - \mathbf{E}_{x } [f_{\mathsf{multi}}(x)  ] \Vert_1\le \epsilon$, 
\item $| \mathbf{E}_{x } [ \langle f, P_t  f \rangle ] - \mathbf{E}_{x } [\langle f_{\mathsf{multi}}, P_t f_{\mathsf{multi}}\rangle ] | \le \epsilon$.
\item $\Pr_{x } [x \in \mathsf{Collision}(f_{\mathsf{multi}})] \le \Pr_{x } [x \in \mathsf{Collision}(f)]  + \epsilon$.
\end{itemize}
Here $\ell = n \cdot (k^2/d \epsilon)^{3d} \cdot d^2$. 
\end{lemma}

\begin{proof}
    Let $f = \mathsf{PTF}(p^{(1)}, \ldots, p^{(k)})$, where $p^{(1)}, \dots, p^{(k)}$ are $(d, \epsilon)$-balanced.
    Let $p^{(i)} = \sum_{q=0}^d I_q (f_q^{(i)})$ for $f_q^{(i)} \in \mathcal{H}^{\otimes q}$.
    For some $T \in \mathbb{N}$ which will be fixed later, consider a
    collection of $\ell = n \cdot T$ variables denoted by  $\{x_{i,j}\}_{1 \le
    i \le n, 1 \le j \le T}$.
    For $s = 1, \dots, k$, we define the polynomial $r^{(s)}$ (in the variables $\{x_{i,j}\}$) by
    \[
        r^{(s)}(\{x_{i,j}\}) = p^{(s)} \left(
            \frac{x_{1,1} + \cdots + x_{1,T}}{\sqrt T}, \dots, \frac{x_{n,1} + \cdots + x_{n,T}}{\sqrt T}
        \right).
    \]
    In terms of the tensor representations, if
    $f_q^{(i)} = \sum_{1 \le j_1, \ldots, j_i \le n} \alpha^{(q)}_{j_1, \ldots, j_i} e_{j_1} \otimes \cdots \otimes e_{j_i}$
    and we define $g^{(i)}_q$ by $r^{(i)} = \sum_{q=0}^d I_q(g^{(i)}_q)$ then
$$
g_{q}^{(i)}= \sum_{1 \le j_1, \ldots, j_q \le n} \alpha^{(i)}_{j_1, \ldots, j_q} \frac{e_{j_1,1} + \cdots + e_{j_1,T}}{\sqrt{T}} \otimes \cdots \otimes \frac{e_{j_q,1} + \cdots + e_{j_q,T}}{\sqrt{T}}. 
$$
Next, define $f' : \mathbb{R}^{\ell} \rightarrow [k]$ as $f'= \mathsf{PTF}(r^{(1)}, \ldots, r^{(k)})$. Since the joint distribution
of $r^{(1)}, \dots, r^{(k)}$ under $\gamma_\ell$ equals the joint distribution of $p^{(1)}, \dots, p^{(k)}$ under $\gamma_n$,
we have that for $1\le i \le k$,
\begin{eqnarray}
\Pr_{x \sim \gamma_n}[f(x)=i] &=&\Pr_{x \sim \gamma_\ell}[f'(x)=i], \nonumber \\ 
\mathbf{E}_{x \sim \gamma_n} [\langle f, P_t f \rangle]&=&\mathbf{E}_{x \sim \gamma_\ell} [\langle f', P_t f' \rangle], \nonumber \\
\textrm{and} \ \Pr_{x \sim \gamma_n} [x \in \mathsf{Collision}(f)] &=& \Pr_{x \sim \gamma_\ell} [x \in \mathsf{Collision}(f')]. \label{eq:equality} 
\end{eqnarray}

To construct a multilinear polynomial, let $h_q^{(i)}$ be the orthogonal
projection of $g_q^{(i)}$ on the space of square-free tensors, and define the
polynomials $w^{(1)}, \ldots, w^{(k)}$ by
$
w^{(i)} = \sum_{q=0}^d I_q(h_{q}^{(i)}). 
$
Finally, we set $f_{\mathsf{multi}} : \mathbb{R}^{\ell} \rightarrow [k]$ as $f_{\mathsf{multi}} = \mathsf{PTF} (w^{(1)}, \ldots, w^{(k)})$
and we claim that it has the desired properties. By its construction, $f_{\mathsf{multi}}$ is a degree-$d$ multilinear PTF\@.
Next,  for $1 \le j_1, \ldots, j_q \le n$, let us define $\mathcal{S}_{j_1, \ldots, j_q}$ as 
$$
\mathcal{S}_{j_1, \ldots, j_q} = \{(s_1, \ldots, s_q) \in [T]^q : | \cup_{b=1}^q (j_b, s_b) | <  q\}. 
$$
With this definition, observe that
\[
r^{(i)} - w^{(i)}
=  \sum_{q=1}^d \quad \sum_{1 \le j_1, \ldots, j_q \le n}\quad \sum_{(s_1, \ldots, s_q) \in \mathcal{S}_{j_1, \ldots, j_q} } \alpha^{(i)}_{j_1, \ldots, j_q}   \cdot \frac{1}{T^{q/2}} \cdot I_q \big(\otimes_{u=1}^q e_{j_u, s_u} \big).
\]
As the tensors $\{\otimes_{u=1}^T e_{j_u, s_u} \}_{j_u \in [n], s_u \in [T]}$ are orthonormal, using \expref{Proposition}{prop:4}, we get that
$$
\mathsf{Var}[r^{(i)} -  w^{(i)} ] = \sum_{q=1}^d\ \sum_{1 \le j_1, \ldots, j_q \le n}  (\alpha^{(i)}_{j_1, \ldots, j_q})^2 \cdot \frac{|\mathcal{S}_{j_1, \ldots, j_q}|}{T^q}. 
$$
On the other hand, 
$$
\mathsf{Var}[r^{(i)}] = \sum_{q=1}^d\ \sum_{1 \le j_1, \ldots, j_q \le n}  (\alpha^{(i)}_{j_1, \ldots, j_q})^2.
$$
It is easy to see that $|\mathcal{S}_{j_1, \ldots, j_q}| \le |\mathcal{S}_{1,\ldots,1}| \le q^2 T^{q-1} \le d^2 T^{q-1}$. Thus,
$$
\mathsf{Var}[r^{(i)} - w^{(i)} ] \leq \mathsf{Var}[r^{(i)} ] \cdot \frac{d^2}{T}. 
$$
By applying \expref{Theorem}{thm:combine-hyper}, we get that if $T \ge (\tau/d)^{3d} \cdot d^2$, then 
$$
\Pr_{x \sim \gamma_\ell} [\mathsf{sign} (r^{(i)})  \neq  \mathsf{sign} (w^{(i)})] \le \tau. 
$$
By a union bound, it follows that
$
\Pr_{x \sim \gamma_\ell} [f_{\mathsf{multi}}(x) \neq  f'(x)]   \le k \cdot \tau
$
and
\[
\Pr_{x \sim \gamma_\ell} [x \in \mathsf{Collision} (f_{\mathsf{multi}})] \leq \Pr_{x \sim \gamma_\ell} [x \in \mathsf{Collision} (f')] + k \cdot \tau. 
\]
Combining this with (\ref{eq:equality}) and setting $\tau = \epsilon/k$ gives us the claim. Also, by an application of Cauchy-Schwarz inequality, we have that 
\[
\sqrt{\mathsf{Var}(w^{(i)})} \ge \sqrt{\mathsf{Var}(r^{(i)})}  \cdot \bigg( 1- \frac{d}{\sqrt{T}}\bigg) \ge \sqrt{\mathsf{Var}(r^{(i)})}  \cdot \big( 1- \epsilon\big).
\]
Hence, by rescaling the polynomials $w^{(i)}$,  we can ensure that $f_{\mathsf{multi}}$ is a $(d,2\epsilon)$-balanced PTF.
\end{proof}

\subsection{Eigenvalues of tensors and polynomials}
We will now define the notion of eigenvalues of tensors and the corresponding polynomials. 
\begin{definition}
Let $f \in \mathcal{H}^{\odot d}$ and $d \ge 2$. For a partition of $[d]$ into $S, \overline{S}$, define
$$
\lambda_{S, \overline{S}}(f) = \max_{g \in \mathcal{H}^{\otimes |S|}, h \in \mathcal{H}^{\otimes |\overline{S}|}} \frac{\langle f, g \otimes h \rangle}{\Vert g \Vert_F \cdot \Vert h \Vert_F}.
$$
Define $\lambda_{\max}(f)$ by
$$
\lambda_{\max}(f) = \max_{S: 0< |S|< d} \lambda_{S, \overline{S}} (f). 
$$
\end{definition}

\begin{definition}\label{def:eigen-tensor}
    Let $p = \sum_{q=0}^d I_q(f_q)$ be a multilinear polynomial of
degree $d>0$.  and let $(f_0, \ldots, f_d) \in \mathcal{H}^{\odot 0} \times
\cdots \times \mathcal{H}^{\odot d}$ be the tensor associated with $p$. Define
$$
\lambda_{\max}(p)  = \max_{1 < q \le d} \lambda_{\max}(f_q). 
$$
(Note that we exclude $\lambda_{\max}(f_1)$.)
Further, we say that $p$ is $\delta$-eigenregular if $\frac{\lambda_{\max}(p)}{\sqrt{\Var(p)}} \le \delta$. 
\end{definition}

It is easy to check that eigenvalues of tensors are invariant under unitary transformations, from which it follows that
eigenregularity is also unaffected by certain changes of variables:
\begin{fact}
If $p(x)$ is $\delta$-eigenregular then for any $0 < \rho < 1$,  $p(\rho \cdot x + \sqrt{1-\rho^2} y)$ is also $\delta$-eigenregular. 
\end{fact}
The next fact states that the contraction product with an eigenregular tensor is significantly contractive. 
\begin{fact}~\label{fact:contraction1}
Let $f \in \mathcal{H}^{\odot d}$  satisfy $\lambda_{\max} (f) \le \kappa$. For any $g \in \mathcal{H}^{\otimes d'}$ and any $0< r \le \min \{ d-1,d' \}$, 
$\Vert f \otimes_r g \Vert_F \le \kappa$. 
\end{fact}
\begin{proof}
    For a fixed orthonormal basis of $\mathcal{H}$, we represent $f$ and $g$
    with the matrices $M_f$ and $M_g$ as discussed above, where $M_f$ is an
    $n^{d-r} \times n^{r}$ matrix and $M_g$ is $n^r \times n^{d'-r}$.
    The condition $\lambda_{\max}(f) \le \kappa$ implies that $\|M_f\|_{op} \le \kappa$,
    where $\|\cdot\|_{op}$ denotes the operator norm.
For
$\mathcal{S} \in [n]^{d'-r}$, let $M_{g, \mathcal{S}}$ denote the
corresponding column of $M_{g}$. Then the corresponding column of
$M_f \cdot M_g$ is given by $M_f M_{g, \mathcal{S}}$. Thus, 
\[
\Vert M_f \cdot M_g \Vert_F^2  = \sum_{\mathcal{S} \in [n]^{d'-r}} \Vert M_f \cdot  M_{g, \mathcal{S}} \Vert_F^2 \le \sum_{\mathcal{S} \in [n]^{d'-r}} \kappa^2 \cdot \Vert M_{g, \mathcal{S}} \Vert_F^2 = \kappa^2 \cdot \Vert M_g \Vert_F^2.
\]
This finishes the proof. 
\end{proof}

We next have the following easy proposition which says the sum of copies of the same polynomial over disjoint variables is eigenregular. 
\begin{fact}\label{fact:eigen-2}
Let $p: \mathbb{R}^{n} \rightarrow \mathbb{R}$ and define $\tilde p : \mathbb{R}^{n k} \rightarrow \mathbb{R}$ by
$$
\tilde p(X_1, \ldots, X_k) = \frac{1}{\sqrt{k}} \cdot \big( p(X_1) + \cdots + p(X_k)), 
$$
where each $X_i = (x_{i,1}, \dots, x_{i,n})$ is a disjoint ``block'' of $n$ variables. Then $\tilde p$ is $\frac{1}{\sqrt{k}}$-eigenregular. 
\end{fact}
\begin{proof}
Let $d$ be the degree of $\tilde p$ and define the tensors $f_q$ by $q(X_1, \ldots, X_k) =  \sum_{q=0}^{d}I_q(f_q)$.
Let $\mathcal{H} = \mathbb{R}^{n k}$ and let $\mathcal{H} = \mathcal{H}_1
\oplus \cdots \oplus \mathcal{H}_k$ where $\mathcal{H}_i = \mathbb{R}^n$
corresponds to the coordinates in $X_i$. Note that $f_q$ has a block diagonal
structure on $\mathcal{H}_1 \oplus \cdots \oplus \mathcal{H}_k$, in the sense that
$f_q(h_1, \dots, h_q) = 0$ whenever two of the $h_i$ belong to different components $\mathcal{H}_i$.

For any $g = \mathcal{H}^{\otimes q}$, let $g^{(i)} \in \mathcal{H}_i^{\otimes q}$ denote the $i$th ``diagonal block''
of $g$, defined by
\[
    g^{(i)}(h_1, \dots, h_q) = g(\iota(h_1), \dots, \iota(h_q)),
\]
where $\iota: \mathcal{H}_i \to \mathcal{H}$ is the inclusion map.
The definition of $\tilde p$ ensures that $\|f_q^{(1)}\|_F^2 = \cdots = \|f_q^{(q)}\|_F^2$
for every $q$, and the fact that $f_q$ is block diagonal ensures that
$\|f_q\|_F^2 = \sum_{i=1}^k \|f_q^{(i)}\|_F^2 = k \|f_q^{(1)}\|_F^2$.
Assuming without loss of generality that $\Var(\tilde p) = \Var(p) = 1$,
we have $\|f_q^{(i)}\|_F \le k^{-1/2} \|f_q\|_F^2 \le k^{-1/2}$ for every $i$.


For any non-trivial partition $S, \overline{S}$ of $[q]$, we have
\begin{eqnarray*}
\lambda_{S, \overline{S}}(f_q) &=&  \max_{g \in \mathcal{H}^{\odot S}, h \in \mathcal{H}^{\odot \overline{S}}} \frac{\langle f_q , g \otimes h \rangle}{\Vert g \Vert_F \cdot \Vert h \Vert_F} \\
&=& \max_{g \in \mathcal{H}^{\odot S}, h \in \mathcal{H}^{\odot \overline{S}}} \frac{\sum_{i=1}^q \langle f_q^{(i)}, g^{(i)} \otimes h^{(i)} \rangle }{\Vert g \Vert_F \cdot \Vert h \Vert_F} \\
&\le& \max_{g \in \mathcal{H}^{\odot S}, h \in \mathcal{H}^{\odot \overline{S}}} \frac{\sum_{i=1}^q \langle f_q^{(i)}, g^{(i)} \otimes h^{(i)} \rangle }{ \sqrt{\sum_{i=1}^q \Vert g^{(i)} \Vert_F^2} \cdot  \sqrt{\sum_{i=1}^q \Vert h^{(i)} \Vert_F^2} } \\
 &\le& \max_{g \in \mathcal{H}^{\odot S}, h \in \mathcal{H}^{\odot \overline{S}}} \frac{\sum_{i=1}^q \frac{1}{\sqrt{k}} \Vert g^{(i)} \otimes h^{(i)} \Vert_F }{ \sqrt{\sum_{i=1}^q \Vert g^{(i)} \Vert_F^2} \cdot  \sqrt{\sum_{i=1}^q \Vert h^{(i)} \Vert_F^2} }\\
 &=& \max_{g \in \mathcal{H}^{\odot S}, h \in \mathcal{H}^{\odot \overline{S}}} \frac{\sum_{i=1}^q \frac{1}{\sqrt{k}} \Vert g^{(i)} \Vert_F \cdot \Vert  h^{(i)} \Vert_F }{ \sqrt{\sum_{i=1}^q \Vert g^{(i)} \Vert_F^2} \cdot  \sqrt{\sum_{i=1}^q \Vert h^{(i)} \Vert_F^2} } \le \frac{1}{\sqrt{k}},
\end{eqnarray*}
where the last equality uses the fact that $\Vert A \otimes B \Vert_F = \Vert A
\Vert_F \cdot \Vert B \Vert_F$, and the last inequality is a consequence of
Cauchy-Schwarz. This implies that $\lambda_{\max}(f_q) \le {1}/{\sqrt k}$ for all $q$,
and hence $\lambda_{\max}(\tilde p) \le {1}/{\sqrt k}$.
\end{proof}

One of the two main results that we will use from~\cite{DS14} is a Central Limit Theorem for Gaussian polynomials:

\begin{theorem}\label{thm:CLT}
Let $d >1$ and let $p_1, \ldots, p_t: \mathbb{R}^n \rightarrow \mathbb{R}$ be $t$ degree-$d$ polynomials such that for $1\le i \le t$, $\mathbf{E}[p_i]=0$ and $\mathsf{Var}(p_i) \le 1$ and each $p_i$ is $\epsilon$-eigenregular. Let $F = (p_1(x), \ldots, p_t(x))$ where $x \sim \gamma_n$ and let $C$ denote the covariance matrix of $F$, \ie, $C[i,j] = \mathbf{E}_{x \sim \gamma_n} [p_i(x) \cdot p_j(x)]$.  Let $Z= (Z_1, \ldots, Z_t)$ be a $t$-dimensional Gaussian with mean zero and covariance $C$. 
Then, for any $\alpha:\mathbb{R}^t \rightarrow \mathbb{R}$ such that $\alpha \in \mathcal{C}^2$, 
$$
\big| \mathbf{E}_{x \sim \gamma_n} \big[ \alpha(p_1(x), \ldots, p_t(x)) \big] - \mathbf{E}_{Z} \big[ \alpha(Z_1, \ldots, Z_t) \big] \big|  \le 2^{O( d \log d)} \cdot t^2 \cdot \sqrt{\epsilon} \cdot \Vert \alpha '' \Vert_{\infty}. 
$$
\end{theorem}

\subsection{Polynomial regularity lemma}

The next result we will need from~\cite{DS14} is a regularity lemma for low-degree polynomials.
Suppose that we begin with a collection of $k(d+1)$ polynomials $\{p_{s,q}: 1 \le s \le k, 0 \le q \le d\}$,
where $p_{s,q} \in {\cal W}^q$ and $\Var[p_{s,q}] = 1$ for all $s$ and $q$. We will consider ``decompositions''
of these polynomials consisting of the following objects:
\begin{itemize}
    \item for every $1 \le s \le k$ and $0 \le q \le d$, an ``outer'' polynomial $\Outer(p_{s,q})$ of arity $\num(s,q) \in \N$; and
    \item for every $q \le s \le k$, $0 \le q \le d$, and $1 \le \ell \le \num(s,q)$, an ``inner'' polynomial $\Inner(p_{s,q})_\ell$.
\end{itemize}
We also introduce the following statistics of these polynomials:
\begin{itemize}
    \item the ``total arity,'' $\mathrm{Num} = \sum_{s=1}^k \sum_{q=0}^d \num(s,q)$;
    \item $\mathrm{Coeff}(p_{s,q})$, defined to be the sum of the absolute values of the coefficients of $\Outer(p_{s,q})$; and
    \item $\mathrm{Coeff} = \sum_{s=1}^k \sum_{q=0}^d \mathrm{Coeff}(p_{s,q})$.
\end{itemize}

\begin{definition}
    Consider polynomials $\{p_{s,q}\}$ and a function $\beta: [1, \infty) \to (0, 1)$ satisfying $\beta(x) \le 1/x$.
    Given collections of polynomials $\{\Outer(p_{s,q})$ and $\{\Inner(p_{s,q})_\ell\}$ as above, we say that
        these polynomials are a $(\beta, M, N)$ decomposition of $\{p_{s,q}\}$ if the following conditions hold:

\begin{enumerate}
    \item for every $1 \le s \le k$, every $0 \le q \le d$, and every $x \in \R^n$,
    \begin{equation} \label{eq:tildep}
        p_{s,q}(x)=
        \Outer(p_{s,q})\left(
        \Inner(p_{s,q})_1(x),\dots, \Inner(p_{s,q})_{\num({s,q})}(x)\right);
    \end{equation}
    \item for every $1 \le s \le k$, every $1 \leq q \leq d$, and every $1 \le \ell \le \num(s,q)$,
        the polynomial $\Inner(p_{s,q})_\ell$ lies in ${\cal W}^j$ for some $1 \leq j \leq q$ and
        satisfies $\Var[\Inner(p_{s,q})_\ell] = 1$;

    \item for every $1 \le s \le k$ and every $1 \leq q \leq d$,
    $\Outer(p_{s,q})$ is a multilinear polynomial of degree at most $d$;

    \item $\mathrm{Coeff} \le M$ and $\mathrm{Num} \le N$; and

    \item for every $1 \le s \le k$, every $1 \leq q \leq d$, and every $1 \le \ell \le \num(s,q)$,
        $\Inner(p_{s,q})_\ell$ is $\beta(\mathrm{Num} + \mathrm{Coeff})$-eigenregular.
\end{enumerate}
\end{definition}

Intuitively, the above definition states the following: given a decreasing
function $\beta(\cdot)$ and a collection of $k$ multilinear polynomials of degree
$d$, a $(\beta, M,N)$ decomposition of expresses the polynomials as a ``outer
polynomial'' composed with ``inner polynomials.'' The quality of this
decomposition is captured by the following parameters: 
\begin{itemize}
\item the sum of the arities (denoted by $\mathrm{Num}$, and bounded by $N$) and the sum of
    absolute values of coefficients (denoted by $\mathrm{Coeff}$, and bounded by $M$) of the outer
    polynomials; and
\item the eigenregularity of the inner polynomial (bounded in terms of $\beta$).
\end{itemize}

The main decomposition theorem of~\cite{DS14} says that approximate $(\beta, M, N)$
decompositions exist and can be computed.

\begin{theorem}~\label{thm:decomp} 
Fix $d \geq 2$ and fix any non-increasing computable
function $\beta:  [1,\infty) \to (0,1)$ that satisfies $\beta(x) \leq 1/x$.
There is an algorithm {\bf MultiRegularize-Many-Wieners}$_{d,\beta}$
with the following properties.
The algorithm takes as input:
\begin{itemize}
    \item a collection $\{p_{s,q}: 1 \le s \le k, 0 \le q \le d\}$ of multilinear polynomials
        satisfying $p_{s,q} \in {\cal W}^q$ and $\Var[p_{s,q}] = 1$ for every $s, q$; and
    \item a parameter $\tau > 0$.
\end{itemize}
The output of the algorithm consists of
\begin{itemize}
    \item a collection $\{\tilde p_{s,q}: 1 \le s \le k, 0 \le q \le d\}$
        of polynomials
        satisfying (for every $s$ and $q$) $\tilde p_{s,q} \in {\cal W}^q$,
    $\tilde p_{s,0} = p_{s,0}$, and $\Var[p_{s,q} - \tilde p_{s,q}] \le \tau$;
    \item two numbers $M = M_\beta(k, d, \tau)$ and $N = N_\beta(k, d, \tau)$; and
    \item a $(\beta, M, N)$ decomposition of $\{\tilde p_{s,q}\}$.
\end{itemize}
The functions $M_{\beta}(k,d,\tau)$ and $N_{\beta}(k,d,\tau)$ are computable but not necessarily primitive recursive.
\end{theorem}

We remark that our use of \expref{Theorem}{thm:decomp} from~\cite{DS14} is precisely the reason why the quantity $n_0$ in
\expref{Theorem}{thm:main-stab} is not primitive recursive. 

\subsection{Second moments and approximate equivalence}

As we mentioned in the outline of the proof of \expref{Theorem}{thm:reduction}, one
of the main steps in the proof is (after applying a $(\beta, M, N)$
decomposition) to swap out one collection of inner polynomials for a different
one. The main result of this section essentially says that in order to do so,
it suffices to preserve the second moments of all the inner polynomials:

\begin{lemma}~\label{lem:variance-lowerbound}
Let $\Psi:\mathbb{R}^{\mathsf{Num}_\Psi} \rightarrow \mathbb{R}$ be a
degree-$d$ polynomial, and let $\mathsf{Coeff}_\Psi$ be the sum of the absolute
values of its coefficients.
Let $\{A_i\}_{1 \le i \le
\mathsf{Num}_\Psi}$ and $\{B_i \}_{1 \le i \le \mathsf{Num}_\Psi}$  be centered
families of polynomials of variance $1$ such that for all $1 \le i \le m$,
$\exists 0< j \le d$ such that $A_i, B_i \in \mathcal{W}^j$. Assume, moreover, that for all
$1 \le i \le j \le \mathsf{Num}_\Psi$, $\mathbf{E}[A_i A_j] = \mathbf{E}[B_i
B_j]$. If each $A_i, B_i$ is $\zeta$-eigenregular where 
$
\zeta \cdot \mathsf{Coeff}_\Psi \cdot 2^{\mathsf{Num}_\Psi \cdot (\mathsf{Num}_\Psi \cdot d+1)} \leq \kappa, 
$
then $|\mathbf{E}[\Psi(A_1, \ldots, A_{\mathsf{Num}})] - \mathbf{E}[\Psi(B_1, \ldots, B_{\mathsf{Num}})]| \le \kappa$. 
\end{lemma}

Our main tool in the proof of \expref{Lemma}{lem:variance-lowerbound} is the following lemma,
which says that the expectation of a product of eigenregular polynomials is essentially
determined by their second moments.

\begin{lemma}\label{lem:monomial}
    Let $p_1, \ldots, p_t \in \mathbb{R}^n \rightarrow \mathbb{R}$ be $\kappa$-eigenregular polynomials with variance 1, such that $p_i \in \mathcal{W}^{q_i}$ for every $i$.
    Define $C \in \mathbb{R}^{t \times t}$ by $C(i,j) = \langle p_i, p_j \rangle$. 
There exists a function $F : \mathbb{R}^{t \times t} \times \mathbb{Z}^{t} \rightarrow \mathbb{R}$ such that 
\[
\bigg| \mathbf{E}\big[\prod_{i=1}^t p_i \big] - F(C, q_1, \ldots, q_t) \bigg| \le  2^{t (q_1 + \cdots + q_t+1)} \cdot \kappa. 
\]
In other words, up to the error term of $2^{t (q_1 + \cdots + q_t)} \cdot \kappa$, the expectation of the product $\prod_{i=1}^t p_i$ is just dependent on the degrees of the polynomials and the covariance matrix of the polynomials. 
\end{lemma}

\begin{proof}[Proof of \expref{Lemma}{lem:variance-lowerbound}]
Applying \expref{Lemma}{lem:monomial} and triangle inequality, we have
$$
|\mathbf{E}[\Psi(A_1, \ldots, A_{\mathsf{Num}})] - \mathbf{E}[\Psi(B_1, \ldots, B_{\mathsf{Num}})]|  \le \zeta \cdot\mathsf{Coeff}_\Psi \cdot 2^{\mathsf{Num}_\Psi \cdot (\mathsf{Num}_\Psi \cdot d+1)}. 
$$
Applying the assumption on $\zeta$ completes the proof.
\end{proof}

Before proving \expref{Lemma}{lem:monomial}, we will explore some further consequences.
The following lemma is along the same lines as \expref{Lemma}{lem:variance-lowerbound},
but it concerns the sign of $\Psi$ instead of just its expectation.

\begin{lemma}~\label{lem:sign-match}
    Fix $\epsilon > 0$ and $\eta > 0$.
    Let $\Psi$, $A_i$, and $B_i$ satisfy the assumptions of the previous lemma, but
    assume that $A_i$ and $B_i$ are $\zeta$-eigenregular for $\zeta$ satisfying
    \begin{eqnarray*}
    \zeta \cdot \mathsf{Coeff}^2_\Psi \cdot 2^{\mathsf{Num}_\Psi \cdot (\mathsf{Num}_\Psi \cdot 2d+1)} \cdot 2^{O(d^3)} \leq \frac{\eta}{4} , \ \ 
     2^{O(d \log d)} \cdot \mathsf{Num}_{\Psi}^2 \cdot (4c^2) \cdot \zeta \le \epsilon,
    \end{eqnarray*}
    where $c$ is defined by
    \[
    B = \Omega\bigg( \ln  \frac{\mathsf{Num}_\Psi \cdot d}{\epsilon} \bigg)^{d/2} ;  \ \delta = \bigg( \frac{\epsilon \cdot \eta^{1/2d}}{d \cdot B \cdot \sqrt{\mathsf{Num}_\Psi} \cdot \mathsf{Coeff}_\Psi^{1/d}} \bigg)^d; \   c = \frac{\mathsf{Num}_{\Psi}}{\delta \cdot \sqrt{\epsilon}}.
    \]
    If $\mathsf{Var}(\Psi(A_1, \ldots,
    A_{\mathsf{Num}_\Psi})) \ge \eta$ and $|\mathbf{E}[\Psi(A_1, \ldots,
    A_{\mathsf{Num}_\Psi})]| \le \eta$ then
    \[
        |\mathbf{E}[\mathsf{sign}(\Psi(A_1, \ldots, A_{\mathsf{Num}}))] - \mathbf{E}[\mathsf{sign}(\Psi(B_1, \ldots, B_{\mathsf{Num}}))]| \le O(\epsilon).
    \]
\end{lemma}

One step in the proof of \expref{Lemma}{lem:sign-match} is a kind of mollification for the sign function:

\begin{lemma}\label{lem:mollification}
    Fix $\epsilon > 0$, $\eta > 0$. Let $\Psi$ be a degree-$d$ polynomial of
    arity $\mathsf{Num}_\Psi$, and
    let $B$, $\delta$, and $c$ be defined as in \expref{Lemma}{lem:sign-match}.

    There is a function $\tilde{g}_c: \mathbb{R}^{\mathsf{Num}_\Psi}
    \rightarrow [0,1]$ satisfying $\|\tilde g_c''\|_\infty \le 4 c^2$ such that
    for every collection $\{A_i\}_{1 \le i \le \mathsf{Num}_\Psi}$ of centered,
    unit-variance polynomials in $n$ variables that satisfy
    $\mathsf{Var}(\Psi(A_1, \ldots, A_m)) \ge \eta$, we have

\[
\mathop{\mathbf{E}}_{x \sim \gamma_n} [|\mathsf{sign}(\Psi(A_1, \ldots, A_m)) - \tilde{g}_c(A_1, \ldots, A_m)|] \le O(\epsilon). 
\]
\end{lemma}

\begin{proof}
Let $R$ denote the set $\{x: \mathsf{sign}(\Psi(x)) = 1\}$. By a standard argument based
on convolution with a mollifier, there is a function $\tilde{g}_c :
\R^{\mathsf{Num}_\Psi} \to [0, 1]$ such that $\| \tilde g_c''\|_\infty \le 4
c^2$ and
\begin{equation}\label{eq:diffmoli} 
    \big| \mathsf{sign}(\Psi(x)) - \tilde{g}_c(x) \big| \leq \min \bigg\{ 1,
    \frac{\mathsf{Num}_\Psi^2}{c^2 \cdot \mathsf{dist}^2(x, \partial R)}\bigg\}
\end{equation}
for all $x$, where $\mathsf{dist} (x, \partial R)$  denotes the Euclidean
distance of $x$ from the boundary of $R$.

It is not hard to check (and follows immediately from Claim 49 in~\cite{DS14})
that if $x \in \mathbb{R}^{\Num_\Psi} $ satisfies $\Vert x \Vert_{\infty} \le B$ and
$|\phi(x)| \ge d \cdot (2B)^{d} \cdot \delta \cdot \mathsf{Coeff}_\Psi \cdot \Num_\Psi^{d/2}$
then $\mathsf{dist}(x, \partial R)\ge \delta$, in which case~\eqref{eq:diffmoli} implies
that $|\mathsf{sign}(\Psi(x)) - \tilde g_c(x)| \le \mathsf{Num}_\Psi^2 / (c^2 \delta^2)$.
Rearranging the logic, if $|\mathsf{sign}(\Psi(x)) - \tilde g_c(x)| > \mathsf{Num}_\Psi^2 / (c^2 \delta^2)$ then either $\max_j |A_j(x)| > B$ or $\Psi(A_1(x), \dots A_{\mathsf{Num}_\Psi}(x)) \le d (2B)^d \delta \cdot \Coeff_\Psi \cdot \Num_\Psi^{d/2}$. Hence,
\begin{multline*}
    \mathbf{E}_{x \sim \gamma_n} [|\mathsf{sign}(\phi(A_1, \ldots, A_{\Num_\Psi})) -
    \tilde{g}_c(A_1, \ldots, A_{\Num_\Psi})|] \\
    \le
    \Pr_{x \sim \gamma_n} \big[\max_j |A_j| > B\big] + \Pr_{x \sim \gamma_n} \big[ |\phi(A_1, \ldots, A_{\Num_\Psi})| \le d \cdot (2B)^{d} \cdot \delta \cdot \Coeff_\Psi \cdot \Num_\Psi^{d/2}\big] + \frac{\Num_\Psi^2}{c^2 \delta^2}. 
\end{multline*}
We now bound each term on the right hand side. For the first term, each $A_j$
is a centered, unit-variance, degree-$d$ polynomial and so
\expref{Theorem}{thm:hyper} gives
$$
\Pr_{x \sim \gamma_n} \big[\max_j |A_j(x)| > B\big]  \le \Num_\Psi \cdot d \cdot \exp \big( -B^{2/d}\big).
$$
To bound the second term, apply \expref{Theorem}{thm:CW} and the assumption
that $\mathsf{Var}(\Psi(A_1, \ldots, A_{\Num_\Psi})) \ge \eta$:
\[
\Pr_{x \sim \gamma_n} \big[
    |\phi(A_1, \ldots, A_m)| \le d \cdot (2B)^{d/2} \cdot \delta \cdot \Coeff_\Psi \cdot \Num_\Psi^{d/2}\big] \le \frac{d \cdot (2B) \cdot \delta^{1/d} \cdot \Coeff_\Psi^{1/d} \cdot \sqrt{\Num_\Psi}}{\eta^{1/2d}}.
\]
Plugging in our choice of parameters completes the proof.
\end{proof}


\begin{proof}[Proof of \expref{Lemma}{lem:sign-match}]
    Consider the polynomial $\Upsilon(z_1, \ldots, z_{\mathsf{Num}_\Psi}) =
    (\Psi(z_1, \ldots, z_{\mathsf{Num}_\Psi}))^2$; let $\mathsf{Num}_\Upsilon$
be its arity and let
$\mathsf{Coeff}_\Upsilon$ be the sum of the absolute values of its
coefficients. Observe that $\Upsilon$ is a degree-$2d$ polynomial, and that
$\mathsf{Coeff}_{\Upsilon} \le \mathsf{Coeff}^2_{\Psi}$ and
$\mathsf{Num}_\Upsilon  = \mathsf{Num}_\Psi$. Now, applying
\expref{Lemma}{lem:variance-lowerbound} to the function $\Upsilon$, %
\begin{equation}\label{eq:variance-diff1}
\big|\mathbf{E} [ \Psi^2(A_1, \ldots, A_{\mathsf{Num}})] -  \mathbf{E} [ \Psi^2(B_1, \ldots, B_{\mathsf{Num}})]  \big| \le \eta/4. 
\end{equation}
By \expref{Lemma}{lem:variance-lowerbound} applied to the function $\Psi$,
\begin{equation}\label{eq:variance-diff1a}
  \begin{split}
  \big|\mathbf{E} [ \Psi(A_1, \ldots, A_{\mathsf{Num}})] -  \mathbf{E} [ \Psi(B_1, \ldots, B_{\mathsf{Num}})]  \big| &\le \zeta \cdot \mathsf{Coeff}_\Psi \cdot 2^{\mathsf{Num}_\Psi \cdot (\mathsf{Num}_\Psi \cdot d+1)}\\ &\le \frac{\eta}{4 \cdot \mathsf{Coeff}_\Psi \cdot 2^{\mathsf{Num}_\Psi^2 \cdot d} \cdot 2^{O(d^3)}}.
\end{split}
\end{equation}

H\"older's inequality and the Gaussian hypercontractive inequality (see Theorem 5.10 in~\cite{Janson:97}) imply that
for any monomial in $A_i$,
$\E [\prod_{j=1}^d A_{i_j}] \le \prod_{j=1}^d \E[|A_{i_j}|^d]^{1/d} \le d^{2d} \le 2^{O(d^2)}$;
by the triangle inequality, it follows that $\E[\Psi(A_1, \dots, A_{\mathsf{Num}})] \le \mathsf{Coeff}_\Psi 2^{O(d^3)}$. The same bound applies with $B$ in place of $A$, and together
with~\eqref{eq:variance-diff1a} this yields
\begin{equation}\label{eq:variance-diff2}
\big|\big(\mathbf{E} [ \Psi(A_1, \ldots, A_{\mathsf{Num}})]\big)^2- \big(\mathbf{E} [ \Psi(B_1, \ldots, B_{\mathsf{Num}})]\big)^2 \big| \le \eta/4. 
\end{equation}

Combining (\ref{eq:variance-diff1}) and (\ref{eq:variance-diff2}), we get that 
$
|\mathsf{Var}(\Psi(A_1, \ldots, A_{\mathsf{Num}})) - \mathsf{Var}(\Psi(A_1, \ldots, A_{\mathsf{Num}}))| \le \eta/2
$ and thus, $$\mathsf{Var}(\Psi(B_1, \ldots, B_{\mathsf{Num}})) \ge \eta/2.$$
Applying \expref{Lemma}{lem:mollification}, we obtain a function
$\tilde{g}_c : \mathbb{R}^{\mathsf{Num}} \rightarrow [0,1]$ such that
$\Vert \tilde{g}_c'' \Vert_{\infty} \le 4c^2$,
\begin{align*}
    \big|\mathbf{E} [ \mathsf{sign}(\Psi(A_1, \ldots, A_{\mathsf{Num}}))]  - \mathbf{E} [ \tilde{g}_c(\Psi(A_1, \ldots, A_{\mathsf{Num}}))] \big| &\le O(\epsilon), \text{ and } \\
    \big|\mathbf{E} [ \mathsf{sign}(\Psi(B_1, \ldots, B_{\mathsf{Num}}))]  - \mathbf{E} [ \tilde{g}_c(\Psi(B_1, \ldots, B_{\mathsf{Num}}))] \big| &\le O(\epsilon).
\end{align*}
Finally, applying \expref{Theorem}{thm:CLT}, we have 
\[
\big| \mathbf{E} [ \tilde{g}_c(\Psi(B_1, \ldots, B_{\mathsf{Num}}))]- \mathbf{E} [ \tilde{g}_c(\Psi(A_1, \ldots, A_{\mathsf{Num}}))] \big| \le 2^{O(d \log d)} \cdot \mathsf{Num}_\Psi^2 \cdot (4c^2) \cdot \zeta \le \epsilon.
\]
The last inequality uses the second condition on $\zeta$.  Combining these three inequalities, we obtain the proof. 
\end{proof}


In the remainder of this section, we prove \expref{Lemma}{lem:monomial}. Let $p_i$ and $q_i$
be as in the statement of \expref{Lemma}{lem:monomial}, and let $h_i \in \mathcal{H}^{\otimes q_i}$
be defined by $p_i = I_{q_i}(h_i)$.
To prove the lemma, we will express the product $p_1 \cdot \cdots \cdot p_t$ in terms of the tensors $\{h_1, \ldots, h_t\}$ using successive applications of It\^{o}'s multiplication formula.  Towards this, we first have the following definition. 
\begin{definition} 
Given $p_1, \ldots, p_t$ as above, a tuple $\mathbf{r} = (r_1, \ldots, r_{t-1})$ is said to be a contraction sequence if it satisfies the following definition: 
\begin{equation}\label{eq:sequence}
\textrm{For all } 1\le i \le t-1, \  0 \le r_i \le \min \bigg\{ q_{i+1}, q_i+ \sum_{j=1}^{i-1} (q_j - 2 r_j) \bigg\}.
\end{equation}
\end{definition}
Let $\mathcal{S}_{\mathsf{valid}}$ be the set of all contraction sequences. 
Also for a contraction sequence $\mathbf{r}=(r_1, \ldots, r_t)$,
define the quantity 
\[
\gamma_{\mathbf{r}} = \prod_{i=1}^{t-1} r_i! \cdot \binom{q_{i+1}}{r_i} \cdot \binom{q_i+ \sum_{j=1}^{i-1} (q_j - 2 r_j) }{r_i} \cdot \frac{\sqrt{(q_{i+1}+ \sum_{j=1}^{i} (q_j - 2 r_j) )!}}{\sqrt{q_{i+1}!} \cdot \sqrt{(q_i+ \sum_{j=1}^{i-1} (q_j - 2 r_j))!}}. 
\]
Also, let us define $g^{\mathbf{r}}_1 = h_1$ and for $1 < i  \le t$, define $g^{\mathbf{r}}_i $ iteratively as $g^{\mathbf{r}}_i  = g^{\mathbf{r}}_{i-1} \widetilde{\otimes}_{r_{i-1}} h_i$. Define  $m^{\mathbf{r}}_i = q_{i+1} + \sum_{j=1}^{i} (q_j - 2 r_j)$ and  observe that $g^{\mathbf{r}}_{i+1}  \in \mathcal{H}^{\odot m^{\mathbf{r}}_i}$. 
Applying It\^{o}'s multiplication formula (\expref{Proposition}{prop:Ito}) iteratively, we obtain that
\begin{equation}\label{eq:Ito1}
\prod_{i=1}^t p_i = \sum_{\mathbf{r} =(r_1, \ldots, r_{t-1}) \in \mathcal{S}_{\mathsf{valid}}} \gamma_{\mathbf{r}} I_{m^{\mathbf{r}}_{t-1}}(g^{\mathbf{r}}_{t}).
\end{equation}
The next claim establishes a simple upper bound on $\gamma_{\mathbf{r}}$. 
\begin{claim}~\label{clm:bound-easy1}
$\sum_{\mathbf{r} =(r_1, \ldots, r_{t-1}) \in \mathcal{S}_{\mathsf{valid}}} |\gamma_{\mathbf{r}}| \le 2^{t (q_1 + \cdots + q_t +1 )}. $
\end{claim}
\begin{proof} By applying the fact that $\binom{n}{x} \le 2^n$, it is easy to show that 
\[
r_i! \cdot \binom{q_{i+1}}{r_i} \cdot \binom{q_i+ \sum_{j=1}^{i-1} (q_j - 2 r_j) }{r_i} \cdot \frac{\sqrt{(q_{i+1}+ \sum_{j=1}^{i} (q_j - 2 r_j) )!}}{\sqrt{q_{i+1}!} \cdot \sqrt{(q_i+ \sum_{j=1}^{i-1} (q_j - 2 r_j))!}} \le 2^{\sum_{j=1}^i (q_j - 2 r_j) + q_{i+1} + r_i }. 
\] 
This implies that 
\[
 |\gamma_{\mathbf{r}}|  \le 2^{\sum_{i=1}^{t-1} (\sum_{j=1}^i (q_j - 2 r_j) + q_{i+1} + r_i ) } \le 2^{t(q_1 +\cdots + q_t) - \sum_{i=1}^{t-1} r_{i}}. 
\]
As a result, we obtain 
\[
\sum_{\mathbf{r} =(r_1, \ldots, r_{t-1}) \in \mathcal{S}_{\mathsf{valid}}} |\gamma_{\mathbf{r}}| \le \sum_{\mathbf{r} =(r_1, \ldots, r_{t-1}) \in \mathcal{S}_{\mathsf{valid}}}2^{t(q_1 +\cdots + q_t) - \sum_{i=1}^{t-1} r_{i}} \le 2^{t(q_1 + \cdots + q_t +1)}. \qedhere
\]
\end{proof}
Our strategy for proving \expref{Lemma}{lem:monomial} is as follows. We will first partition the sum in (\ref{eq:Ito1}) into two sets and show that the terms belonging to the first set are all bounded by $\kappa$. We will then further partition the terms in the  second set into two sets: Again, the terms in the first set will be bounded by $\kappa$ whereas terms in the second set will just be a function of $(C, q_1, \ldots, q_t)$ which will conclude our proof. Towards this, let us partition $\mathcal{S}_{\mathsf{valid}}$ into two partitions, $\mathcal{S}_z$ and $\mathcal{S}_{\mathsf{valid}} \setminus \mathcal{S}_z$ where the first set is defined as 
 $\mathcal{S}_z = \{\mathbf{r} \in \mathcal{S}_{\mathsf{valid}}:$ For all $i <t$, $(r_i=0) \vee (r_i = q_{i+1}) \}$.
Next, prove the following proposition. 
\begin{proposition}~\label{prop:approx-error2}
Let $h_1, \ldots, h_t$ be as above such that for all $1 \le i \le t$,  $ \Vert h_1 \Vert_F = \ldots = \Vert h_t \Vert_F=1$. Further, each $h_i$ is $\kappa$-eigenregular. If $\mathbf{r} \not \in \mathcal{S}_z$, then $\Vert g^{\mathbf{r}}_{t} \Vert_F \le \kappa$. 
\end{proposition}
\begin{proof}
If $\mathbf{r} \not \in \mathcal{S}_z$, then there is a coordinate $1 \le j \le t-1$ such that $r_j \neq  0$ and $r_j \neq  q_{j+1}$. Appealing to \expref{Fact}{fact:contraction}, it is easy to see that  $\Vert g^{\mathbf{r}}_{j} \Vert_F \le 1$. Next, we appeal to \expref{Fact}{fact:contraction1}, we get $\Vert g^{\mathbf{r}}_{j} {\otimes}_{r_{j}} h_{j+1} \Vert_F \le \kappa$. By the triangle inequality, symmetrization decreases norm and so we obtain that $\Vert g^{\mathbf{r}}_{j} \widetilde{\otimes}_{r_{j}} h_{j+1} \Vert_F \le \kappa$. Finally, again using that $\Vert h_{j+2} \Vert_F, \ldots, \Vert h_t \Vert_F \le 1$ and applying \expref{Fact}{fact:contraction} iteratively, we obtain $\Vert g^{\mathbf{r}}_{t} \Vert_F \le \kappa$. This finishes the proof.
\end{proof}
To define further partitions of $\mathcal{S}_z$, we will need a somewhat more elaborate definition. To do this, consider any $\mathbf{r} \in \mathcal{S}_z$. Recall that $m^{\mathbf{r}}_i = q_{i+1} + \sum_{j=1}^{i} (q_j - 2 r_j)$ and $g_{i+1} \in \mathcal{H}^{\odot m^{\mathbf{r}}_i}$. For any $\ell \in \mathbb{N}$, let $\mathbb{S}_\ell$ denote the symmetric group on $\ell$ elements. Consider a tuple of permutations $\boldsymbol{\sigma}= (\sigma_1, \ldots, \sigma_{t-1})$ such that $\sigma_i \in \mathbb{S}_{m^{\mathbf{r}}_i}$. For $\boldsymbol{\sigma}$ and $\mathbf{r}$, we define $g^{\mathbf{r}, \boldsymbol{\sigma}}_i$ iteratively as 
$ 
g^{\mathbf{r}, \boldsymbol{\sigma}}_1 = h_1 \ \  \textrm{and} \ \ g^{\mathbf{r}, \boldsymbol{\sigma}}_i = \sigma_{i-1}(g^{\mathbf{r}, \boldsymbol{\sigma}}_{i-1} {\otimes}_{r_{i-1}} h_i).
$
Let $\Lambda_{\mathbf{r}} = \{\boldsymbol{\sigma} : \boldsymbol{\sigma} = ( \sigma_1, \ldots, \sigma_{t-1})\}$. With this notation, we have that 
$
g^{\mathbf{r}}_t = ({1}/{|\Lambda_{\mathbf{r}}|}) \cdot \sum_{\boldsymbol{\sigma} \in \Lambda_{\mathbf{r}}} g^{\mathbf{r}, \boldsymbol{\sigma}}_t. 
$
We will now partition pairs $(\mathbf{r}, \boldsymbol{\sigma})$ (where $\mathbf{r} \in \mathcal{S}_z$). 
For this subsequent partitioning, we will need to associate a graph with the pair $(\mathbf{r}, \boldsymbol{\sigma})$ where $\mathbf{r} \in \mathcal{S}_z$ and then partition based on this graph. To understand what this graph is meant to capture, note that to obtain $g^{\mathbf{r}, \boldsymbol{\sigma}}_{i+1}$ from $g^{\mathbf{r}, \boldsymbol{\sigma}}_{i}$, one of the following two events happen: (a) If $r_i=0$, then we first compute $g^{\mathbf{r}, \boldsymbol{\sigma}}_{i} \otimes h_{i+1}$ and then permute the indices of the resulting tensor by $\sigma_i$. Note that each time an index is added, we can associate a unique $1 \le j \le t$ with this index. (b) If $r_i =q_{i+1}$, we compute the contraction product of $g^{\mathbf{r}, \boldsymbol{\sigma}}_{i}$ with $h_{i+1}$ and subsequently permute the indices of the resulting tensor by $\sigma_i$. Note that crucially, the ``contraction'' is along the last $r_i$ indices of $g^{\mathbf{r}, \boldsymbol{\sigma}}_{i}$. If we do a contraction at $i=i_0$, we can associate $i_0$ with all the indices that got collapsed at $i=i_0$. In other words, when an index gets created, we associate an element of $[t]$ and another one, when it gets annihilated. The graph structure is meant to capture this relation as formalized below. 

\begin{figure}[htb]
\hrule
\vline
\begin{minipage}[t]{0.98\linewidth}
\vspace{10 pt}
\begin{center}
\begin{minipage}[h]{0.95\linewidth}

%
%
%
%
%
%
%

\vspace{3 pt}
\underline{\textsf{Description of graph}}
\begin{enumerate}
\item Initialize bipartite graph with $L = R = \phi$. The vertex set $L$ will be ordered at all points in the algorithm. 

\item Add $q_1$ vertices (in order), each colored as `$1$' to $L$. 

\item For $i = 2$ to $i = t$, 
\begin{enumerate}
\item If $r_{i-1} =0$, then add $q_i$ vertices to $L$ with color `$i$' (in order). Permute the `unmatched' vertices in $L$ with permutation $\sigma_{i-1}$. 
\item If $r_{i-1} = q_i$, then add $q_i$ vertices to $R$ with color `$i$' (in order) and match them to the `last' $q_i$ vertices in $L$. Permute the unmatched vertices in $L$ with permutation $\sigma_{i-1}$. 
\end{enumerate}
\end{enumerate}

\vspace{5 pt}

\end{minipage}
\end{center}

\end{minipage}
\hfill \vline
\hrule
%
\label{fig:DS}
\end{figure}%
%
%
%
%
%
%
%
%
%
~The correspondence between the graph constructed above and the contraction process used to generate $g^{\mathbf{r}, \boldsymbol{\sigma}}_{i+1}$ is obvious. For example, we have the following easy observation. 
\begin{obs}
The dimension of the tensor $g^{\mathbf{r}, \boldsymbol{\sigma}}_t$ is precisely $n^{|U|}$ where $U$ is the set of unmatched vertices in $L$. Consequently, 
$\mathbf{E}[g^{\mathbf{r}, \boldsymbol{\sigma}}_t] \neq 0$ only if $U \neq  \phi$. 
\end{obs} 
Since we are interested in the expectation of $\mathbf{E}[\prod_{i=1}^t p_i]$, thus from now onwards we only focus on $(\mathbf{r}, \boldsymbol{\sigma})$ such that there are no unmatched vertices in the graph generated above. Note that the property of having no unmatched vertices is just a property of $\mathbf{r}$ and not $\boldsymbol{\sigma}$. The next definition forms the crux  of our basis of partitioning the pairs $(\mathbf{r}, \boldsymbol{\sigma})$. 
\begin{definition}
For a pair $(\mathbf{r}, \boldsymbol{\sigma})$ such that the corresponding graph has no unmatched vertices, we say that $(\mathbf{r}, \boldsymbol{\sigma})$ is ``aligned'' if and only each color `$i$', there is a unique color `$j$' such that vertices of color $i$ are only matched to vertices of color `$j$' and vice-versa. Else, it is said to be ``unaligned.''
\end{definition}
We will now prove the following claim. 
\begin{claim}\label{clm:approx-error2}
Let $(\mathbf{r}, \boldsymbol{\sigma})$ be an unaligned pair. Then,  $\Vert g^{\mathbf{r}, \boldsymbol{\sigma}}_t \Vert \le \kappa$.  \end{claim}
\begin{proof}
Note that we are assuming that there are no unmatched vertices in the graph. If the pair is $(\mathbf{r}, \boldsymbol{\sigma})$ is unaligned, then there is a unique smallest integer $1 < N\le t$ such that at the $N$-{th} step of the algorithm to form the graph, the pair becomes unaligned. Namely, $N$ is the smallest integer such that one of the following events happen:

\smallskip
\noindent
\textbf{Event A:} At the beginning of Step (3), when $i=N$, we have $r_{i-1} = q_{i}$ and the last $q_i$ vertices in $L$ are not of the same color. 

\smallskip
\noindent
\textbf{Event B:} At the beginning of Step (3), when $i=N$, we have $r_{i-1} = q_{i}$, the last $q_i$ vertices in $L$ are of the same color (say $j$) but there is also another unmatched vertex in $L$ colored $j$. 

We ask the reader to verify that  $(\mathbf{r}, \boldsymbol{\sigma})$ is unaligned if and only if one of Event A or B happen. We will now prove that in both cases $A$ and $B$, our conclusion holds.

\smallskip
\noindent
\textbf{Proof for Event A: } Let $\mathcal{I}_1 = \{1\} \cup \{N>i > 1: r_{i-1} =0 \}$ and $\mathcal{I}_2 =  \{N>i > 1: r_{i-1}  \neq  0 \}$. By definition of $N$, it is easy to verify that there is a one-one mapping $\nu$ from $\mathcal{I}_2$ to $\mathcal{I}_1$ such that for any $a \in \mathcal{I}_1$, vertices of color $a$ are matched with vertices of color $\nu(a)$ and vice-versa. 
Let us define $\mathcal{I}_3 = \mathcal{I}_1 \setminus \nu(\mathcal{I}_2)$. 
As a consequence, we have $$g^{\mathbf{r}, \boldsymbol{\sigma}}_{N-1} =\prod_{a \in \mathcal{I}_2} \langle h_a , h_{\nu(a)} \rangle \cdot \mu (\otimes_{j \in \mathcal{I}_3} h_j),$$
where $\mu$ is a permutation of the indices of  $\otimes_{j \in \mathcal{I}_3} h_j$. Further, since at $i=N$, the last $q_N$ vertices of $L$ are of at least two different colors, we have that 
the last $q_N$ indices of $\mu (\otimes_{j \in \mathcal{I}_3} h_j)$ belong to more than one $h_j$ (for $j \in \mathcal{I}_3$). Let $A$ be the set of all indices of $\otimes_{j \in \mathcal{I}_3} h_j$ which do not map to the last $q_N$ positions under $\mu$. Consider some assignment $\boldsymbol{\rho}$ for these indices. Then, we have that 
\[
\Vert g^{\mathbf{r}, \boldsymbol{\sigma}}_{N} \Vert_F^2 = \sum_{\boldsymbol{\rho}} \Vert \prod_{a \in \mathcal{I}_2} \langle h_a , h_{\nu(a)} \rangle \cdot \mu (\otimes_{j \in \mathcal{I}_3} h_{j, \boldsymbol{\rho}}) \otimes_{q_N} h_N \Vert_F^2.\]
Here $h_{j, \boldsymbol{\rho}}$ is tensor obtained by restricting the indices in $A$ to $\rho$. Since $\lambda_{\max}(h_N) \le \kappa$, we get
\[
\Vert \mu (\otimes_{j \in \mathcal{I}_3} h_{j, \boldsymbol{\rho}}) \otimes_{q_N} h_N \Vert_F^2 \le \kappa^2 \cdot \Vert \otimes_{j \in \mathcal{I}_3} h_{j, \boldsymbol{\rho}}\Vert_F^2. 
\]
This implies that 
\[
\Vert g^{\mathbf{r}, \boldsymbol{\sigma}}_{N} \Vert_F^2 \le \kappa^2  \sum_{\boldsymbol{\rho}}  \prod_{a \in \mathcal{I}_2} \langle h_a , h_{\nu(a)} \rangle^2 \cdot  \Vert \otimes_{j \in \mathcal{I}_3} h_{j, \boldsymbol{\rho}}\Vert_F^2 \le  \kappa^2 \prod_{a \in \mathcal{I}_2} \langle h_a , h_{\nu(a)} \rangle^2  \le \kappa^2. 
\]
The penultimate inequality uses the fact that $\sum_{\boldsymbol{\rho}}  \Vert \otimes_{j \in \mathcal{I}_3} h_{j, \boldsymbol{\rho}}\Vert_F^2=1$. Using the fact that $\Vert g_j \Vert_F=1$ for all $j>N$ and applying \expref{Fact}{fact:contraction}, we get that $|g^{\mathbf{r}, \boldsymbol{\sigma}}_{t}|\le \kappa$. This finishes the proof.

\smallskip
\noindent
\textbf{Proof for Event B: } Let us define $\mathcal{I}_1$, $\mathcal{I}_2$ and $\mathcal{I}_3$ as in Event A. As before, we have
\[
g^{\mathbf{r}, \boldsymbol{\sigma}}_{N-1} =\prod_{a \in \mathcal{I}_2} \langle h_a , h_{\nu(a)} \rangle \cdot \mu (\otimes_{j \in \mathcal{I}_3} h_j),
\]
where $\mu$ is a permutation of the indices of  $\otimes_{j \in \mathcal{I}_3} h_j$. The crucial difference is that there exists $j_0 \in \mathcal{I}_3$ such that if $\mathcal{B}$ denotes the set of positions where the indices of $h_{j_0}$ get mapped under $\mu$, then this includes the last $q_N$ indices as a proper subset. 
Using $\lambda_{\max}(h_{j_0}) \le \kappa$, we get that 
$\Vert \mu (\otimes_{j \in \mathcal{I}_3} h_{j}) \otimes_{q_N} h_N \Vert_F^2 \le \kappa^2 \cdot \Vert \otimes_{j \in \mathcal{I}_3 \setminus \{j_0\}} h_{j}\Vert_F^2. $
However, $\Vert \otimes_{j \in \mathcal{I}_3 \setminus \{j_0\}} h_{j}\Vert_F^2. =1$, and hence we get $\Vert g^{\mathbf{r}, \boldsymbol{\sigma}}_{N} \Vert_F^2 \le \kappa^2$.
 Using the fact that $\Vert g_j \Vert_F=1$ for all $j>N$ and applying \expref{Fact}{fact:contraction}, we get that $|g^{\mathbf{r}, \boldsymbol{\sigma}}_{t}|\le \kappa$. This finishes the proof.
\end{proof}

\begin{claim}
Let $(\mathbf{r}, \boldsymbol{\sigma})$ be an aligned pair. Then, $g^{\mathbf{r}, \boldsymbol{\sigma}}_{N}$ is just dependent on the covariance matrix $C$ (and the pair $(\mathbf{r}, \boldsymbol{\sigma})$). 
\end{claim}
\begin{proof}
Since $(\mathbf{r}, \boldsymbol{\sigma})$ is an aligned pair, it means that the set $\{1, \ldots, t\}$ can be partitioned into two sets given by $L$ and $R$ and a bijection $\pi : L \rightarrow R$  such that 
$g^{\mathbf{r}, \boldsymbol{\sigma}}_{t} = \prod_{a \in L} \langle h_a, h_{\pi(a)} \rangle$. Note that the mapping $\pi$ is just a function of $(\mathbf{r}, \boldsymbol{\sigma})$ and once $\pi$ is fixed, $g^{\mathbf{r}, \boldsymbol{\sigma}}_{t} $ is just a function of the matrix $C$. 
\end{proof}

\begin{proof}[Proof of \expref{Lemma}{lem:monomial}]
Let $B = \{(\mathbf{r}, \boldsymbol{\sigma}): \mathbf{r} \in \mathcal{S}_z \ \textrm{and} \ (\mathbf{r}, \boldsymbol{\sigma}) \ \textrm{ is aligned}\}$.
Applying (\ref{eq:Ito1}), we get that
\[
    \bigg| \mathbf{E}[\prod_{i=1}^t p_i] - \frac{1}{|\Lambda_{\mathbf{r}}|}\sum_{(\mathbf{r}, \boldsymbol{\sigma}) \in B} \gamma_{\mathbf{r}} I_{m_{t-1}^{\mathbf{r}}}(g^{\mathbf{r}, \boldsymbol{\sigma}}_{t})\bigg| \leq
\sum_{\mathbf{r} \not \in \mathcal{S}_z} \|\gamma_{\mathbf{r}} \cdot
g^{\mathbf{r}}_{t} \| +  \frac{1}{|\Lambda_{\mathbf{r}}|}\sum_{(\mathbf{r},
\boldsymbol{\sigma}) \not \in B :\mathbf{r} \in \mathcal{S}_z}
\|\gamma_{\mathbf{r}} \cdot g^{\mathbf{r}, \boldsymbol{\sigma}}_{t} \|.
\]
Note that $({1}/{|\Lambda_{\mathbf{r}}|}) \sum_{(\mathbf{r}, \boldsymbol{\sigma}) \in B} \gamma_{\mathbf{r}} g^{\mathbf{r}, \boldsymbol{\sigma}}_{t}$ is dependent just on $C$ and $\mathbf{r}$. Further, applying \expref{Claim}{clm:approx-error2}, \expref{Proposition}{prop:approx-error2} and \expref{Claim}{clm:bound-easy1}, we see that the right hand side can be bound by
$2^{t(q_1 + \cdots + q_t +1)} \cdot \kappa$. This finishes the proof. 
\end{proof}


\subsection{Proof of \expref{Theorem}{thm:reduction}}
Let us assume that $f= \mathsf{PTF}(p_1, \ldots, p_k)$ 
where $f$ is the PTF appearing in the statement of \expref{Theorem}{thm:reduction}
After applying \expref{Lemma}{lem:multilinear}, we can pretend that the PTF $f =  \mathsf{PTF}(p_1, \ldots, p_k)$ in \expref{Theorem}{thm:reduction} is multilinear. 
After scaling the polynomials $p_1, \ldots, p_k$,  we can assume that $\mathsf{Var}(p_s)=1$ for all $1 \le s \le k$. With this, we note that there are constants $c_{s,q}$ for $1 \le s \le k$ and $1 \le q \le d$ such that 
\[
p_{s} = p_{s,0} + \sum_{q=1}^d c_{s,q} p_{s,q}, 
\]
where $p_{s,q} \in \mathcal{W}^q$, $\mathsf{Var}(p_{s,q}) =1$ for all $s \in [k]$, $1 \le q \le d$ and $\sum_{q=1}^d c_{s,q}^2=1$. 

Let $\beta: [0, \infty) \to [0, 1)$ be a quickly decreasing function, with
properties that will be specified later. Apply \expref{Theorem}{thm:decomp}
(with $\tau = (\epsilon / (kd))^{3d}$), and let $\tilde p_{s,q}$ (for $0 \le s \le k$ and $1 \le q \le d$) be the resulting polynomials. Defining $\tilde p_s = \tilde p_{s,0} + \sum_{q=1}^d c_{s,q} \tilde p_{s,q}$, \expref{Theorem}{thm:decomp} implies that $\E[p_s - \tilde p_s] = 0$ and
\[
\mathsf{Var}(p_s - \tilde{p}_s) = \sum_{q=1}^d c_{s,q}^2 \mathsf{Var}(p_{s,q} - \tilde{p}_{s,q})  \le \tau. 
\]
Define the PTF $\tilde{f} = \mathsf{PTF}(\tilde{p}_1, \ldots, \tilde{p}_k)$. By
\expref{Theorem}{thm:combine-hyper} and a union bound, 
\begin{align*}
\Vert \mathbf{E}_{x \sim \gamma_n} [\tilde{f}(x) ] - \mathbf{E}_{x \sim \gamma_\ell} [f_{}(x)  ] \Vert&\le\epsilon,  \\
\big| \mathbf{E}_{x \sim \gamma_n} [ \langle \tilde{f}, P_t  \tilde{f} \rangle ] - \mathbf{E}_{x \sim \gamma_{\ell}} [\langle f_{}, P_t f_{}\rangle ] \big| &\le \epsilon, \\ 
\Pr_{x \sim \gamma_{\ell}} [x \in \mathsf{Collision}(\tilde{f})] &\le \Pr_{x \sim \gamma_{n}} [x \in \mathsf{Collision}(f_{})] + \epsilon.
\end{align*}

{}From now on, we will forget about $f$ and work with $\tilde{f}$. Recalling (from \expref{Theorem}{thm:decomp}) that
$$
\tilde{p}_{s,q} = \Outer(\tilde p_{s,q})\left(
\Inner(\tilde p_{s,q})_1(x),\dots, \Inner(\tilde p_{s,q})_{\num({s,q})}(x)\right),
$$
our next task is to ``swap out'' the inner polynomials for a new collection.
Let $\mathcal{P}$ denote the collection of polynomials
$\{\Inner(\tilde p_{s,q})_\ell: 1 \le s \le k, 1 \le q \le d, 1 \le \ell \le \num(s,q)\}$, and
we will introduce another family of polynomials $\mathcal{R} =
\{\Inner(r_{s,q})_\ell: 1 \le s \le k, 1 \le q \le d, 1 \le \ell \le \num(s,q)\}$
that satisfies certain conditions:
\begin{cond}\label{cond:junta} $ $
\begin{itemize}
    \item Each polynomial $\Inner(r_{s,q})_\ell$ has arity $n_0 = O_{s,k,\beta,\epsilon}(1)$.
\item For every $s$, $q$, and $\ell$, $\Inner(r_{s,q})_\ell \in \mathcal{W}^j$ if and only if $\Inner(p_{s,q})_\ell \in \mathcal{W}^j$.
\item $\mathcal{P}$ and $\mathcal{R}$ have the same covariances:
    for every $1 \le s_1, s_2 \le k$, $1 \le q_1, q_2 \le d$, and $1 \le \ell_i \leq \num(s_i, q_i)$ $($for $i \in \{1,2\})$, 
\[
\mathbf{E}[\Inner(\tilde p_{s_1,q_1})_{\ell_1}(x) \cdot \Inner(\tilde p_{s_2,q_2})_{\ell_2}(x) ] = \mathbf{E}[\Inner(r_{s_1,q_1})_{\ell_1}(x) \cdot \Inner(r_{s_2,q_2})_{\ell_2}(x) ].
\]
\item Every polynomial in the family $\mathcal{R}$ is $\beta(\mathrm{Num} + \mathrm{Coeff})$-eigenregular. 
\end{itemize}
\end{cond}

It turns out to be possible to satisfy these conditions:
\begin{lemma}~\label{lem:junta-construct}
There exists a family $\mathcal{R}$ of polynomials on
$n_0= \mathsf{poly}(d, \mathrm{Num}, \beta(\mathrm{Num} + \mathrm{Coeff}))$ variables that
meets the requirements of Condition~\ref{cond:junta}.
\end{lemma}

\begin{proof}
For $1 \le i \le d$, let $\mathcal{P}_i = \{\Inner(p_{s,q})_\ell\}_{\ell =
1,\dots, \num({s,q})} \cap  \mathcal{W}^i$. Let $m_i$ denote the size of
$\mathcal{P}_i$. We will construct $\mathcal{R}$ by constructing the
corresponding set $\mathcal{R}_i$ for each $1 \le i \le d$. Note that
individually constructing $\mathcal{R}_i $ for each $1 \le i \le d$ suffices
for our construction, because if $r_1 \in \mathcal{W}^i$, $r_2 \in
\mathcal{W}^j$ and $i \neq j$ then $\mathbf{E}[r_1(x) \cdot r_2(x)]=0$. 

For convenience of notation, let us enumerate the elements of $\mathcal{P}_i$
as $\{p_1, \ldots, p_{m_i}\}$, and we will construct $\mathcal{R}_i = \{r_1,
\ldots, r_{m_i}\}$. Define $h_j \in \mathcal{H}^{\odot i}$ by $p_j = I_i(h_j)$,
and note that $h_1, \dots, h_{m_i}$ span a space of dimension at most $m_i$.
Setting $\mathcal{H}_1 = \R^{m_i}$, \expref{Proposition}{prop:basis} implies
that $\dim(\mathcal{H}_1^{\odot i}) \ge m_i$, and so there exist
$v_1,\dots, v_{m_i} \in \mathcal{H}_1^{\odot i}$ satisfying $\inr{v_j}{v_{j'}}
= \inr{h_j}{h_{j'}}$ for all $j, j'$.

Let $\delta = \beta(\Num + \Coeff)$, $\kappa = \lceil 1/\delta^2 \rceil$, and $n_0 = m_i \kappa$;
note that $n_0$ is indeed polynomial in $d$, $\Num$, and $\beta(\Num + \Coeff)$.
For $1 \le j \le m_i$, define
$r_j: \mathbb{R}^{n_0} \rightarrow \mathbb{R}$ as follows: divide $x \in
\mathbb{R}^{n_0}$ into $\kappa$ blocks of size $m_i$ each (call the blocks $X_1,
\ldots, X_{\kappa}$) and define
\[
r_j = \frac{1}{\sqrt{\kappa}} \cdot \big( q_{j}(X_1) + \cdots  + q_{j}(X_\kappa)\big). 
\]
Then
\[
\mathbf{E}[r_{j_1} \cdot r_{j_2}] = \sum_{j=1}^\kappa \frac{1}{\kappa} \cdot  \mathbf{E}[q_{j_1}(X_j) \cdot q_{j_2}(X_j)] = \sum_{j=1}^\kappa \frac{1}{\kappa} \cdot \langle v_{j_1}, v_{j_2} \rangle=\sum_{j=1}^\kappa \frac{1}{\kappa} \cdot \langle h_{j_1}, h_{j_2} \rangle = \mathbf{E} [p_{j_1} \cdot p_{j_2}].
\]
The first equality relies on the observation that $q_{j_1}, q_{j_2} \in
\mathcal{W}^{i}$ for $i>1$ and thus $\mathbf{E}[q_{j_1}(X_j) \cdot
q_{j_2}(X_\ell)]= \E[q_{j_1}(X_j)] \E [q_{j_2}(X_\ell)] = 0$ if $j \neq \ell$;
the second and fourth equality use \expref{Proposition}{prop:4}.
Finally, \expref{Fact}{fact:eigen-2} implies that the polynomials $\{r_{j}\}_{1 \le
j \le m_i}$ are $\delta = \beta(\mathrm{Num} + \mathrm{Coeff})$-eigenregular. 
\end{proof}

Having constructed the family $\mathcal{R}$, for $1 \le s \le k$ and $1 \le q \le d$, we define
\[
r_{s,q} =\Outer(\tilde p_{s,q}) \left(\Inner(r_{s,q})_1(x),\dots, \Inner(r_{s,q})_{\num({s,q})}(x)\right).
\]
That is, compared to $\tilde p_{s,q}$, the inner polynomials have changed but the outer polynomial is the same.
For $1 \le s \le k$, we define
\[
    r_s =  \tilde p_{s,0} + \sum_{q=1}^d c_{s,q} r_{s,q},
\]
Finally, we define $\fh = \mathsf{PTF}({r}_1, \ldots, {r}_k)$. Note that $\fh: \mathbb{R}^{n_0} \rightarrow \{0,1\}$.

We will now show that this construction has all the properties stated in \expref{Theorem}{thm:reduction}. 
First of all, note that as long as the function $\beta(\cdot)$ is computable, $n_0$ is a computable function of $d$ and $\epsilon$.
To examine the noise stability of $\fh$, we introduce
a new family of polynomials which will essentially be
``noise-attenuated'' versions of the existing ones. These polynomials will be over
either $2n$ or $2n_0$ variables.
For $1 \le s \le k$, $1 \le q \le d$ and $1 \le \ell \le \num(s,q)$, we define 
\begin{align*}
    \Inner(u_{s,q})_\ell(x,y) &= \Inner(\tilde p_{s,q})_\ell(\eee^{-t}x + \sqrt{1 - \eee^{-2t}}y), \\
    \Inner(v_{s,q})_\ell(x,y) &= \Inner(r_{s,q})_\ell(\eee^{-t}x + \sqrt{1 - \eee^{-2t}}y),
\end{align*}
and
\begin{align*}
    u_{s,q} &= \Outer(\tilde p_{s,q}) \left(\Inner(u_{s,q})_1(x,y),\dots, \Inner(u_{s,q})_{\num({s,q})}(x,y)\right), \\
    v_{s,q} &=\Outer(\tilde p_{s,q}) \left(\Inner(v_{s,q})_1(x,y),\dots, \Inner(v_{s,q})_{\num({s,q})}(x,y)\right),
\end{align*}
and finally
\begin{align*}
    u_s &= u_{s,0} + \sum_{q=1}^d c_{s,q} u_{s,k}, \\
    v_s &= v_{s,0} + \sum_{q=1}^d c_{s,q} v_{s,k},
\end{align*}
where $u_{s,0} = v_{s,0} = p_{s,0}= r_{s,0}$.
To help keep the notation straight, note that every polynomial involving
``$u$'' is on $2n$ variables, while every polynomial involving ``$v$'' is on $2
n_0$ variables. The PTFs corresponding to $u$ and $v$ will be denoted
$f_u = \mathsf{PTF}(u_1,\dots, u_k)$ and $f_v = \mathsf{PTF}(v_1, \dots, v_k)$.

The purpose of introducing these new polynomials is to have the following expressions
for noise stability:
$$
\mathop{\mathbf{E}}_{x \sim \gamma_n} [\langle \tilde{f}(x) , P_{t} \tilde{f}(x) \rangle] = \mathop{\mathbf{E}}_{x \sim \gamma_n, y \sim \gamma_n}[\langle \tilde{f}(x), f_u(x,y) \rangle], \ 
$$ 
$$
\mathop{\mathbf{E}}_{x \sim \gamma_n} [\langle \fh(x) , P_{t} \fh(x)\rangle] = \mathop{\mathbf{E}}_{x \sim \gamma_n, y \sim \gamma_n}[\langle \fh(x), f_v(x,y) \rangle].
$$
The next claim establishes relations between the pairwise correlations of polynomials in the families
$\{\Inner(u_{s,q})_\ell\}$ and $\{\Inner(v_{s,q})_\ell\}$. 

\begin{claim}\label{clm:const-corr}
For all $1 \le s_1, s_2 \le k$, $1 \le q_1, q_2 \le d$ and $1 \le \ell_i \le \mathsf{num}(s_i, q_i)$ (for $i \in \{1,2\}$), we have
\[
\mathbf{E}[\Inner(u_{s_1,q_1})_{\ell_1}(x,y) \cdot \Inner(u_{s_2,q_2})_{\ell_2}(x,y) ] = \mathbf{E}[\Inner(v_{s_1,q_1})_{\ell_1}(x,y) \cdot \Inner(v_{s_2,q_2})_{\ell_2}(x,y) ]
\]
and
\[
\mathbf{E}[\Inner(p_{s_1,q_1})_{\ell_1}(x) \cdot \Inner(u_{s_2,q_2})_{\ell_2}(x,y) ] = \mathbf{E}[\Inner(r_{s_1,q_1})_{\ell_1}(x) \cdot \Inner(v_{s_2,q_2})_{\ell_2}(x,y) ].
\]
\end{claim}
\begin{proof} $ $
Note that $\Inner(u_{s_1,q_1})_{\ell_1}$ and 
$\Inner(u_{s_2,q_2})_{\ell_2}$ are obtained by applying a unitary transformation (on the space of variables) to the polynomials $\Inner(p_{s_1,q_1})_{\ell_1}$ and 
$\Inner(p_{s_2,q_2})_{\ell_2}$. Since the standard Gaussian measure is rotationally invariant,
\[
\mathbf{E}[\Inner(u_{s_1,q_1})_{\ell_1}(x,y) \cdot \Inner(u_{s_2,q_2})_{\ell_2}(x,y) ] = \mathbf{E}[\Inner(\tilde p_{s_1,q_1})_{\ell_1}(x) \cdot \Inner(\tilde p_{s_2,q_2})_{\ell_2}(x) ].
\]
Likewise, 
\[
\mathbf{E}[\Inner(v_{s_1,q_1})_{\ell_1}(x) \cdot \Inner(v_{s_2,q_2})_{\ell_2}(x) ] = \mathbf{E}[\Inner(r_{s_1,q_1})_{\ell_1}(x,y) \cdot \Inner(r_{s_2,q_2})_{\ell_2}(x,y) ].
\]
Since the right hand sides above are equal (by the construction of the family $\{\Inner(r_{s,q})_{\ell}\}$), the first claim follows.

To prove the second claim, note that that there exist $j_1, j_2 \in \mathbb{N}$, such that 
$$
\Inner(p_{s_1,q_1})_{\ell_1}(x), \ \Inner(r_{s_1,q_1})_{\ell_1}(x) \in \mathcal{W}^{j_1} \text{ and } \Inner(u_{s_2,q_2})_{\ell_2}(x,y), \ \Inner(v_{s_2,q_2})_{\ell_2}(x,y) \in \mathcal{W}^{j_2}. 
$$
If $j_1 \ne j_2$, then the second claim holds trivially because both sides are zero;
we will assume, therefore, that $j_1 = j_2 = j$.
Next, we observe that 
\[
\mathbf{E}[\Inner(p_{s_1,q_1})_{\ell_1}(x) \cdot \Inner(u_{s_2,q_2})_{\ell_2}(x,y) ]  = \langle \Inner(p_{s_1,q_1})_{\ell_1}(x),\  P_t \  \Inner(p_{s_2,q_2})_{\ell_2}(x) \rangle,
\]
and recall that $P_t w = \eee^{-jt} w$ for all $w \in \mathcal{W}^j$. Hence,
\begin{align*}
    \mathbf{E}[\Inner(p_{s_1,q_1})_{\ell_1}(x) \cdot \Inner(u_{s_2,q_2})_{\ell_2}(x,y) ] &= \eee^{-jt} \mathbf{E}[\Inner(p_{s_1,q_1})_{\ell_1}(x) \cdot \Inner(p_{s_2,q_2})_{\ell_2}(x) ], \\
\mathbf{E}[\Inner(r_{s_1,q_1})_{\ell_1}(x) \cdot \Inner(v_{s_2,q_2})_{\ell_2}(x,y) ] 
                                                                                         &= \eee^{-jt} \mathbf{E}[\Inner(r_{s_1,q_1})_{\ell_1}(x) \cdot \Inner(r_{s_2,q_2})_{\ell_2}(x) ],
\end{align*}
and as before, the right hand sides are equal.
\end{proof}



We will now list the conditions that we will require on the function $\beta(\cdot)$.
These conditions look somewhat intimidating, but it is easy to see that they
will be satisfied as long as $\beta$ decreases sufficiently quickly.
\begin{cond}\label{cond:enforce-beta}
The function $\beta: [1, \infty) \rightarrow [0,1]$ should satisfy the following conditions: 
\begin{enumerate}
\item For $\xi = {\epsilon}/{(40k^2)}$ and 
for $B^{(1)}, \delta^{(1)}, c_{(1)}$ defined as \[
B^{(1)} = \Omega\bigg( \ln  \frac{\mathsf{Num} \cdot d}{\xi} \bigg)^{d/2} ;  \ \delta^{(1)} = \bigg( \frac{\xi}{d \cdot B^{(1)} \cdot \sqrt{\mathsf{Num}} \cdot \mathsf{Coeff}^{1/d}} \bigg)^d; \   c_{(1)} = \frac{\mathsf{Num}}{\delta^{(1)} \cdot \sqrt{\xi}},
\]
we  require
\begin{align*}
    \beta(\mathsf{Coeff} + \mathsf{Num}) \cdot \mathsf{Coeff}^2 \cdot 2^{\mathsf{Num} \cdot (\mathsf{Num}\cdot 2d+1)} \cdot 2^{O(d^3)} &\leq \frac{1}{4} \text{ and }\\
    2^{O(d \log d)} \cdot \mathsf{Num}^2 \cdot 4c_{(1)}^{2} \cdot \beta(\mathsf{Coeff} + \mathsf{Num}) &\le \xi.
\end{align*}
\item Define  $\xi = {\epsilon}/{(40k^2)}$, $L = 4 \cdot 9^{d+1} \cdot
    (d+1)^2 \cdot \log^{d/2} (k \cdot d/\epsilon) $ and $\vartheta = ({L}/{2}) \cdot
    L^{-2^{2d}}$. Further,  define $B^{(2)}$, $\delta^{(2)}$
    and $c_{(2)}$ as
\[
B^{(2)} = \Omega\bigg( \ln  \frac{\mathsf{2\cdot Num} \cdot 2d}{\xi} \bigg)^{d} ;  \ \delta^{(2)} = \bigg( \frac{\xi \cdot \vartheta^{1/2d}}{2d \cdot B^{(2)} \cdot \sqrt{2 \cdot \mathsf{Num}} \cdot \mathsf{Coeff}^{2/d}} \bigg)^{2d}; \   c_{(2)} = \frac{2 \cdot \mathsf{Num}}{\delta^{(2)} \cdot \sqrt{\xi}}.
\]
We require
\begin{align*}
    \beta(\mathsf{Coeff}^2 + 2 \cdot \mathsf{Num}) \cdot \mathsf{Coeff}^4 \cdot 2^{2\mathsf{Num} \cdot (\mathsf{Num}\cdot 4d+1)} \cdot 2^{O(d^3)} &\leq \frac{\vartheta}{4} \text{ and } \\
    2^{O(d \log d)} \cdot \mathsf{Num}^2 \cdot 4c_{(2)}^{2} \cdot \beta(\mathsf{Coeff}^2 + 2 \cdot \mathsf{Num}) &\le \xi.
 \end{align*}
\end{enumerate}
\end{cond}

Applying the above conditions on $\beta$ will allow us to conclude that our pairs
of polynomials have approximately the same joint distribution of signs. For
the rest of this section, set $\xi = {\epsilon}/{(40k^2)}$.
\begin{corollary}~\label{corr:bounds}
    If $\beta$ satisfies \expref{Condition}{cond:enforce-beta} then
for all $1 \le s, s' \le k$, 
\begin{itemize}
\item[(a)] $\big| \Pr_{x \sim \gamma_{n}}[ \tilde{p}_s \ge 0 ] - \Pr_{x \sim \gamma_{n_0}}[ r_s \ge 0]\big| \le \xi$,
\item[(b)] $\big| \Pr_{x,y \sim \gamma_{n}}[ u_s \ge 0 ] -\Pr_{x,y \sim \gamma_{n_0}}[ v_s \ge 0]\big| \le \xi$,
\item[(c)] $\big| \Pr_{x \sim \gamma_{n}}[ \tilde{p}_s \cdot \tilde{p}_{s'} \ge 0 ] - \Pr_{x \sim \gamma_{n_0}}[ {r}_s \cdot {r}_{s'} \ge 0]\big| \le \xi$,
\item[(d)] $\big| \Pr_{x,y \sim \gamma_{n}}[ {u}_s \cdot {u}_{s'} \ge 0 ] - \Pr_{x,y \sim \gamma_{n_0}}[ {v}_s \cdot {v}_{s'} \ge 0]\big| \le \xi$,
\item[(e)] $\big| \Pr_{x,y \sim \gamma_{n}}[ \tilde{p}_s \cdot {u}_{s'} \ge 0 ] - \Pr_{x,y \sim \gamma_{n_0}}[ {r}_s \cdot {v}_{s'} \ge 0]\big| \le \xi$.
\end{itemize}
\end{corollary}
\begin{proof}
We first note that it suffices to prove (a), (c) and (e). This is because the polynomials $\{{u}_s\}$ (resp.\ $\{{v}_s\}$) are obtained by a unitary transformation of variables 
on the polynomials $\{\tilde{p}_s\}$ (resp.\ $\{ {r}_s\}$). For $1 \le s \le k$, note that there is a $\Psi_s: \mathbb{R}^{\mathsf{Num}} \rightarrow \mathbb{R}$
%
such that
%
\[
\tilde{p}_s = \Psi_s( \{\Inner(\tilde p_{s,q})_1(x),\dots, \Inner(\tilde p_{s,q})_{\num({s,q})}(x)\}_{ 1 \le q \le d}),
\]
\[
{r}_s = \Psi_s( \{\Inner(r_{s,q})_1(x),\dots, \Inner(r_{s,q})_{\num({s,q})}(x)\}_{ 1 \le q \le d}).
\]

\smallskip
\noindent
\textbf{Proof of Item (a):} Recall that $\mathsf{Var}(\tilde{p}_s)=1$. 
Define $B^{(1)}, \delta^{(1)}, c_{(1)}$ as in
\expref{Condition}{cond:enforce-beta}. We now apply \expref{Lemma}{lem:sign-match} with
$\Psi = \Psi_s$, $\eta=1$, $\{A_i\} = \mathcal{P}$, and $\{B_i\} =
\mathcal{R}$. \expref{Condition}{cond:enforce-beta} implies that the $B$ polynomials
(\ie, $\Inner(r_{s,q})_\ell$) satisfy the eigenregularity assumption
in \expref{Lemma}{lem:sign-match}. The conclusion of \expref{Lemma}{lem:sign-match}
proves (a).

\smallskip
\noindent
\textbf{Proof of Item (c):} For $1 \le s, s' \le k$, note that there is a function $\Psi_{s,s'}: \mathbb{R}^{2 \cdot \mathsf{Num}} \rightarrow \mathbb{R}$ such that 
\begin{align*}
    \tilde{p}_s \cdot \tilde{p}_{s'} &= \Psi_{s,s'}( \{\Inner(\tilde p_{s,q})_\ell(x)\}_{q,\ell}, \{\Inner(\tilde p_{s',q})(x)\}_{q,\ell}), \\
    r_s \cdot r_{s'} &= \Psi_{s,s'}( \{\Inner(r_{s,q})_\ell(x)\}_{q,\ell}, \{\Inner(r_{s',q})(x)\}_{q,\ell}).
\end{align*}
Note that $\Psi_{s,s'}$ is a degree-$2d$ polynomial with
$\mathsf{Coeff}_{\Psi_{s,s'}} \le \mathsf{Coeff}^2$ and
$\mathsf{Num}_{\Psi_{s,s'}} \le 2 \mathsf{Num}$. In order to apply \expref{Lemma}{lem:sign-match}, we will need a lower bound on the variance of $\tilde p \cdot \tilde p_s$:
\[
\mathsf{Var}(\tilde{p}_s \cdot \tilde{p}_{s'}) \ge \frac{L}{2} \cdot L^{-2^{2d}},
\]
where $L = 4 \cdot 9^{d+1} \cdot (d+1)^2 \cdot \log^{d/2} (k \cdot d/\epsilon)$.
(This lower bound will be proven as \expref{Lemma}{lemma:variance-product-lower}.) Note
that the right hand side above is exactly what we called $\vartheta$ in
\expref{Condition}{cond:enforce-beta}.
Now we will apply \expref{Lemma}{lem:sign-match} with $\eta = \vartheta$:
\expref{Claim}{clm:const-corr}
implies that the second moments match, while the second part of
\expref{Condition}{cond:enforce-beta} ensures that the polynomials
$\Inner(r_{t,q})_\ell$ are sufficiently eigenregular.
The conclusion of \expref{Lemma}{lem:sign-match} proves item (c).

\smallskip
\noindent
\textbf{Proof of Item (e):} The proof of item (e) is the same as the proof of item (c) with $u_{s'}$ taking over the role of $\tilde{p}_{s'}$. We omit the details.
\end{proof}

\begin{lemma}~\label{lemma:variance-product-lower}
Let $p, q: \mathbb{R}^n \rightarrow \mathbb{R}$ be degree-$d$ polynomials such that $\mathsf{Var}(p) = \mathsf{Var}(q)=1$. Let $T = \max \{\mathbf{E}[p^2], \mathbf{E}[q^2]\}$ and 
$L=4 \cdot T \cdot 9^{d+1} \cdot (d+1)^2$. 
Then
\[
\mathsf{Var}(p \cdot q) \ge \frac{L}{2} \cdot (L)^{2^{-2d}}. 
\] 
\end{lemma}

We will begin with an easier lower bound, in terms of the ``highest-level weights:''
\begin{lemma}\label{lem:variance-bound}
    If $p = \sum_{i=0}^{d_1} I_i(f_i)$ and $q = \sum_{i=0}^{d_2} I_i(g_i)$
    for $f_i \in \mathcal{H}^{\odot i}$ and $g_i \in \mathcal{H}^{\odot i}$ then
\[
\mathsf{Var}(p \cdot q) \ge {\Vert f_{d_1} \Vert_F^2 \cdot \Vert g_{d_2} \Vert_F^2}.
\] 
\end{lemma}
\begin{proof} 
\begin{eqnarray*}
p \cdot q &=& \sum_{r_1=0}^{d_1} \sum_{r_2=0}^{d_2} I_{r_1}(f_{r_1}) \cdot I_{r_2}(g_{r_2})  \\
&=& \sum_{r_1=0}^{d_1} \sum_{r_2=0}^{d_2} \sum_{r=0}^{r_1 \wedge r_2} r! \cdot \binom{r_1}{r} \cdot \binom{r_2}{r} \cdot \frac{\sqrt{(r_1+r_2 - 2r)!}}{\sqrt{r_1!}\sqrt{r_2!}} \cdot I_{r_1+r_2 - 2r} (f_{r_1} \widetilde{\otimes_r} g_{r_2}). 
\end{eqnarray*}
The second equality uses It\^{o}'s multiplication formula
(\expref{Proposition}{prop:Ito}). The term with $r=0$, $r_1= d_1$ and
$r_2=d_2$ is orthogonal to all the other terms because it is the only
one belonging to the $(r_1 + r_2)$th Wiener chaos; hence,
\[
\mathsf{Var}(p \cdot q) \ge \mathsf{Var} \bigg( \frac{\sqrt{(d_1+d_2 )!}}{\sqrt{d_1!}\sqrt{d_2!}}  \cdot I_{d_1+d_2} (f_{d_1} \widetilde{\otimes_0} g_{d_2}) \bigg)
= \frac{(d_1 + d_2)!}{d_1! \cdot d_2!} \|f_{d_1} \widetilde{\otimes_r} g_{d_2}\|_F^2.
\]
Finally, Neuberger~\cite{Neuberger:74} showed that
\[
  \|f_{d_1} \widetilde{\otimes_r} g_{d_2}\|_F^2 \ge \frac{1}{\binom{d_1 + d_2}{d_1}} \|f_{d_1}\|_F^2 \|g_{d_2}\|_F^2,
\]
which completes the proof.
\end{proof}


\begin{proof}[Proof of \expref{Lemma}{lemma:variance-product-lower}]
    Let us express $p$ and $q$ in terms of iterated It\^{o} integrals. 
\[
p = \sum_{r=0}^{d} I_r(f_r) \ \textrm{ and } \ q = \sum_{r=0}^{d} I_r(g_r).
\]
Let $\alpha = \mathbf{E}[p^2]$ and $\beta= \mathbf{E}[q^2]$. 
Let $\Gamma : \mathbb{N} \rightarrow [0,1]$ be a function (which we shall fix later). Consider the following iterative process:
\begin{itemize}
\item Start with $i=d$ and $j=d$. If $\Vert f_i \Vert_F \cdot \Vert g_j \Vert_F \ge \Gamma(i+j)$, then stop the process. 
\item If $\Vert f_i \Vert_F \le \sqrt{\Gamma(i+j)}$, then $i \leftarrow i-1$, else $j \leftarrow  j-1$. 
\item If any of $i$ or $j$ reaches zero, terminate the process. 
\end{itemize}
Recall that  $T = \max \{\alpha, \beta\} $ and $L = 4 \cdot T \cdot 9^{d+1} \cdot (d+1)^2$. We let $\Gamma (x) = L \cdot L^{-2^{x}}$. First, we claim that for $\Gamma(\cdot)$ as chosen here, the iterative process above terminates with $(i,j)$ such that neither of them is zero. Towards a contradiction, assume that the above process terminates with $i=0$ and $j>0$. Further, for $d \le \ell \le 1$, when the value of $i$ drops from $\ell$ to $\ell-1$, let the value of $i+j$ be $\kappa_\ell$. Note that $\kappa_1 < \cdots < \kappa_d$ and $\kappa_1 \ge 2$.  Thus, if the process terminates with $i=0$, 
then 
\[
\mathsf{Var}(p)= \sum_{\ell=1}^d \Vert f_\ell \Vert_F^2  \le \sum_{\ell=1}^d  \Gamma(\kappa_\ell) \le 2 \cdot \Gamma(\kappa_1)  <1.
\]
This results in a contradiction which means that both $i$ and $j$ must be non-zero at the end of the process.  Here the penultimate inequality uses that for our choice of $\Gamma(\cdot)$, $\Gamma(x+1) \le \Gamma(x)/2$ and the last inequality uses that $\Gamma(\kappa_1) <1/2$. 

Let us assume that $(i_0, j_0)$ is the pair returned by the above iterative process and define $\tilde{p} = \sum_{r=0}^{i_0} I_r(f_r)$ and $\tilde{q} = \sum_{r=0}^{j_0} I_r(g_r)$. Applying \expref{Lemma}{lem:variance-bound}, we have
$
\mathsf{Var}(\tilde{p} \cdot \tilde{q}) \ge {\Gamma^2(i_0+j_0)}. 
$
Observe that $p \cdot q = \tilde{p} \cdot \tilde{q} + (p - \tilde{p}) \cdot q + \tilde{p} \cdot (q-\tilde{q})$. Thus, it follows that $\mathsf{Var}(p \cdot q - \tilde{p} \cdot \tilde{q}) = \mathsf{Var}((p - \tilde{p}) \cdot q + \tilde{p} \cdot (q-\tilde{q}))$. Applying Jensen's inequality, we get
\[
\sqrt{\mathsf{Var}(p \cdot q)} \ge \sqrt{\mathsf{Var}(\tilde{p} \cdot \tilde{q})} - \sqrt{\mathsf{Var}(\tilde{p} \cdot {(q-\tilde{q})})} - \sqrt{\mathsf{Var}(q \cdot {(p-\tilde{p})})}.\]
Hence, it suffices to bound $\mathsf{Var}(\tilde{p} \cdot {(q-\tilde{q})})$ and $\mathsf{Var}(q \cdot {(p-\tilde{p})})$. To do this, notice that by definition of the process, 
\[
\mathsf{Var}(q-\tilde{q}), \ \mathsf{Var}(p-\tilde{p}) < \sum_{\ell = i_0 +j_0+1}^{d} \Gamma(\ell) <  2 \cdot \Gamma(i_0+j_0+1). 
\]
Now, Cauchy-Schwarz and the Gaussian hypercontractive inequality imply that
$\E[(r\cdot s)^2] \le 9^d \E[r^2] \E [s^2]$ for any degree-$d$ polynomials $r$ and $s$.
Hence,
\[
\mathsf{Var}(\tilde p \cdot {(q-\tilde{q})}), \  \mathsf{Var}(q \cdot {(p-\tilde{p})})\le 2 \cdot \Gamma(i_0+j_0+1) \cdot d \cdot (d+1) \cdot 8^d \cdot T. 
\]
Thus, 
\[
\sqrt{\mathsf{Var}(p \cdot q)} \ge {\Gamma(i_0+j_0)} - 2 \sqrt{2} \cdot (d+1) \cdot (2\sqrt{2})^d \cdot \sqrt{T} \cdot \sqrt{\Gamma(i_0+j_0+1)}.
\]
However, note that for any $x \ge 0$, $\sqrt{\Gamma(x+1) } = \Gamma(x) / \sqrt{L}$. Plugging the value of $L$, we obtain that 
\[
\sqrt{\mathsf{Var}(p \cdot q)}  \ge {\Gamma(i_0+j_0)}/2 \ge \frac{L}{2} \cdot L^{2^{-2d}}. \qedhere
\]
%
%
%
%
%
%
%
%
%
%
%
%
%
%
%
%
%
%
%
%
%
%
\end{proof}


Before continuing with the proof, let us summarize what we have shown so far:
\begin{theorem}\label{thm:summary}
Fix $\epsilon > 0$ and $t > 0$, and let $p_1, \ldots, p_k: \mathbb{R}^n \rightarrow \mathbb{R}$ be degree-$d$ polynomials
satisfying $\Var(p_s) = 1$ and $|\E[p_s]| \le \log^{d/2} (k \cdot d/\epsilon)$ for 
all $1 \le s \le k$.
There exists a computable $n_0 = n_0(k, d, t, \epsilon)$
and polynomials $r_1, \dots, r_k: \R^{n_0} \to \R$ such that
for every $1 \le s, s' \le k$:
\begin{enumerate}
\item[(a)] $\big|\Pr_{x \sim \gamma_n} [{p}_s \ge 0] - \Pr_{x \sim \gamma_n} [{r}_s \ge 0]\big| \le \epsilon$,
\item[(b)] $\big|\Pr_{x,y \sim \gamma_{n}} [{u}_s \ge 0] - \Pr_{x,y \sim \gamma_{n_0}} [{v}_s \ge 0]\big| \le \epsilon$,
\item[(c)] $\big|\Pr_{x \sim \gamma_n} [{p}_s \cdot {p}_{s'} \ge 0] - 
\Pr_{x \sim \gamma_{n_0}} [{r}_s \cdot {r}_{s'} \ge 0]\big| \le \epsilon$,
\item[(d)] $\big|\Pr_{x,y \sim \gamma_n} [{u}_s\cdot {u}_{s'} \ge 0] - 
\Pr_{x,y \sim \gamma_{n_0}} [{v}_s \cdot {v}_{s'} \ge 0]\big| \le \epsilon$,
\item[(e)] $\big|\Pr_{x,y \sim \gamma_n} [{p}_s\cdot {u}_{s'} \ge 0] - 
\Pr_{x,y \sim \gamma_{n_0}} [{r}_s \cdot {v}_{s'} \ge 0]\big| \le \epsilon$,
\end{enumerate}
where $u_s(x, y) = p_s(\eee^{-t} x + \sqrt{1-\eee^{-2t}} y)$ and
$v_s(x, y) = r_s(\eee^{-t} x + \sqrt{1-\eee^{-2t}} y)$.
\end{theorem}

Indeed, \expref{Theorem}{thm:summary} follows directly from \expref{Corollary}{corr:bounds}
and the fact that $\mathsf{sign}(p_s)$ and $\mathsf{sign}(\tilde p_s)$ disagree with probability at most $O(\epsilon)$.

Returning to the notation of \expref{Corollary}{corr:bounds}, our next goal is to
prove
that $\Pr_{x \sim \gamma_{n_0}} [x \in \mathsf{Collision}(\fh)]$ is small.

\begin{claim}~\label{clm:collision-upper}
$\displaystyle \Pr_{x \sim \gamma_{n_0}} [x \in \mathsf{Collision}(\fh)]\le  (k^2+1) \cdot \Pr[x \in \mathsf{Collision}(f)] + k \cdot \xi + \frac{3k^2 \xi}{2}. $ \end{claim}

\begin{proof}
For
a multivariate PTF $f = \mathsf{PTF}(p_1, \ldots, p_k)$, define
$\mathsf{Unique}(f,i) = \overline{\mathsf{Collision}(f)} \cap \{x : p_i(x)  \ge
0\}$. In other words, it is the set of all points such that $i$ is the unique
index such that $p_i(x) \ge 0$.
Since we haven't mentioned $f$ and $\fh$ for a while, recall that $f =
\mathsf{PTF}(p_1, \dots, p_k)$ and $\fh = \mathsf{PTF}(r_1, \dots, r_k)$.

Next, observe that for real numbers $a, b \ne 0$,
\begin{equation}
  \label{eq:identity}
  \mathbf{I} \ [a \ge 0 \wedge b \ge 0] = \frac12 \big( \mathbf{I} \ [ a \cdot b \ge 0] + \mathbf{I} \ [a \ge 0] + \mathbf{I} \ [b \ge 0] - 1\big)
\end{equation}
(where $\mathbf{I}(P)$ is 1 if $P$ is true and $0$ otherwise).
Since $p_s(x) \ne 0$ almost surely for every $s$, the identity above implies that
\begin{align}
&| \Pr [p_{s}(x) \ge 0 \wedge p_{s'}(x) \ge 0]  - \Pr [r_{s}(x) \ge 0 \wedge r_{s'}(x) \ge 0] | \notag \\
& \le \frac{1}{2} \big( \big| \Pr[p_{s}(x) \ge 0 ]  - \Pr [r_{s}(x) \ge 0 ] \big| +  \big| \Pr_{x \sim \gamma_n} [p_{s'}(x) \ge 0 ]  - \Pr_{x \sim \gamma_n} [r_{s'}(x) \ge 0 ] \big|\big) \notag \\
& \qquad +\frac{1}{2}  \big( \big| \Pr_{x \sim \gamma_n} [p_{s'}(x) \cdot p_s(x) \ge 0 ]  - \Pr_{x \sim \gamma_n} [r_{s'}(x)  \cdot r_s(x)\ge 0 ] \big|\big) \notag \\
&\le \frac{3\xi}{2},
\label{eq:gamma1}
\end{align}
where the last inequality follows from the first and third items of \expref{Corollary}{corr:bounds}. Now observe that 
\begin{align*}
\Pr_{x \sim \gamma_n} [p_s(x) \ge 0] - \sum_{s' \neq s} \Pr_{x \sim \gamma_n} [p_s(x) \ge 0 \wedge p_{s'}(x) \ge 0]  &\leq \Pr_{x \sim \gamma_n} [x \in \mathsf{Unique}(f,s)] \le \Pr_{x \sim \gamma_n} [p_s(x) \ge 0],\\
\Pr_{x \sim \gamma_n} [r_s(x) \ge 0] - \sum_{s' \neq s} \Pr_{x \sim \gamma_n} [r_s(x) \ge 0 \wedge r_{s'}(x) \ge 0]  &\leq \Pr_{x \sim \gamma_n} [x \in \mathsf{Unique}(\fh,s)] \le \Pr_{x \sim \gamma_n} [r_s(x) \ge 0].
\end{align*}
Subtracting these two and applying (\ref{eq:gamma1}), we obtain
\begin{eqnarray*}
\big|\Pr_{x \sim \gamma_n} [x \in \mathsf{Unique}(f,s)]- \Pr_{x \sim \gamma_n} [x \in \mathsf{Unique}(\fh,s)] \big| &\leq& \big| \Pr_{x \sim \gamma_n} [p_s(x) \ge 0] -\Pr_{x \sim \gamma_n} [r_s(x) \ge 0]\big| \\ &+& \big|\sum_{s' \neq s} \Pr_{x \sim \gamma_n} [p_s(x) \ge 0 \wedge p_{s'}(x) \ge 0]\big| + \frac{3k \xi}{2}.
\end{eqnarray*}
Adding over all $s$ and applying the first item of \expref{Corollary}{corr:bounds},  we obtain 
\[
\sum_{s=1}^k \big|\Pr_{x \sim \gamma_n} [x \in \mathsf{Unique}(f,s)]- \Pr_{x \sim \gamma_n} [x \in \mathsf{Unique}(\fh,s)] \big| \le k\cdot \xi + \frac{3k^2 \xi}{2} +\sum_{s' \neq s} \Pr_{x \sim \gamma_n} [p_s(x) \ge 0 \wedge p_{s'}(x) \ge 0]
\]

Noting that for all $(s,s')$, $\Pr [p_s(x) \ge 0 \wedge p_{s'}(x) \ge 0] \le \Pr [x \in \mathsf{Collision}(f)]$, we obtain
\[
\sum_{s=1}^k \big|\Pr_{x \sim \gamma_n} [x \in \mathsf{Unique}(f,s)]- \Pr_{x \sim \gamma_n} [x \in \mathsf{Unique}(\fh,s)] \big| \le k^2 \cdot \Pr[x \in \mathsf{Collision}(f)] + k \cdot \xi + \frac{3k^2 \xi}{2}. 
\]
This however immediately implies that 
\[
\Pr_x [x \in \mathsf{Collision}(\fh)] \le  (k^2+1) \cdot \Pr[x \in \mathsf{Collision}(f)] + k \cdot \xi + \frac{3k^2 \xi}{2}. 
\qedhere
\]
\end{proof}

\begin{lemma}~\label{lem:noise-stab}
\[
\big| \mathbf{E}[\langle f, P_t f \rangle]  - \mathbf{E}[\langle \fh, P_t \fh \rangle]\big| \le 2\Pr[x \in \mathsf{Collision}(f)] +2 \Pr[x \in \mathsf{Collision}(\fh)] + \frac{3k\xi}{2}. 
\]
\end{lemma}
\begin{proof}
We begin by noting the following two inequalities :
\[
\big| \mathbf{E}[\langle f, P_t f \rangle]  -\sum_{s=1}^k \Pr[p_s(x) \ge 0 \wedge u_s(x,y) \ge 0] \big| \le 2 \Pr[x \in \mathsf{Collision}(f)],
\]
\[
\big| \mathbf{E}[\langle \fh, P_t \fh \rangle]  -\sum_{s=1}^k \Pr[r_s(x) \ge 0 \wedge v_s(x,y) \ge 0] \big| \le 2 \Pr[x \in \mathsf{Collision}(\fh)].
\]
Thus, 
\begin{eqnarray}
\big| \mathbf{E}[\langle f, P_t f \rangle] - \mathbf{E}[\langle \fh, P_t \fh \rangle] \big| &\le& 2\Pr[x \in \mathsf{Collision}(f)] +2 \Pr[x \in \mathsf{Collision}(\fh)] \nonumber \\ 
&+& \sum_{s=1}^k \big|\Pr[p_s(x) \ge 0 \wedge u_s(x,y) \ge 0]- \Pr[r_s(x) \ge 0 \wedge v_s(x,y) \ge 0] \big|. \nonumber \\ \label{eq:noise-junta-1}
\end{eqnarray}
We now seek to bound $\Pr[p_s(x) \ge 0 \wedge u_s(x) \ge 0]- \Pr[r_s(x) \ge 0 \wedge v_s(x) \ge 0]$. As before, using that the polynomials $p_s, u_s, r_s$ and $v_s$ are zero only
on a measure zero set and using the identity from \eqref{eq:identity}, we obtain that 
\begin{eqnarray*}
&& \big|\Pr[p_s(x) \ge 0 \wedge u_s(x,y) \ge 0]- \Pr[r_s(x) \ge 0 \wedge v_s(x,y) \ge 0] \big| \\ &\le& \frac{1}{2} \cdot  \big|\Pr[p_s(x) \ge 0]- \Pr[r_s(x) \ge 0] \big| +  \frac{1}{2} \cdot  \big|\Pr[u_s(x,y) \ge 0]- \Pr[v_s(x,y) \ge 0] \big|  \\ &+&  \frac{1}{2} \cdot  \big|\Pr[p_s(x) \cdot u_s(x,y) \ge 0]- \Pr[r_s(x) \cdot v_s(x,y) \ge 0] \big|.
\end{eqnarray*} 
Applying \expref{Corollary}{corr:bounds} to bound all the terms, we obtain that 
$$
 \big|\Pr[p_s(x) \ge 0 \wedge u_s(x,y) \ge 0]- \Pr[r_s(x) \ge 0 \wedge v_s(x,y) \ge 0] \big| \le \frac{3 \xi}{2}. 
$$
Plugging this bound back to (\ref{eq:noise-junta-1}), we obtain that 
\[
\big| \mathbf{E}[\langle f, P_t f \rangle] - \mathbf{E}[\langle \fh, P_t \fh \rangle] \big| \le   2\Pr[x \in \mathsf{Collision}(f)] +2 \Pr[x \in \mathsf{Collision}(\fh)] + \frac{3k\xi}{2}. 
\qedhere
\]
\end{proof}

We are now in a position to finish the proof of \expref{Theorem}{thm:reduction} and
show that the construction $\fh$ indeed satisfies the required properties. Note
that by construction, $\fh$ is a degree-$d$ PTF over $n_0$ variables where as
we have said before, once $\beta (\cdot)$ is a computable function, $n_0 = n_0(k,
\epsilon, d, t)$ is a computable function. Further, it is easy to see that
$\beta(\cdot)$ can be chosen to be computable; we just need to decrease sufficiently quickly
to satisfy
\expref{Condition}{cond:enforce-beta}. 

We first bound $\Vert \mathbf{E}[\fh] - \boldsymbol{\mu}\Vert_1 = \Vert \mathbf{E}[\fh] - \mathbf{E}[f] \Vert_1$. To do this, note that 
\begin{align*}
  \Vert \mathbf{E}[\fh] - \mathbf{E}[f] \Vert_1 &\le \sum_{s=1}^k |\Pr[\tilde{p}_s(x) \ge 0] - \Pr[{r}_s(x) \ge 0] | + \Pr[x \in \mathsf{Collision}(f)]\\
  &\phantom{\le \sum_{s=1}^k |\Pr[\tilde{p}_s(x) \ge 0] - \Pr[{r}_s(x) \ge 0] }\:\:\:\: + \Pr[x \in \mathsf{Collision}(\fh)] \\ &\le k \cdot \xi + \Pr[ x \in \mathsf{Collision}(f)] + (k^2 +1) \cdot \Pr[ x \in \mathsf{Collision}(f)] + k \cdot \xi + \frac{3 k^2 \cdot \xi}{2}. 
\end{align*}
Here the last inequality follows by applying the first item of \expref{Corollary}{corr:bounds} and \expref{Claim}{clm:collision-upper}. Plugging in the value of $\xi = {\epsilon}/{(20k^2)}$ and the upper bound on $\Pr[x \in \mathsf{Collision}(f)]$, we obtain that $\Vert \mathbf{E}[\fh] - \mathbf{E}[f] \Vert_1 \le \epsilon$. 

For the second part, note that applying \expref{Lemma}{lem:noise-stab}, we obtain
\begin{eqnarray*}
\langle \fh, P_t \fh \rangle \ge \langle f, P_t f \rangle - 2\Pr[x \in \mathsf{Collision}(f)] -2 \Pr[x \in \mathsf{Collision}(\fh)] - \frac{3k\xi}{2}. 
\end{eqnarray*}
Again plugging in the value $\xi = {\epsilon}/{(20k^2)}$ and the upper bound on $\Pr[x \in \mathsf{Collision}(f)]$, we obtain that 
$$
\langle \fh, P_t \fh \rangle \ge \langle f, P_t f \rangle - \epsilon.
$$
This completes the proof of \expref{Theorem}{thm:reduction}.

\nocite{DMN17-conf}   %
  %
\nocite{Kalai:02}     %

\bibliographystyle{tocplain}   %

%
%
%
%
\bibliography{v015a006}

%
%
%
%
%
\begin{tocauthors}
\begin{tocinfo}[de]
 Anindya De\\
 Assistant professor\\
 University of Pennsylvania\\
 anindyad\tocat{}seas\tocdot{}upenn\tocdot{}edu \\   %
 \url{http://www.seas.upenn.edu/~anindya}      %
\end{tocinfo}
\begin{tocinfo}[mossel]
  Elchanan Mossel\\
  Professor\\
  Massachusetts Institute of Technology\\
  elmos\tocat{}mit\tocdot{}edu\\
  \url{http://www.math.mit.edu/~elmos/}
\end{tocinfo}
\begin{tocinfo}[neeman]
  Joe Neeman\\
  Assistant professor\\
  University of Texas, Austin\\
  jneeman\tocat{}math\tocdot{}utexas\tocdot{}edu\\
  \url{https://web.ma.utexas.edu/users/jneeman/}
\end{tocinfo}
\end{tocauthors}
%

%

\begin{tocaboutauthors}
\begin{tocabout}[de]
\textsc{Anindya De} is an assistant professor of computer science at
the  \href{www.upenn.edu}{University of Pennsylvania}. He received his
\phd\ from \href{http://www.berkeley.edu/index.html}{U.\,C. Berkeley}
in 2013, advised by Luca Trevisan.
He was a postdoc at the \href{https://www.ias.edu/}{Institute for
Advanced Study} and \href{dimacs.rutgers.edu/}{DIMACS} from
2013 to 2015 and an assistant professor at
\href{www.northwestern.edu}{Northwestern University} from
2015 to 2018.  His research interests are in complexity theory,
learning theory, discrete analysis and applied probability.
\end{tocabout}
\begin{tocabout}[mossel]
\textsc{Elchanan Mossel} is a professor of mathematics at MIT\@.
He received his \phd\ in mathematics from Hebrew University in 2000,
advised by Yuval Peres.  After his \phd, he was a postdoc in
the Theory Group of Microsoft Research, Redmond,
and a Miller fellow in Statistics and Computer Science at U.\,C. Berkeley.
His research interests include combinatorial statistics,
discrete Fourier analysis and influences, randomized algorithms,
computational complexity, MCMC, Markov random fields,
social choice, game theory and evolution.
\end{tocabout}
\begin{tocabout}[neeman]
\textsc{Joe Neeman} is a professor of mathematics at the
\href{www.math.uni-bonn.de}{University of Bonn} and
an assistant professor of mathematics at the
\href{www.math.utexas.edu}{University of Texas, Austin}.
He received his \phd\ in 2013 from
\href{http://www.stat.berkeley.edu}{U.\,C. Berkeley},
advised by Elchanan Mossel, after which he was a postdoc at U.\,T. Austin.
He works in various areas of probability and on applications to computer
science.
\end{tocabout}

\end{tocaboutauthors}


\end{document}

%
%
%
%
